% Lecture notes on the classification of homogeneous plane continua,
% following the formal (Lean 4) proof in this repository.
% Volume 6: The indecomposable case.
%
% Compile with:  pdflatex main && pdflatex main   (two runs for the contents and references)
\documentclass[11pt,oneside]{book}

\usepackage[T1]{fontenc}
\usepackage{lmodern}
\usepackage[margin=1.15in]{geometry}
\usepackage{amsmath,amssymb,amsthm}
\usepackage{enumitem}
\usepackage{booktabs}
\usepackage{longtable}
\usepackage{tikz}
\usepackage[hidelinks]{hyperref}

% ---------------------------------------------------------------- theorem styles
\theoremstyle{plain}
\newtheorem{theorem}{Theorem}[chapter]
\newtheorem{lemma}[theorem]{Lemma}
\newtheorem{proposition}[theorem]{Proposition}
\newtheorem{corollary}[theorem]{Corollary}
\newtheorem{claim}[theorem]{Claim}
\theoremstyle{definition}
\newtheorem{definition}[theorem]{Definition}
\newtheorem{example}[theorem]{Example}
\newtheorem{exercise}[theorem]{Exercise}
\theoremstyle{remark}
\newtheorem{remark}[theorem]{Remark}

% ---------------------------------------------------------------- notation
\newcommand{\N}{\mathbb{N}}
\newcommand{\Z}{\mathbb{Z}}
\newcommand{\Q}{\mathbb{Q}}
\newcommand{\R}{\mathbb{R}}
\newcommand{\C}{\mathbb{C}}
\newcommand{\diam}{\operatorname{diam}}
\newcommand{\dist}{\operatorname{dist}}
\newcommand{\interior}{\operatorname{int}}
\newcommand{\cl}[1]{\overline{#1}}
\newcommand{\abs}[1]{\left\lvert #1\right\rvert}
\newcommand{\set}[1]{\left\{ #1 \right\}}
\newcommand{\nrm}[1]{\lVert #1\rVert}
% the composant of p in M
\newcommand{\cmp}[2]{\operatorname{cmp}_{#1}(#2)}
% the frontier of a set
\newcommand{\fr}{\operatorname{fr}}
% the interior of a compact set of the plane (union of the bounded complementary components)
\newcommand{\inn}{\operatorname{ins}}
% the layer k^{-1}(s) of an admissible map
\newcommand{\lay}[2]{#1^{-1}(#2)}
% the Effros neighbourhood
\newcommand{\Eff}{\mathcal{W}}
% coordinate projections
\newcommand{\pr}[1]{\pi_{#1}}
% the diagonal
\newcommand{\Diag}{\Delta}
% the parity of a closed polygon at a point
\newcommand{\parity}{\operatorname{par}}
% the relative frontier of a subgraph
\newcommand{\rfr}{\partial}
% a side of a set over a graph
\newcommand{\side}{\operatorname{side}}
% the unit interval
\newcommand{\I}{[0,1]}

% Names of declarations in the Lean development (typewriter, underscores allowed).
\newcommand{\lean}[1]{\texttt{\detokenize{#1}}}
% A Lean name ending in a subscript digit.
\newcommand{\leansub}[2]{\texttt{\detokenize{#1}}\ensuremath{_{#2}}}
% A paragraph pointing to the formal counterpart of a result.
\newcommand{\formal}[1]{\par\smallskip\noindent{\small\textit{Formal counterpart:} #1}\par\smallskip}
% References to numbered results of the earlier volumes.
\newcommand{\vone}[1]{Volume~1, #1}
\newcommand{\vtwo}[1]{Volume~2, #1}
\newcommand{\vthree}[1]{Volume~3, #1}
\newcommand{\vfour}[1]{Volume~4, #1}
\newcommand{\vfive}[1]{Volume~5, #1}

\setlist{itemsep=2pt,topsep=4pt}

\title{\textbf{Homogeneous Plane Continua}\\[6pt]
\Large Lecture notes following the formal proof\\[18pt]
\large Volume 6\\[4pt]
\Large The indecomposable case:\\ hereditary indecomposability, span zero\\ and snake-likeness}
\author{}
\date{}

\begin{document}

\frontmatter
\maketitle
\tableofcontents

\mainmatter
% The introduction is Chapter 0, as in Volumes 1--5.
\setcounter{chapter}{-1}
%======================================================================
\chapter{Introduction}\label{ch:intro}
%======================================================================

\section{What this volume is about}

The classification of homogeneous plane continua (\vone{Theorem~0.1}) splits into an indecomposable and a decomposable case. This
volume treats the indecomposable case:

\begin{theorem}[Theorem~\ref{thm:indec-pseudoarc}]
A nondegenerate plane continuum is homogeneous and indecomposable if and only if it is a pseudo-arc.
\end{theorem}

The ``if'' direction is Bing's (Volume~2). The ``only if'' direction has three steps, each a theorem from the literature:
\begin{enumerate}
\item \textbf{Hereditary indecomposability} (Hagopian 1976; Theorem~\ref{thm:HI}): every homogeneous indecomposable plane continuum is hereditarily
indecomposable. The proof used here is Hagopian's of 1984 (following a plan of Jones): if not, a monotone decomposition produces a homogeneous plane
continuum that is indecomposable and contains an arc, and no such continuum exists (Bing 1960, Jones 1975; Theorem~\ref{thm:noarc}).
\item \textbf{Span zero} (Theorem~C of Hoehn--Oversteegen, from Oversteegen--Tymchatyn 1982; Theorem~\ref{thm:theoremC}): every homogeneous indecomposable plane
continuum has span zero.
\item \textbf{Snake-likeness} (Hoehn--Oversteegen 2016, Theorem~1; Theorem~\ref{thm:ho1}): a hereditarily indecomposable plane continuum of span zero is snake-like.
\end{enumerate}
A nondegenerate snake-like hereditarily indecomposable continuum is a pseudo-arc by definition (\vone{Definition~2.14}).

The volume has ten chapters.

\begin{itemize}
\item \textbf{Chapter~\ref{ch:homog}: First consequences of homogeneity in the plane.} Composants, Effros neighbourhoods, Moore's triod theorem; such a continuum is
atriodic and hereditarily unicoherent, and no proper subcontinuum separates the plane (Hagopian 1975, 1976).
\item \textbf{Chapter~\ref{ch:typelambda}: Continua of type $\lambda$.} Irreducible continua with a monotone map onto $\I$ whose layers have empty interior (Kuratowski \S48);
their existence in a decomposable subcontinuum, and Hagopian's Lemma~3 of 1976.
\item \textbf{Chapter~\ref{ch:decomp}: Hagopian's decomposition.} If $M$ is not hereditarily indecomposable, a monotone map of $M$ onto a homogeneous continuum containing an arc
(Hagopian 1984); proofs in outline.
\item \textbf{Chapter~\ref{ch:noarc}: No arc, and hereditary indecomposability.} Jones's arc components and lines, a recurrence argument, a trapping argument in a collar;
Hagopian's theorem.
\item \textbf{Chapter~\ref{ch:span}: Span zero.} Span, Cook's confluence theorem, Davis's theorem, and Oversteegen--Tymchatyn's Theorem~3.6 and Corollary~3.5.
\item \textbf{Chapter~\ref{ch:strips}: The strip theorem.} Oversteegen--Tymchatyn's Theorem~3.2, by a parity argument for closed polygons.
\item \textbf{Chapter~\ref{ch:snake}: Snake-likeness: the plan of Hoehn and Oversteegen.} Theorem~1 from Theorem~20, via Theorem~4 and grid approximations.
\item \textbf{Chapter~\ref{ch:folds}: Simple folds.} The Krasinkiewicz--Minc lemma, Theorem~8 (maps lift through folds), PL graphs, Lemma~6, Proposition~7.
\item \textbf{Chapter~\ref{ch:stairwells}: Stairwells, and Theorem~20 from Theorem~15.} Sides, straight sets, stairwell structures, Propositions~10--19, Lemma~18.
\item \textbf{Chapter~\ref{ch:ladders}: Theorem~15: ladders.} A separator contains a stairwell of odd height, by an explicit construction from a sign function on a grid.
\end{itemize}

The main results are:

\begin{theorem}[Hagopian 1975--1976; Theorems~\ref{thm:hagopian-lemma1} and~\ref{thm:nosep}]
A homogeneous indecomposable plane continuum is atriodic and hereditarily unicoherent, and none of its proper subcontinua separates the plane.
\end{theorem}

\begin{theorem}[Moore's triod theorem; Theorem~\ref{thm:moore-triod}]
Every family of pairwise disjoint triods in the plane is countable.
\end{theorem}

\begin{theorem}[Hagopian 1984; Theorem~\ref{thm:h84}]
A homogeneous indecomposable plane continuum that is not hereditarily indecomposable has a monotone map onto a homogeneous continuum containing an arc, with
proper fibres.
\end{theorem}

\begin{theorem}[Bing 1960, Jones 1975; Theorem~\ref{thm:noarc}]
No homogeneous indecomposable plane continuum contains an arc.
\end{theorem}

\begin{theorem}[Hagopian 1976; Theorem~\ref{thm:HI}]
Every homogeneous indecomposable plane continuum is hereditarily indecomposable.
\end{theorem}

\begin{theorem}[Theorem~C; Theorem~\ref{thm:theoremC}]
Every homogeneous indecomposable plane continuum has span zero.
\end{theorem}

\begin{theorem}[Hoehn--Oversteegen, Theorem~1; Theorem~\ref{thm:ho1}]
A hereditarily indecomposable plane continuum with span zero is snake-like.
\end{theorem}

\begin{theorem}[Hoehn--Oversteegen, Theorem~20; Theorem~\ref{thm:ho20}]
Let $X$ be a hereditarily indecomposable compactum, $f\colon X\to G$ a map to a PL graph, $M\subseteq G\times(0,1)$ a separator of $G\times\set0$ from $G\times\set1$ in $G\times\I$, and $U\supseteq M$
open. Then $f$ is approximated arbitrarily well by $\pi\circ h$ for maps $h\colon X\to U\cap(G\times\I)$.
\end{theorem}

\section{Where these results are used}

Theorem~\ref{thm:indec-pseudoarc} is the indecomposable case of the classification, assembled in Volume~8. Several tools are used again in the decomposable case
(Volume~7): the facts of Chapter~\ref{ch:homog} (atriodicity, hereditary unicoherence, Effros neighbourhoods), type $\lambda$ continua, and the no-arc theorem in
the form of Theorem~\ref{thm:noarc}.

\section{Where the formal proof differs from the literature}\label{sec:intro-diff}

\begin{enumerate}
\item \textbf{Hagopian's papers.} The topological invariance of plane separation, which Hagopian uses without statement, is replaced by \vfour{Theorem~5.22} in the proof that no
proper subcontinuum separates the plane (Theorem~\ref{thm:nosep}). Two theorems quoted from the literature on irreducible continua are false as cited (Remarks~\ref{rem:end-sets-false} and~\ref{rem:thomas21});
the formal development proves counterexamples and uses correct forms. Several of Hagopian's 1984 lemmas are proved by different routes (Chapter~\ref{ch:decomp}).
\item \textbf{Jones's no-arc theorem.} Jones's property (f) cannot be proved from what his argument uses (the dyadic solenoid satisfies all of it), and is replaced by two
properties (R) and (O); a gap in his density argument is replaced by a recurrence argument; the final planar contradiction is a new trapping argument
(Chapter~\ref{ch:noarc}).
\item \textbf{Oversteegen and Tymchatyn.} Several steps of the proofs of Corollary~3.5 and Theorem~3.2 are unargued or impossible as printed (Chapter~\ref{ch:span}). The strip
theorem is proved, in a slightly weaker form that suffices, by a parity argument for closed polygons, without the Jordan curve theorem (Chapter~\ref{ch:strips}).
\item \textbf{Hoehn and Oversteegen.} Theorem~1 as restated on their p.~211, and Theorem~$1'$, are false for a one-point continuum (Remark~\ref{rem:ho1-point}). Folds are built
as subsets of $E\times\R$ rather than as quotients (Chapter~\ref{ch:folds}). Theorem~8 uses a linear push towards fan centres rather than Tietze extensions; the fans must
be disjoint, which is not said. Proposition~7 is proved without boundary bumping and needs a lemma that is not stated. In Proposition~10 a neighbourhood must be shrunk. The
remarks on pp.~193--194 need connectedness. One sentence of the proof of Proposition~19 is false (Remark~\ref{rem:p208}). Claim~15.1 and the irreducible separator in the proof of
Theorem~15 are replaced by ladders (Chapter~\ref{ch:ladders}).
\item \textbf{Proofs in outline.} Chapter~\ref{ch:decomp} and parts of Chapters~\ref{ch:stairwells} and~\ref{ch:ladders} give long case analyses in outline, marked as such. The complete
arguments are in the formal development, at the declarations named in the ``Formal counterpart'' lines.
\end{enumerate}

\section{Prerequisites and conventions}\label{sec:intro-conv}

The reader should know the material of Volume~3 (irreducible continua, the cut-wire theorem, boundary bumping, upper semicontinuous decompositions,
snake-like continua and the $\varepsilon$-map criterion), the Effros theorem (Volume~3), composants (Volume~2), the plane topology of Volume~4 (crossing parity of polygons, Janiszewski's
theorem) and Volume~5 (the Jordan curve theorem, monotone decompositions of the plane, collars).

\begin{itemize}
\item \textbf{Continua.} A continuum is compact, connected and metric; it may be a point unless ``nondegenerate'' is said.
\item \textbf{The plane.} As in Volumes~4 and~5, $\R^2$ carries the sup metric, so that balls are squares.
\item \textbf{Graphs.} In Chapters~\ref{ch:snake}--\ref{ch:ladders}, $E$ is a real normed space, $G\subseteq E$ a PL graph, $G^\star=G\times\I$, and $\pi$ the first projection of $E\times\R$; $\rfr A$ is
the relative frontier of $A$ in the ambient graph.
\item \textbf{Span.} ``Span zero'' means span zero in the sense of Hoehn and Oversteegen (Definition~\ref{def:span}).
\end{itemize}

\paragraph{Formal counterparts.} As in the earlier volumes, many results are followed by a line ``Formal counterpart:'' naming the corresponding declaration of the
Lean development (all names begin with \lean{HPC.}). These lines may be skipped. Appendix~\ref{app:formal} collects them.

\paragraph{Exercises.} Each chapter ends with exercises. None of them is used later.

%======================================================================
\chapter{First consequences of homogeneity in the plane}\label{ch:homog}
%======================================================================

This chapter collects the facts about a homogeneous indecomposable plane continuum that the rest of
the volume uses over and over. They come from the first lemma of Hagopian's 1975 paper
\emph{Homogeneous plane continua} and from the first page of his 1976 paper. All of them rest on one
idea: a homogeneous indecomposable continuum has uncountably many pairwise disjoint composants, each
containing a homeomorphic copy of any given proper subcontinuum, and the plane does not have room for
uncountably many disjoint copies of certain configurations.

\section{Standing conventions}

A \emph{continuum} is a compact connected metric space (or a compact connected subset of a metric
space). The formal development allows a continuum to be a single point; ``nondegenerate'' is said
explicitly when it is needed. A continuum $M$ is \emph{decomposable} if $M=A\cup B$ with $A,B$ proper
subcontinua, \emph{indecomposable} otherwise, and \emph{hereditarily indecomposable} if every
subcontinuum is indecomposable. A space $Y$ is \emph{homogeneous} if for any $x,y\in Y$ there is a
homeomorphism $h\colon Y\to Y$ with $h(x)=y$.

\formal{\lean{HPC.Fold.Q78.IsContinuum}, \lean{HPC.Fold.Q78.IsIndecomposable},
\lean{HPC.Jones.IsDecomposableSpace}, \lean{HPC.Plane.IsHereditarilyIndecomposable},
\lean{HPC.Jones.IsHomogeneous}.}

A compact set $K\subseteq\R^2$ \emph{separates the plane} if $\R^2\setminus K$ is not connected
(\formal{\lean{HPC.Plane.SeparatesPlane}.}) Triods, atriodic continua and (hereditarily) unicoherent
continua are as in \vthree{Definitions~4.1 and~4.2}; in particular a triod is a continuum $T$ with a
subcontinuum $N$ (the core) such that $T\setminus N$ is the union of three nonempty pairwise separated
sets.

Throughout the volume, unless said otherwise, $M\subseteq\R^2$ denotes a \emph{homogeneous
indecomposable plane continuum}, in its own topology; homeomorphisms are homeomorphisms of $M$ onto
itself. We write $\cmp Mp$ for the composant of $p$ in $M$ (\vtwo{Definition~1.1}): the union of the
proper subcontinua of $M$ containing $p$.

\section{Composants}

\begin{theorem}[Moore's theorems on composants]\label{thm:composants}
Let $X$ be a continuum.
\begin{enumerate}[label=(\alph*)]
\item If $X$ is nondegenerate, every composant of $X$ is dense in $X$, and $p\in\cmp Xp$.
\item If $X$ is indecomposable and $H,K$ are proper subcontinua with a common point, $H\cup K$ is a
proper subcontinuum.
\item If $X$ is indecomposable, two composants are equal or disjoint.
\item If $X$ is indecomposable, every proper subcontinuum has empty interior in $X$, and every
composant is a countable union of nowhere dense closed sets (it is meagre).
\item If $X$ is nondegenerate and indecomposable, $X$ has uncountably many composants.
\end{enumerate}
\end{theorem}

\begin{proof}
(a) $\set p$ is a proper subcontinuum, so $p\in\cmp Xp$. Let $U$ be open with $U\cap X\neq\emptyset$. As $X$
is connected and nondegenerate, it has no isolated points, so $U$ contains a point $y\in X$ with $y\ne
p$; choose $s>0$ with $s<d(p,y)$ and $\bar B(y,s)\subseteq U$. By boundary bumping in metric form
(\vthree{Corollary~1.12}) some subcontinuum $K$ of $X\setminus B(y,s)$ contains $p$ and meets $\bar
B(y,s)\subseteq U$. $K$ misses $y$, so it is proper, and $K\subseteq\cmp Xp$ meets $U$.

(b) $H\cup K$ is a continuum (two continua with a common point). If it were $X$, $X$ would be the union of
the two proper subcontinua $H$ and $K$.

(c) Suppose $z\in\cmp Xx\cap\cmp Xy$: there are proper subcontinua $H\ni x,z$ and $K\ni y,z$. If $w\in\cmp
Xx$, take a proper subcontinuum $L\ni x,w$. By (b) twice, $L\cup H\cup K$ is a proper subcontinuum, and it
contains $y$ and $w$; so $w\in\cmp Xy$. By symmetry the composants are equal.

(d) A proper subcontinuum of an indecomposable continuum has empty interior (\vtwo{Theorem~1.7}). Let
$\set{b_n}$ be a countable base of $X$ and $C_n$ the component of $x$ in the closed set $X\setminus b_n$
(when $x\notin b_n$). Each $C_n$ is a subcontinuum missing the nonempty open set $b_n$ (if $b_n\neq\emptyset$), so
it is proper, closed and nowhere dense. If $K$ is a proper subcontinuum containing $x$, pick $y\in X\setminus
K$ and $b_n\ni y$ with $\cl{b_n}\cap K=\emptyset$; then $K\subseteq C_n$. So $\cmp Xx=\bigcup_n C_n$.

(e) $X$ is a compact metric space, hence a Baire space. By (a) and (c), $X$ is the union of its composants;
if there were countably many, by (d) $X$ would be a countable union of nowhere dense closed sets.
\end{proof}

\formal{\lean{HPC.Hagopian.composant_dense'}, \lean{HPC.Hagopian.union_ne_univ_of_indec},
\lean{HPC.Hagopian.composant_eq_or_disjoint'}, \lean{HPC.Hagopian.interior_eq_empty_of_proper},
\lean{HPC.Hagopian.composant_isMeagre'}, \lean{HPC.Hagopian.uncountable_composants'} (in
\lean{HPC/Hagopian/H75Composants.lean}).}

\begin{lemma}[Copies in every composant]\label{lem:copies}
Let $K$ be a proper subcontinuum of $M$ and $x_0\in K$. For every composant $\kappa$ of $M$ and every
$y\in\kappa$ there is a homeomorphism $h$ with $h(x_0)=y$, and then $h(K)$ is a proper subcontinuum of $M$
contained in $\kappa$. Copies $h(K)\subseteq\kappa$ and $h'(K)\subseteq\kappa'$ in different composants are disjoint.
\end{lemma}

\begin{proof}
$h$ exists by homogeneity; $h(K)$ is a proper subcontinuum containing $y$, so $h(K)\subseteq\cmp My=\kappa$.
Different composants are disjoint by Theorem~\ref{thm:composants}(c).
\end{proof}

\formal{\lean{HPC.Hagopian.isContinuum_homeo_image}, \lean{HPC.Hagopian.homeo_image_ne_univ},
\lean{HPC.Hagopian.subset_composant}.}

\section{Effros neighbourhoods}

The Effros theorem in the form of \vthree{Theorem~7.15} is the basic tool for using homogeneity
locally. Hagopian states it in two ways.

\begin{proposition}[Hagopian 1975, Lemma~4; 1976, Lemma~2; 1984, Lemma~1]\label{prop:effros-nbhd}
Let $Y$ be a homogeneous compact metric space, $\varepsilon>0$ and $x\in Y$.
\begin{enumerate}[label=(\alph*)]
\item There is an open set $W\ni x$ such that for all $y,z\in W$ some homeomorphism $h$ of $Y$ has
$h(y)=z$ and $d(v,h(v))<\varepsilon$ for all $v\in Y$.
\item The set $\Eff(x,\varepsilon)$ of points $h(x)$, $h$ a homeomorphism moving no point by $\varepsilon$ or
more (an \emph{$\varepsilon$-homeomorphism}), is open and contains $x$.
\end{enumerate}
\end{proposition}

\begin{proof}
(a) Take $\delta$ for $\varepsilon$ as in \vthree{Theorem~7.15} and $W=B(x,\delta/2)$.

(b) $x\in\Eff(x,\varepsilon)$ by the identity. Let $y=h(x)$ with $h$ moving points by less than $\varepsilon$.
By compactness $\max_v d(v,h(v))=\varepsilon'<\varepsilon$. Take $\delta$ for $\varepsilon-\varepsilon'$ as in
\vthree{Theorem~7.15}. If $d(y,z)<\delta$, some $g$ with $g(y)=z$ moves points by less than
$\varepsilon-\varepsilon'$, and $g\circ h$ takes $x$ to $z$ and moves points by less than $\varepsilon$. So
$B(y,\delta)\subseteq\Eff(x,\varepsilon)$.
\end{proof}

\formal{\lean{HPC.Hagopian.hagopian75_lemma4}, \lean{HPC.Hagopian.hagopian76_lemma2},
\lean{HPC.Hagopian.hagopian84_lemma1}; all from \lean{HPC.Effros.effros_homogeneous}.}

\section{A non-unicoherent continuum separates the plane}

\begin{proposition}[Moore 1962, Theorem~22]\label{prop:moore22}
Let $A,B\subseteq\R^2$ be compact and connected with $A\cap B$ not connected. Then $A\cup B$ separates the
plane.
\end{proposition}

\begin{proof}
This is Eilenberg's argument, with Borsuk's theorem in place of a logarithm on the whole plane. Write
$A\cap B=P\cup Q$ with $P,Q$ nonempty, closed and disjoint. By Urysohn's lemma there is a continuous
$\varphi\colon\R^2\to[0,1]$ with $\varphi=0$ on $P$ and $\varphi=1$ on $Q$. Define $f=e^{2\pi i\varphi}$ on $A$ and
$f=1$ on $B$; on $A\cap B$ both formulas give $1$, so $f$ is continuous and nowhere zero on $A\cup B$. If
$\R^2\setminus(A\cup B)$ were connected, Borsuk's theorem (\vfour{Theorem~2.10}) would give a logarithm $g$
of $f$ on $A\cup B$. On the connected set $A$, $g$ and $2\pi i\varphi$ are logarithms of the same function,
so they differ by a constant (\vfour{Lemma~1.3}); on the connected set $B$, $g$ is a logarithm of $1$,
so it is constant. Comparing at $p\in P$ and $q\in Q$: $g(q)-g(p)=2\pi i(\varphi(q)-\varphi(p))=2\pi i$ from
$A$, and $g(q)-g(p)=0$ from $B$. This is a contradiction.
\end{proof}

\formal{\lean{HPC.Hagopian.separatesPlane_of_not_isPreconnected_inter}.}

\section{Separating continua in uncountably many composants}

The next lemma is the common core of the hereditary unicoherence of $M$ and of the fact that no
proper subcontinuum of $M$ separates the plane.

\begin{lemma}[Disjoint complementary domains]\label{lem:disjoint-domains}
Let $F_1,F_2\subseteq M$ be disjoint closed sets with $F_1\ne\emptyset$, and let $U_i$ be a component of
$\R^2\setminus F_i$ with $U_i\cap M=\emptyset$ ($i=1,2$). Then $U_1\cap U_2=\emptyset$.
\end{lemma}

\begin{proof}
Suppose $U_1\cap U_2\neq\emptyset$. $U_1$ is connected and misses $M\supseteq F_2$, so it lies in a component of
$\R^2\setminus F_2$; it meets $U_2$, so $U_1\subseteq U_2$. Symmetrically $U_2\subseteq U_1$, and $U=U_1=U_2$. The
frontier of a component of the complement of a closed set lies in that set, so $\fr U\subseteq F_1\cap
F_2=\emptyset$. Then $U$ is a nonempty clopen subset of the plane, so $U=\R^2$; but $U$ misses
$F_1\neq\emptyset$.
\end{proof}

\formal{\lean{HPC.Hagopian.disjoint_components_missing}.}

\begin{lemma}\label{lem:sep-composants}
Suppose that every composant $\kappa$ of $M$ contains a proper subcontinuum $S_\kappa$ that separates the
plane. Then $M$ is degenerate or decomposable; that is, under the standing hypotheses this is
impossible.
\end{lemma}

\begin{proof}
Assume $M$ is nondegenerate (it is indecomposable by hypothesis). The $S_\kappa$ lie in distinct composants,
so they are pairwise disjoint (Theorem~\ref{thm:composants}(c)), and there are uncountably many of them
(Theorem~\ref{thm:composants}(e)).

\emph{Step 1: some $S_\kappa$ has all its complementary domains meeting $M$.} Otherwise choose, for every
$\kappa$, a component $U_\kappa$ of $\R^2\setminus S_\kappa$ missing $M$. By Lemma~\ref{lem:disjoint-domains}
the $U_\kappa$ are pairwise disjoint nonempty open sets, uncountably many of them. Each contains a point
with rational coordinates, and these points are distinct: a contradiction.

\emph{Step 2.} Fix such a $\kappa$ and $S=S_\kappa$. As $S$ separates the plane it has two complementary
domains $U\neq V$, both meeting $M$. $U\cap M$ is a nonempty open subset of the Baire space $M$, and $\kappa$
is meagre (Theorem~\ref{thm:composants}(d)); so there is $p\in U\cap M$ with $p\notin\kappa$. The composant
$\cmp Mp$ is dense (Theorem~\ref{thm:composants}(a)), so it contains a point $q$ of the open set $V\cap M$,
and there is a proper subcontinuum $L$ containing $p$ and $q$. $L$ is connected and meets the two
components $U$ and $V$ of $\R^2\setminus S$, so it meets $S$. By Theorem~\ref{thm:composants}(b), $L\cup S$ is
a proper subcontinuum; it contains $p$ and a point of $\kappa$, so $p\in\kappa$. This is a contradiction.
\end{proof}

\formal{The argument appears twice in the formal development: in
\lean{HPC.Hagopian.hereditarilyUnicoherent_of_homogeneous_indecomposable} (\lean{HPC/Hagopian/H75Lemma1.lean})
and in \lean{HPC.Hagopian.no_proper_subcontinuum_separates_aux} (\lean{HPC/Hagopian/H76NoSep.lean}); the
countability step is \lean{HPC.Jones.countable_of_pairwiseDisjoint_open}, the Baire step
\lean{HPC.Hagopian.exists_notMem_of_isMeagre}.}

\section{Hereditary unicoherence and the non-separation theorem}

\begin{theorem}[Hagopian 1975, Lemma~1, second half]\label{thm:hu}
$M$ is hereditarily unicoherent.
\end{theorem}

\begin{proof}
If $M$ is a point there is nothing to prove. Suppose $F=H\cup L$ with $H,L$ subcontinua and $H\cap L$ not
connected. Then $F\ne M$: if $H\cup L=M$, indecomposability forces $H=M$ or $L=M$, and then $H\cap L$ is
$L$ or $H$, which is connected. Fix $x_0\in F$. For each composant $\kappa$ take $y\in\kappa$ and $h$ with
$h(x_0)=y$ (Lemma~\ref{lem:copies}); $S_\kappa=h(F)=h(H)\cup h(L)$ is a proper subcontinuum in $\kappa$, and
$h(H)\cap h(L)=h(H\cap L)$ is not connected. By Proposition~\ref{prop:moore22}, $S_\kappa$ separates the
plane. Lemma~\ref{lem:sep-composants} gives a contradiction.
\end{proof}

\formal{\lean{HPC.Hagopian.hereditarilyUnicoherent_of_homogeneous_indecomposable},
\lean{HPC.Hagopian.hagopian75_lemma1_hu}.}

Note that here the copies $h(F)$ separate the plane for the same reason as $F$: they are again unions of
two continua with a disconnected intersection. For a general separating proper subcontinuum one needs that
separation passes to homeomorphic copies. Volume~4 proves this in full (Borsuk's invariance theorem,
\vfour{Theorem~2.16}) and also in a weaker form that suffices here and is the one used by the formal
proof: \emph{if a compact set $K$ separates the plane and $K'$ is a plane continuum homeomorphic to $K$,
some subcontinuum of $K'$ separates the plane} (\vfour{Theorem~5.22}, by Bing's theorem that a continuum
none of whose subcontinua separates the plane is tree-like).

\begin{theorem}[Hagopian 1976, p.~344]\label{thm:nosep}
No proper subcontinuum of $M$ separates the plane.
\end{theorem}

\begin{proof}
Let $K$ be a proper subcontinuum that separates the plane, and $x_0\in K$. For each composant $\kappa$, a
copy $h(K)\subseteq\kappa$ (Lemma~\ref{lem:copies}) contains a subcontinuum $S_\kappa$ that separates the plane
(\vfour{Theorem~5.22}). $S_\kappa$ is proper, as it lies in $h(K)$. Lemma~\ref{lem:sep-composants} gives a
contradiction.
\end{proof}

\formal{\lean{HPC.Hagopian.exists_separating_subcontinuum_of_homeomorph},
\lean{HPC.Hagopian.separating_in_copy}, \lean{HPC.Hagopian.no_proper_subcontinuum_separates}.}

\begin{remark}
Hagopian's 1976 paper cites ``[1975, Lemma~1, second paragraph of proof]'' for Theorem~\ref{thm:nosep}.
That paragraph treats the non-unicoherent case only; for a general separating subcontinuum the step
``each copy separates the plane'' is the topological invariance of plane separation, which neither paper
states. The formal proof supplies it as described above.
\end{remark}

\section{Moore's triod theorem}

\begin{theorem}[Moore 1962, Theorem~84]\label{thm:moore-triod}
Every family of pairwise disjoint triods in the plane is countable.
\end{theorem}

The proof splits according to whether the cores separate the plane. Both cases are a pigeonhole
argument over countably many ``rational'' choices.

\begin{proposition}[Separating cores; Moore's Theorem~73]\label{prop:moore73}
Let $(K_i)_{i\in I}$ be pairwise disjoint connected subsets of the plane, and $Z_i\subseteq K_i$ closed sets
that separate the plane with $K_i\setminus Z_i\neq\emptyset$. Then $I$ is countable.
\end{proposition}

\begin{proof}
For each $i$ choose $p_i\in K_i\setminus Z_i$, a ball $Q_i\ni p_i$ with rational centre and radius missing $Z_i$,
and a point $o_i$ with rational coordinates in a component $C_i$ of $\R^2\setminus Z_i$ not containing $p_i$
(there are at least two components, and components of open sets are open). Then $Q_i$, which is connected
and misses $Z_i$, lies in the component of $p_i$, so $Q_i\cap C_i=\emptyset$. If $I$ were uncountable, two
indices $i\neq j$ would have $Q_i=Q_j=Q$ and $o_i=o_j=o$. Let $C_i,C_j$ be the components of $o$ in
$\R^2\setminus Z_i$ and $\R^2\setminus Z_j$.

\emph{$C_i\subseteq C_j$.} $K_j$ is connected, misses $K_i\supseteq Z_i$, and contains $p_j\in Q$; so $K_j$ lies in the
component of $\R^2\setminus Z_i$ containing $Q$, which is not $C_i$. Hence $Z_j\subseteq K_j$ misses $C_i$, and
$C_i$ is a connected subset of $\R^2\setminus Z_j$ containing $o$. By symmetry $C_j\subseteq C_i$.

So $C=C_i=C_j$, and $\fr C\subseteq Z_i\cap Z_j=\emptyset$. $C$ is a nonempty clopen proper subset of the
plane (it misses $Q$), which is impossible.
\end{proof}

\formal{\lean{HPC.Hagopian.comp_subset_comp}, \lean{HPC.Hagopian.countable_of_separating}.}

For non-separating cores the tool is an interlacing lemma: two disjoint continua inside a closed curve
cannot meet the curve in interlaced points.

\begin{lemma}[Interlacing; Moore's Theorem~76]\label{lem:interlace}
Let $J\subseteq\R^2$ be closed with $\R^2\setminus J=D\cup E$, where $D,E$ are disjoint open sets, $E$ is connected
and $J\subseteq\cl E$. Let $x$ be a continuum and $y$ a connected set, disjoint, both contained in $D\cup J$.
Suppose $J=\alpha\cup\beta$ with $\alpha,\beta$ compact and $\alpha\cap\beta\subseteq x$. Then $y$ does not meet both
$\alpha\setminus\beta$ and $\beta\setminus\alpha$.
\end{lemma}

\begin{proof}
Suppose $y_1\in y\cap(\alpha\setminus\beta)$ and $y_2\in y\cap(\beta\setminus\alpha)$. Put $F_1=x\cup\alpha$ and $F_2=x\cup\beta$;
these are compact, and $F_1\cap F_2=x\cup(\alpha\cap\beta)=x$ is connected. $E$ misses $F_1$ and $F_2$ (it misses $J$
and $x\subseteq D\cup J$); let $U_k$ be the component of $\R^2\setminus F_k$ containing $E$.

(i) \emph{$U_1\cap U_2\cap(D\setminus x)=\emptyset$.} If $w$ were in it, then for $e\in E$ neither $F_1$ nor $F_2$
would separate $w$ from $e$, and by Janiszewski's theorem (\vfour{Theorem~2.7}) $F_1\cup F_2=x\cup J$ would not.
But a connected set missing $J$ cannot meet both $D$ and $E$.

(ii) \emph{A point $z\in J\setminus x$ lies in exactly one of $U_1$, $U_2$:} in $U_2\setminus U_1$ if $z\in\alpha\setminus\beta$,
and in $U_1\setminus U_2$ if $z\in\beta\setminus\alpha$. Say $z\in\alpha\setminus\beta$. Then $z\in F_1$, so $z\notin U_1$; and $z\notin
F_2$, while $z\in\cl E$, so $z$ is in the closure of the connected set $E\subseteq\R^2\setminus F_2$ and off $F_2$; hence
$z\in U_2$.

(iii) \emph{$y\subseteq U_1\cup U_2$.} Points of $y\cap J$ are in $J\setminus x$ and are covered by (ii). Let $w\in y\cap
D$ and suppose $w\notin U_1\cup U_2$. Let $C$ be the component of $w$ in $\R^2\setminus(x\cup J)$; it is open. If a
point $z\in J\setminus x$ were in $\fr C$, then by (ii) $z$ would lie in $U_1$ or in $U_2$, say in $U_k$; the open
set $U_k$ meets $C$, and $C$ is a connected subset of $\R^2\setminus F_k$, so $C\subseteq U_k$ and $w\in U_k$. Hence
$\fr C\subseteq x$. Now $y$ is connected, misses $x\supseteq\fr C$, and meets the open set $C$, so $y\subseteq C$. But
$y_1\in y\cap J$ and $C\cap J=\emptyset$.

By (i)--(iii), $y\cap U_1$ and $y\cap U_2$ are disjoint relatively open sets covering $y$, containing $y_2$ and $y_1$
respectively. This contradicts the connectedness of $y$.
\end{proof}

\formal{\lean{HPC.Hagopian.interlace}, with \lean{HPC.Hagopian.mem_comp_of_mem_closure},
\lean{HPC.Hagopian.mem_of_mem_frontier_comp}, \lean{HPC.Hagopian.subset_of_disjoint_frontier}.}

For a simple closed polygon $J$, with $D=\inn J$ and $E$ the unbounded complementary domain
(\vfour{Theorem~3.15}), $J$ is the image of $[0,2]$ under a continuous $\gamma$ injective on $[0,2)$ with
$\gamma(2)=\gamma(0)$. If $0\le s<u<t<v<2$ with $\gamma(s),\gamma(t)\in x$ and $\gamma(u),\gamma(v)\in y$, the lemma applies to
$\alpha=\gamma[s,t]$ and $\beta=\gamma([t,2]\cup[0,s])$: $\alpha\cap\beta=\set{\gamma(s),\gamma(t)}\subseteq x$, $\gamma(u)\in\alpha\setminus\beta$
and $\gamma(v)\in\beta\setminus\alpha$. So such parameters do not exist.

\formal{\lean{HPC.Hagopian.exists_loop_param}, \lean{HPC.Hagopian.interlace_param}.}

\begin{lemma}\label{lem:interlace-three}
Let $x,y$ be as above, and suppose there are three pairwise ordered intervals $I_1<I_2<I_3$ in $[0,2)$
such that each $I_k$ contains a parameter $p_k$ with $\gamma(p_k)\in x$ and a parameter $q_k$ with
$\gamma(q_k)\in y$. This is impossible.
\end{lemma}

\begin{proof}
If $p_3<q_3$ then $p_1<q_2<p_3<q_3$ are interlaced parameters of the forbidden kind; if $q_3<p_3$ then
$q_1<p_2<q_3<p_3$ are, with the roles of $x$ and $y$ exchanged.
\end{proof}

\formal{\lean{HPC.Hagopian.interlace_three}.}

\begin{proposition}[Non-separating cores; Moore's Theorems~79, 78 and~76]\label{prop:moore79}
A family of pairwise disjoint triods in the plane whose cores do not separate the plane is countable.
\end{proposition}

\begin{proof}[Proof in outline]
Let $T_i$ be a triod with core $Z_i$ and legs $A_i,B_i,C_i$, and choose $a_i\in A_i$, $b_i\in B_i$, $c_i\in C_i$.
By the enclosing polygon theorem (\vfour{Theorem~4.14}, applied with $G=\R^2\setminus\set{a_i,b_i,c_i}$ and a slight
shrinking) there is a simple closed polygon $J_i$, which may be taken with rational vertices, such that
$Z_i\subseteq\inn J_i$ and $a_i,b_i,c_i\notin\inn J_i\cup J_i$. There are countably many such polygons, so if the
family were uncountable, uncountably many triods would share one polygon $J$, with $D=\inn J$.

For such a triod, boundary bumping applied to the continuum $Z_i\cup A_i$ (and likewise for the other legs)
gives a subcontinuum of $T_i\cap(D\cup J)$ through a point of $Z_i$ that meets $J$ in a point of $A_i$; the union
of the three continua so obtained is a continuum $x_i\subseteq T_i\cap(D\cup J)$ meeting $J$ in a point of each
leg. The three points have distinct parameters on $J$, which lie in three pairwise ordered intervals
with rational end points. Again by countability, two triods $i\neq j$ have the same three intervals.
Then $x_i$ and $x_j$ are disjoint continua in $D\cup J$ to which Lemma~\ref{lem:interlace-three} applies:
a contradiction.
\end{proof}

\formal{\lean{HPC.Hagopian.TriodData}, \lean{HPC.Hagopian.exists_leg},
\lean{HPC.Hagopian.exists_cross_continuum}, \lean{HPC.Hagopian.exists_rat_intervals},
\lean{HPC.Hagopian.countable_of_nonseparating}.}

\begin{proof}[Proof of Theorem~\ref{thm:moore-triod}]
The triods whose cores separate the plane are countably many by Proposition~\ref{prop:moore73} (with
$K_i=T_i$, $Z_i$ the core; $T_i\setminus Z_i$ is nonempty). The others are countably many by
Proposition~\ref{prop:moore79}.
\end{proof}

\formal{\lean{HPC.Hagopian.moore_triod_theorem'}, \lean{HPC.Hagopian.mooreTriodStatement}; the statement as a
proposition is \lean{HPC.Hagopian.MooreTriodStatement}.}

\begin{theorem}[Hagopian 1975, Lemma~1]\label{thm:hagopian-lemma1}
$M$ is atriodic and hereditarily unicoherent.
\end{theorem}

\begin{proof}
Hereditary unicoherence is Theorem~\ref{thm:hu}. Let $T\subseteq M$ be a triod. $T\neq M$, since an
indecomposable continuum is not a triod (\vthree{Corollary~4.7}). Fix $x_0\in T$; for each composant
$\kappa$ choose a copy $h_\kappa(T)\subseteq\kappa$ (Lemma~\ref{lem:copies}). These are pairwise disjoint triods in the
plane, uncountably many by Theorem~\ref{thm:composants}(e), contradicting Theorem~\ref{thm:moore-triod}.
\end{proof}

\formal{\lean{HPC.Hagopian.atriodic_of_homogeneous_indecomposable},
\lean{HPC.Hagopian.not_isTriod_univ_of_indecomposable}, \lean{HPC.Hagopian.isTriodSet_image},
\lean{HPC.Hagopian.hagopian75_lemma1}.}

\section{Two consequences used later}

\begin{corollary}\label{cor:nested-or-arc}
\leavevmode
\begin{enumerate}[label=(\alph*)]
\item If $H,K$ are subcontinua of $M$ with a common point, $H\cap K$ is a continuum.
\item \emph{(The triod argument.)} There are no continua $N,X_1,X_2,X_3\subseteq M$ such that each $X_i$ meets
$N$, $X_i\cap X_j\subseteq N$ for $i\neq j$, and no $X_i$ is contained in $N$.
\end{enumerate}
\end{corollary}

\begin{proof}
(a) $H\cup K$ is a subcontinuum, which is unicoherent by Theorem~\ref{thm:hu}. (b) By
\vthree{Lemma~4.5}, $N\cup X_1\cup X_2\cup X_3$ would be a triod in $M$, against
Theorem~\ref{thm:hagopian-lemma1}.
\end{proof}

\formal{\lean{HPC.Hagopian.isContinuum_inter_of_HU}; \lean{HPC.Plane.isTriod_of_core}.}

Most arguments of the next two chapters end with Corollary~\ref{cor:nested-or-arc}(b): one exhibits a
continuum $N$ and three continua through it that pairwise meet only in $N$, and concludes that the
configuration is impossible.

\section*{Exercises}

\begin{exercise}
Show that in Theorem~\ref{thm:composants}(d) the countable base can be replaced by any countable family of
open sets $b_n$ such that every point $y\ne x$ lies in some $b_n$ with $x\notin\cl{b_n}$.
\end{exercise}

\begin{exercise}
Show that the circle is homogeneous and not hereditarily unicoherent, and explain which hypothesis of
Theorem~\ref{thm:hu} fails.
\end{exercise}

\begin{exercise}
Deduce from Proposition~\ref{prop:moore73} that the plane does not contain uncountably many pairwise disjoint
simple closed curves each with a ``tail'' (a continuum $K_i=J_i\cup A_i$, $A_i$ an arc meeting the simple closed
curve $J_i$ in one end point).
\end{exercise}

%======================================================================
\chapter{Continua of type \texorpdfstring{$\lambda$}{lambda}}\label{ch:typelambda}
%======================================================================

A decomposable irreducible continuum looks, from far away, like an arc: it can often be mapped onto
$[0,1]$ with connected fibres, the \emph{layers}, so that it is ``an arc of layers''. Continua with such
a map whose layers are nowhere dense are said to be of \emph{type $\lambda$} (Hagopian writes ``type
A$'$'', following Thomas). They are Hagopian's main tool: if $M$ is not hereditarily indecomposable,
it contains a continuum of type $\lambda$, and the layers of such continua are rigid enough under small
homeomorphisms of $M$ to produce, in Chapter~\ref{ch:decomp}, a decomposition of $M$ whose quotient
contains an arc.

The first half of the chapter is the general theory, which comes from Kuratowski's \emph{Topology},
Volume~II, \S48 (Thomas's memoir, which Hagopian cites, was not used). The second half proves the
existence of type $\lambda$ continua in $M$ and the invariance of their interior layers under small
homeomorphisms (Hagopian 1975, Theorem~1).

\section{Admissible maps and layers}

A continuum $E$ is \emph{irreducible between $a$ and $b$} if $a,b\in E$ and no proper subcontinuum of $E$
contains both (\vthree{Definition~2.1}).

\formal{\lean{HPC.IsIrreducibleBetween}; for the space itself, \lean{HPC.Hagopian.IrrBetween}.}

\begin{definition}\label{def:admissible}
Let $E$ be a compact metric space. An \emph{admissible map} of $E$ is a continuous surjection
$k\colon E\to[0,1]$ all of whose fibres are connected. The fibres $\lay ks$ are the \emph{layers} of $k$;
$\lay k0$ and $\lay k1$ are the \emph{end layers}, the others are \emph{interior layers}. For $A\subseteq[0,1]$
write $\lay kA$ for the union of the layers over $A$. A continuum $E$ is of \emph{type $\lambda$} if it is
irreducible between two of its points and has an admissible map whose layers all have empty interior in
$E$. Such a pair $(E,k)$ is called \emph{minimal admissible}; following Hagopian, $k$ is ``the
decomposition'' of $E$ (by Theorem~\ref{thm:unique-decomp} below its layers are determined by $E$).
\end{definition}

\formal{\lean{HPC.Hagopian.IsAdmissibleMap}, \lean{HPC.Hagopian.layer}, \lean{HPC.Hagopian.IsTypeA'},
\lean{HPC.Hagopian.IsMinimalAdmissible}, \lean{HPC.Hagopian.layerSet}.}

The fibre decomposition of an admissible map is upper semicontinuous, and its quotient space is
homeomorphic to $[0,1]$ (a continuous bijection from a compact space onto a Hausdorff space); so the map
form of the definition agrees with Hagopian's, who speaks of an upper semicontinuous decomposition with
quotient an arc (\formal{\lean{HPC.Hagopian.admissible_quotient_homeomorph}.})

\begin{lemma}\label{lem:adm-basic}
Let $k$ be an admissible map of the compact space $E$.
\begin{enumerate}[label=(\alph*)]
\item If $S\subseteq[0,1]$ is connected, $\lay kS$ is connected. If moreover $S$ is a nonempty closed interval,
$\lay kS$ is a continuum; in particular every layer is a continuum.
\item If $E$ is irreducible between $a$ and $b$, then $\set{k(a),k(b)}=\set{0,1}$.
\item \emph{(The interval lemma.)} If $E$ is irreducible between two points and $K\subseteq E$ is a subcontinuum
meeting $\lay kd$ and $\lay ke$ with $d<e$, then $\lay ku\subseteq K$ for every $u\in(d,e)$.
\end{enumerate}
\end{lemma}

\begin{proof}
(a) $k$ is a closed map (compact to Hausdorff) with connected fibres. If $\lay kS=P\cup Q$ with $P,Q$
disjoint and relatively closed, every fibre over $S$ lies in $P$ or in $Q$, so $S=k(P)\cup k(Q)$ with $k(P)$
and $k(Q)$ disjoint and relatively closed in $S$ ($k$ is closed); one of them is empty.

(b) Let $s=k(a)$, $t=k(b)$ and $I$ the closed interval between them. $\lay kI$ is a continuum containing $a$
and $b$, so it is $E$, and $I=k(E)=[0,1]$.

(c) By (b) we may assume $k(a)=0$ and $k(b)=1$. The continua $\lay k{[0,d]}$, $K$ and $\lay k{[e,1]}$ form a chain
of continua each meeting the next, so their union is a continuum containing $a$ and $b$, hence equal to
$E$. A layer over $u\in(d,e)$ misses the first and the last, so it lies in $K$.
\end{proof}

\formal{\lean{HPC.Hagopian.IsAdmissibleMap.isPreconnected_preimage}, \lean{HPC.Hagopian.isContinuum_img},
\lean{HPC.Hagopian.isContinuum_layer}, \lean{HPC.Hagopian.end_layers},
\lean{HPC.Hagopian.layer_subset_of_meets'}, \lean{HPC.Hagopian.layer_subset_of_meets}.}

\begin{proposition}[End sets]\label{prop:end-sets}
Let $(E,k)$ be minimal admissible. A subcontinuum of $E$ meeting both end layers is $E$.
\end{proposition}

\begin{proof}
By the interval lemma, $K\supseteq\lay k{(0,1)}$. The end layers are closed with empty interior, so their
union has empty interior, and $\lay k{(0,1)}$ is dense in $E$. As $K$ is closed, $K=E$.
\end{proof}

\formal{\lean{HPC.Hagopian.dense_middle}, \lean{HPC.Hagopian.eq_of_meets_ends},
\lean{HPC.Hagopian.admissible_end_sets'}, \lean{HPC.Hagopian.hagopian76_end_sets}.}

\begin{remark}\label{rem:end-sets-false}
Without the condition on the layers this fails, although it is used in this form in the literature (as
``Thomas, Theorem~8''): on $E=[0,2]$, irreducible between $0$ and $2$, the map $k(x)=\max(0,x-1)$ is
admissible with end layer $\lay k0=[0,1]$, and $K=[1,2]$ meets both end layers. The formal development
proves this counterexample (\formal{\lean{HPC.Hagopian.not_admissible_end_sets}.}) Hagopian only uses
end sets of type $\lambda$ continua, where Proposition~\ref{prop:end-sets} applies.
\end{remark}

\begin{lemma}\label{lem:split-conn}
Let $K$ be a connected space, $K=C\cup D$ with $C,D$ closed, and $S\subseteq C$ connected with $C\cap D\subseteq S$.
Then $C$ is connected.
\end{lemma}

\begin{proof}
If $C=C_1\cup C_2$ with $C_1,C_2$ disjoint closed and $S\subseteq C_1$, then $C_2\cap D\subseteq C_2\cap S=\emptyset$, so
$K=(C_1\cup D)\cup C_2$ is a union of two disjoint closed sets, and $C_2=\emptyset$.
\end{proof}

\formal{\lean{HPC.Hagopian.isPreconnected_of_union_closed}.}

\begin{proposition}[Lemma~2 of Hagopian 1975, necessity]\label{prop:lemma2-nec}
In a continuum $(E,k)$ of type $\lambda$, every subcontinuum $K$ with nonempty interior in $E$ is decomposable.
\end{proposition}

\begin{proof}
Let $d=\min k(K)$ and $e=\max k(K)$. If $d=e$, $K$ lies in a layer, which has empty interior. So $d<e$; pick
$m\in(d,e)$ and put $C=K\cap\lay k{[0,m]}$, $D=K\cap\lay k{[m,1]}$. By the interval lemma
$S=\lay k{(d,m]}\subseteq K$, so $S\subseteq C$, $S$ is connected (Lemma~\ref{lem:adm-basic}(a)), and $C\cap
D\subseteq\lay km\subseteq S$. By Lemma~\ref{lem:split-conn}, $C$ is a continuum, and likewise $D$. Neither is $K$: $C$
misses the points of $K$ in $\lay ke$, and $D$ those in $\lay kd$. So $K=C\cup D$ is decomposable.
\end{proof}

\formal{\lean{HPC.Hagopian.not_isIndecomposable_of_typeA'}.}

\begin{theorem}[Uniqueness of the decomposition; Hagopian 1975, Lemma~3]\label{thm:unique-decomp}
Let $(E,k)$ be minimal admissible and $k'$ any admissible map of $E$. Every layer of $k$ lies in a layer of
$k'$. Consequently two minimal admissible maps of $E$ have the same layers.
\end{theorem}

\begin{proof}
Let $L=\lay kt$, a continuum. If $k'$ took two values $d<e$ on $L$, the interval lemma for $k'$ would give
$L\supseteq\lay{k'}{(d,e)}$, a nonempty open subset of $E$; but $L$ has empty interior. The second statement
follows by applying the first in both directions.
\end{proof}

\formal{\lean{HPC.Hagopian.fibre_subset_fibre}, \lean{HPC.Hagopian.hagopian75_lemma3},
\lean{HPC.Hagopian.fibre_eq_iff}, \lean{HPC.Hagopian.hagopian75_lemma3_unique}.}

So a homeomorphism $h$ carries the decomposition of $E$ to that of $h(E)$, and ``the decomposition'' of a
type $\lambda$ continuum is well defined (up to reversing $[0,1]$).

\section{Kuratowski's family \texorpdfstring{$\mathcal D$}{D}}

The converse of Proposition~\ref{prop:lemma2-nec} is the substantial half of Lemma~2. Throughout this
section $Y$ is a continuum irreducible between $a\ne b$. A set $D$ is a \emph{closed domain} if
$D=\cl{\interior D}$. Let $\mathcal D$ consist of $\emptyset$ and the connected closed domains containing $a$.

\formal{\lean{HPC.Hagopian.IsDElt}.}

\begin{lemma}[Kuratowski \S48 II--III]\label{lem:kur}
\leavevmode
\begin{enumerate}[label=(\alph*)]
\item If $K\ni a$ is a subcontinuum with $b\notin K$, then $\cl{Y\setminus K}$ is a continuum.
\item \emph{(No hole.)} Let $K\ni a$ be a subcontinuum with $b\notin K$ and $U\subseteq K$ a nonempty open set with
$a\notin U$ and $\cl U\cap\interior K\subseteq U$. This is impossible.
\item If $K\ni a$ is a subcontinuum with nonempty interior, then $\cl{\interior K}\in\mathcal D$.
\item $\mathcal D$ is a chain under inclusion; if $D_1\subsetneq D_2$ in $\mathcal D$, then $D_1\subseteq\interior D_2$.
\end{enumerate}
\end{lemma}

\begin{proof}
(a) Let $\cl{Y\setminus K}=P\cup Q$ with $P,Q$ disjoint closed and $b\in P$, and suppose $Q\ne\emptyset$. Since $Y\setminus K$
is dense in $\cl{Y\setminus K}$ and $Q$ is relatively open there, $Q\setminus K\ne\emptyset$. We claim $K\cup P$ is connected. If
$K\cup P=A\cup B$ with $A,B$ disjoint closed and $K\subseteq A$, then $B\subseteq P\setminus K$ misses $Q$, so $Y=(A\cup Q)\cup B$ is a
union of disjoint closed sets and $B=\emptyset$. So $K\cup P$ is a continuum containing $a$ and $b$, hence equals $Y$; but
it misses $Q\setminus K\neq\emptyset$.

(b) The component $C$ of $a$ in the compact set $K\setminus U$ reaches the relative frontier of $K\setminus U$ in $K$
(boundary bumping, \vthree{Theorem~1.10}), which lies in $\cl U\setminus U$; by hypothesis such points are not in
$\interior K$, so they lie in $\cl{Y\setminus K}$. By (a), $C\cup\cl{Y\setminus K}$ is a continuum; it contains $a$ and $b$ and
misses $U$ ($\cl{Y\setminus K}$ misses $U\subseteq\interior K$). This contradicts irreducibility.

(c) If $b\in K$ then $K=Y$. Otherwise let $F=\cl{\interior K}$ and suppose $F=F_1\cup F_2$ with $F_1,F_2$ disjoint
closed, $F_2\ne\emptyset$ and $a\notin F_2$ (if $a\notin F$ take $F_1=\emptyset$). Then $U=\interior K\setminus F_1$ is open,
nonempty ($F_2$ is relatively open in $F=\cl{\interior K}$, so it meets $\interior K$), contained in $K$, misses $a$,
and $\cl U\cap\interior K\subseteq F_2\cap\interior K\subseteq U$. This contradicts (b). So $F$ is connected and contains
$a$.

(d) Let $D_1,D_2\in\mathcal D$ be nonempty, neither equal to $Y$, and suppose $D_1\not\subseteq D_2$ and $D_2\not\subseteq D_1$.
Since $D_1$ is a closed domain not contained in $D_2$, there is $x\in\interior D_1\setminus D_2$; likewise $y\in\interior
D_2\setminus D_1$. The continuum $D_1\cup\cl{Y\setminus D_2}$ (it is connected by (a), and $x$ is a common point) contains
$a$ and $b$, so it is $Y$; but $y$ is in neither set. For the second statement, suppose $z\in D_1\setminus\interior
D_2$, so $z\in\cl{Y\setminus D_2}$; again $D_1\cup\cl{Y\setminus D_2}=Y$, while a point of $\interior D_2\setminus D_1$ lies in
neither.
\end{proof}

\formal{\lean{HPC.Hagopian.irr_compl_continuum}, \lean{HPC.Hagopian.irr_no_hole},
\lean{HPC.Hagopian.isDElt_closure_interior}, \lean{HPC.Hagopian.IsDElt.chain},
\lean{HPC.Hagopian.IsDElt.subset_interior}.}

\begin{proposition}[No jumps; Kuratowski \S48 VII, Theorems~1--2]\label{prop:no-jumps}
If $D_1\subsetneq D_2$ in $\mathcal D$, then $\cl{D_2\setminus D_1}$ is a continuum. If every subcontinuum of $Y$ with nonempty
interior is decomposable, then some member of $\mathcal D$ lies strictly between $D_1$ and $D_2$.
\end{proposition}

\begin{proof}[Proof in outline]
The first statement is a separation argument of the kind used in Lemma~\ref{lem:kur}(a), applied to
$\cl{Y\setminus D_1}$, $\cl{Y\setminus D_2}$ and the decomposition $D_2=D_1\cup\cl{D_2\setminus D_1}$. For the second, if nothing
lay strictly between, the continuum $\cl{D_2\setminus D_1}$, which has interior $\interior D_2\setminus D_1\ne\emptyset$, would be
decomposable, $\cl{D_2\setminus D_1}=H\cup K$ with $H,K$ proper subcontinua; one of them, say $H$, meets $D_1$, and then
$\cl{\interior(D_1\cup H)}$ is a member of $\mathcal D$ (Lemma~\ref{lem:kur}(c)) strictly between $D_1$ and $D_2$.
\end{proof}

\formal{\lean{HPC.Hagopian.IsDElt.closure_diff}, \lean{HPC.Hagopian.IsDElt.exists_between}.}

\begin{theorem}[From $\mathcal D$ to a map; Kuratowski \S48 IV and VII, Theorems~3--5]\label{thm:scale}
If $Y$ is irreducible between $a$ and $b$ and every subcontinuum of $Y$ with nonempty interior is
decomposable, then $Y$ has an admissible map all of whose layers have empty interior.
\end{theorem}

\begin{proof}
\emph{A countable dense subfamily.} Let $\set{b_n}$ be a countable base. For each pair $(m,n)$ choose, if there is one,
a member of $\mathcal D$ whose interior contains $b_m$ and which misses $b_n$; the countably many chosen sets form a family
$\mathcal C\subseteq\mathcal D$ that is \emph{dense}: strictly between two members $D_1\subsetneq D_2$ of $\mathcal D$ lies a member of
$\mathcal C$. (By Proposition~\ref{prop:no-jumps} there are $D_1\subsetneq D\subsetneq D'\subsetneq D_2$; basic sets $b_m\subseteq\interior
D\setminus D_1$ and $b_n\subseteq\interior D_2\setminus D'$ exist by Lemma~\ref{lem:kur}(d), and the chosen member for $(m,n)$
lies between $D_1$ and $D_2$ because $\mathcal D$ is a chain.)

\emph{A scale.} With $\emptyset$ and $Y$ removed, $\mathcal C$ is a countable dense linear order without end points, hence
(Cantor) order-isomorphic to $\Q\cap(0,1)$. This gives a family $(E_q)_{q\in\Q}$ in $\mathcal D$, increasing, strictly so on
$(0,1)$, with $E_q=\emptyset$ for $q\le0$, $E_q=Y$ for $q\ge1$, and dense in $\mathcal D$. By Lemma~\ref{lem:kur}(d),
$E_q\subseteq\interior E_{q'}$ for $q<q'$.

\emph{The map.} Put $k(y)=\inf\set{q:y\in E_q}\in[0,1]$. Then $k(y)<t$ iff $y\in E_q$ for some rational $q<t$, and
$k(y)>t$ iff $y\notin E_q$ for some rational $q>t$; using $E_q\subseteq\interior E_{q'}$, both sets are open, so $k$ is
continuous. It is onto because the scale is strictly increasing and $Y$ is connected. The fibre over $t$ is the
intersection of the continua $\cl{E_{q'}\setminus E_q}$, $q<t<q'$ (Proposition~\ref{prop:no-jumps}), a directed family, so it
is connected. If a fibre had interior, the closure of $\bigcup_{q<t}E_q$ and the closure of the interior of
$\bigcap_{q>t}E_q$ would be two members of $\mathcal D$, the first strictly inside the second; by density two
members of the scale would lie between them, and both would have index $t$, which is absurd.
\end{proof}

\formal{\lean{HPC.Hagopian.exists_countable_dense_D}, \lean{HPC.Hagopian.IsDScale},
\lean{HPC.Hagopian.exists_scale}, \lean{HPC.Hagopian.scaleFun}, \lean{HPC.Hagopian.continuous_scaleFun},
\lean{HPC.Hagopian.isPreconnected_scaleFun_fibre}, \lean{HPC.Hagopian.interior_scaleFun_fibre},
\lean{HPC.Hagopian.scaleMap_spec}, \lean{HPC.Hagopian.exists_typeA'_map}.}

\begin{remark}
Kuratowski's Theorem~3 states that $\mathcal D$ has the order type of an interval. The formal proof only uses
the order type of the countable dense subfamily, which is that of the rationals by Cantor's theorem.
\end{remark}

\begin{theorem}[Hagopian 1975, Lemma~2]\label{thm:lemma2}
A continuum $E$ irreducible between two points is of type $\lambda$ if and only if every subcontinuum of $E$ with
nonempty interior (relative to $E$) is decomposable.
\end{theorem}

\begin{proof}
Necessity is Proposition~\ref{prop:lemma2-nec}; sufficiency is Theorem~\ref{thm:scale}, applied to $E$ as a
space.
\end{proof}

\formal{\lean{HPC.Hagopian.isTypeA'_of_decomposable}, \lean{HPC.Hagopian.hagopian75_lemma2}.}

\begin{remark}\label{rem:thomas21}
Hagopian's paragraph~4 of the proof of his Theorem~1 quotes a theorem (Thomas, Theorem~21; Kuratowski
\S48 VII, Theorem~6 in the form cited): \emph{if every interior layer of a type $\lambda$ continuum is degenerate,
the continuum is an arc}. As stated this is false: the $\sin(1/x)$-continuum
$\set{(x,\sin(1/x)):0<x\le1}\cup(\set0\times[-1,1])$ is irreducible between $(1,\sin1)$ and $(0,0)$, the first
coordinate is an admissible map with nowhere dense layers, its interior layers are points, but its end layer
over $0$ is a segment. The formal development proves this counterexample, and the two correct forms: if
\emph{all} layers are points, $k$ is a homeomorphism onto $[0,1]$; if all interior layers are points, $E$
contains an arc. The classification does not use paragraph~4.
\end{remark}

\formal{\lean{HPC.Hagopian.not_arc_of_degenerate_layers}, \lean{HPC.Hagopian.arc_of_degenerate_layers'},
\lean{HPC.Hagopian.exists_arc_of_degenerate_layers}.}

\section{Type \texorpdfstring{$\lambda$}{lambda} continua in \texorpdfstring{$M$}{M}}

We return to the standing hypotheses: $M$ is a homogeneous indecomposable plane continuum. By
Theorem~\ref{thm:hagopian-lemma1} it is atriodic and hereditarily unicoherent.

\begin{theorem}[Hagopian 1975, Theorem~1, paragraphs~1--3]\label{thm:exists-lambda}
Every decomposable subcontinuum $A$ of $M$ contains a continuum of type $\lambda$.
\end{theorem}

\begin{proof}
Write $A=B\cup C$ with $B,C$ proper subcontinua, and choose $b\in B\setminus C$, $c\in C\setminus B$. Let $E\subseteq A$ be a
subcontinuum irreducible between $b$ and $c$ (\vthree{Theorem~2.5}). By Theorem~\ref{thm:lemma2} it suffices to
show that no indecomposable subcontinuum $I$ of $E$ contains a nonempty set $Q=O\cap E$ with $O$ open.

Suppose $I$ is one. The sets $I\cap B$ and $I\cap C$ are continua or empty (Corollary~\ref{cor:nested-or-arc}(a)) and
their union is $I$; as $I$ is indecomposable, $I\subseteq B$ or $I\subseteq C$, say $I\subseteq B$. Then $c\notin I$, and $c\notin Q$.
By boundary bumping there is a continuum $F\subseteq E\setminus O$ containing $c$ and meeting $\cl Q\subseteq I$ in a point
$x$.

Fix $q\in Q$ and $\varepsilon=\min(\dist(c,B),\dist(q,F))>0$, and let $\delta$ be given by the Effros theorem
(\vthree{Theorem~7.15}) for $\varepsilon$. The composants of the indecomposable continuum $I$ are meagre in $I$
(Theorem~\ref{thm:composants}(d)), so the ball $B(x,\delta)$ contains points $x_1,x_2$ of $I$ such that $x,x_1,x_2$ lie
in three different composants of $I$. Let $f,g$ be $\varepsilon$-homeomorphisms with $f(x)=x_1$, $g(x)=x_2$. Then
$f(c),g(c)\notin B$ and $q\notin f(F)\cup g(F)$.

\emph{$F$, $f(F)$, $g(F)$ are pairwise disjoint.} If, say, $F\cap f(F)\ne\emptyset$, then $P=F\cup f(F)$ is a continuum,
and $P\cap I$ is a continuum (Corollary~\ref{cor:nested-or-arc}(a)) containing $x$ and $f(x)$, which lie in different
composants of $I$. So $P\cap I=I\ni q$; but $q\notin F$ (as $F\cap O=\emptyset$) and $q\notin f(F)$. The other pairs are
similar.

\emph{The triod.} Take $N=B$ and $X_1=F$, $X_2=f(F)$, $X_3=g(F)$. Each meets $B$ (at $x$, $f(x)$, $g(x)$, points of $I\subseteq
B$), none lies in $B$ (they contain $c$, $f(c)$, $g(c)$), and they are pairwise disjoint. This contradicts
Corollary~\ref{cor:nested-or-arc}(b).
\end{proof}

\formal{\lean{HPC.Hagopian.no_indecomposable_interior_aux}, \lean{HPC.Hagopian.hagopian75_no_indecomposable_interior},
\lean{HPC.Hagopian.hagopian75_exists_typeA'}.}

The disjointness of $F$, $f(F)$, $g(F)$ is not argued in Hagopian's text; the formal proof supplies the argument
above.

\section{Interior layers under small homeomorphisms}

Let $(E,k)$ be minimal admissible in $M$. The next theorem says that a small homeomorphism of $M$ that moves an
interior layer onto itself at one point moves it onto itself.

\begin{theorem}[Hagopian 1975, Theorem~1, paragraphs~9--11]\label{thm:layer-invariant}
Let $0\le r<s<t\le1$, and let $\varepsilon>0$ be at most the distance from $\lay k{[r,t]}$ to $\lay k0\cup\lay k1$. If $h$ is
a homeomorphism of $M$ moving no point by $\varepsilon$ or more, and $h(\lay ks)$ meets $\lay ks$, then $h(\lay ks)=\lay
ks$.
\end{theorem}

\begin{proof}
Write $Y=\lay ks$ and $Z=\lay k{[r,t]}$; note that $\varepsilon>0$ forces $0<r$ and $t<1$.

\emph{Step 1: $h(Z)\subseteq E$.} $h(Z)$ misses the end layers (by the choice of $\varepsilon$) and meets $E$. Suppose
$h(Z)\not\subseteq E$. Then $J=E\cap h(Z)$ is a continuum (Corollary~\ref{cor:nested-or-arc}(a)) with $k(J)=[d',e']\subseteq(0,1)$.
Put $N=\lay k{[d',e']}$, $X_1=\lay k{[0,d']}$, $X_2=\lay k{[e',1]}$ and $X_3=h(Z)\cup N$. These are continua; $X_1,X_2$ meet $N$
in the layers over $d'$ and $e'$; $X_3\cap X_1$ and $X_3\cap X_2$ lie in $N$ because $h(Z)\cap E=J\subseteq N$; and no $X_i$ lies in
$N$ ($X_3$ because $h(Z)\not\subseteq E$). This contradicts Corollary~\ref{cor:nested-or-arc}(b).

\emph{Step 2: $h(Y)\subseteq Y$.} Suppose not. Then $k(h(Y))=[d,e]$ with $d<e$ (it contains $s$), and by the interval lemma
every layer over $(d,e)$ lies in $h(Y)$. Now $E\setminus Y$ is covered by $\lay k{[0,r]}$, $\lay k{[t,1]}$, $\lay k{[r,s)}$ and
$\lay k{(s,t]}$. The images of the first two are compact and miss $h(Y)$. The images of the last two are connected
subsets of $E$ (Step 1) that miss $\lay k{(d,e)}\subseteq h(Y)$, so each lies in $\lay k{[0,d]}$ or in $\lay k{[e,1]}$. Hence a
point $y$ with $k(y)\in(d,e)$ has a neighbourhood that meets $h(E)$ only in $h(Y)$. Pulling it back by $h$, $Y$ has
nonempty interior in $E$, which is false.

\emph{Step 3.} $h^{-1}$ also moves no point by $\varepsilon$ or more and $h^{-1}(Y)$ meets $Y$; by Step~2, $h^{-1}(Y)\subseteq Y$,
that is $Y\subseteq h(Y)$.
\end{proof}

\formal{\lean{HPC.Hagopian.homeo_image_subset_of_meets}, \lean{HPC.Hagopian.homeo_image_layer_subset},
\lean{HPC.Hagopian.hagopian75_layer_invariant_aux}, \lean{HPC.Hagopian.hagopian75_layer_invariant}.}

\begin{remark}
Combining Theorem~\ref{thm:layer-invariant} with Proposition~\ref{prop:effros-nbhd} along a chain of Effros
neighbourhoods inside a layer shows that every interior layer of $E$ is homogeneous (Hagopian's paragraphs~5--8 and~12;
\formal{\lean{HPC.Hagopian.layer_homogeneous_aux}, \lean{HPC.Hagopian.hagopian75_layer_homogeneous}.}) This is not used
in the classification.
\end{remark}

The form in which Theorem~\ref{thm:layer-invariant} is used later replaces the smallness of $h$ by a condition on
where $h$ sends the end sets.

\begin{theorem}[Hagopian 1976, Lemma~3]\label{thm:h76-lemma3}
Let $(E,k)$ and $(F,d)$ be minimal admissible in $M$, $0<s<1$, and $h$ a homeomorphism of $M$. Let $Y=\lay ks$ and $U=\lay
du$. Suppose that $U$ meets $h(Y)$, that $U$ misses $h(\lay k0\cup\lay k1)$, and that $h(Y)$ misses $\lay d0\cup\lay d1$. Then
$h(Y)=U$.
\end{theorem}

\begin{proof}[Proof in outline]
The hypothesis $h(Y)\cap(\lay d0\cup\lay d1)=\emptyset$ replaces the smallness of $h$. The triod of Step~1 above puts $h(Y)$
inside $F$; then $J=F\cap h(E)$ is a continuum. Let $[\gamma,\delta]=k(h^{-1}(J))$. If $\gamma=s$ (or $\delta=s$), the continuum
$h(\lay k{[0,s]})$ (or $h(\lay k{[s,1]})$) meets $F$ only in $h(Y)$, so by the same triod it lies in $F$, which puts $h(\lay
k0)$ (or $h(\lay k1)$) into the layer of $F$ containing $h(Y)$, against the hypotheses. If $\gamma<s<\delta$ and $h(Y)\ne U$,
then $h(Y)$ contains an open subset of $F$ and Step~2 above gives $Y$ nonempty interior in $E$. So $h(Y)\subseteq U$; the same
argument for $h^{-1}$, $U$ and $Y$ gives $U\subseteq h(Y)$.
\end{proof}

\formal{\lean{HPC.Hagopian.hagopian76_lemma3}, proved in \lean{HPC/Hagopian/H76Lemma3.lean}.}

\section*{Exercises}

\begin{exercise}
Show that an arc is of type $\lambda$, and that its decomposition is into points. Show that the
$\sin(1/x)$-continuum is of type $\lambda$, and find its decomposition.
\end{exercise}

\begin{exercise}
Show that a continuum of type $\lambda$ is not indecomposable. (Use Proposition~\ref{prop:lemma2-nec} with $K=E$.)
\end{exercise}

\begin{exercise}
Let $(E,k)$ be minimal admissible and $0<a<b<1$. Show that $\lay k{[a,b]}$ is of type $\lambda$, and that its decomposition
consists of the layers of $k$ over $[a,b]$.
\end{exercise}

%======================================================================
\chapter{Hagopian's decomposition}\label{ch:decomp}
%======================================================================

This chapter proves the following theorem, which is the first half of the proof that $M$ is hereditarily
indecomposable (the second half, Chapter~\ref{ch:noarc}, shows that the quotient it produces cannot exist).

\begin{theorem}[Hagopian 1984, Theorem~1, with $N=M$]\label{thm:h84}
Let $M$ be a homogeneous indecomposable plane continuum that is not hereditarily indecomposable. There is a
continuous surjection $q\colon M\to Y$ onto a compact metric space $Y$, with connected fibres none of which is $M$,
such that $Y$ is homogeneous and contains an arc.
\end{theorem}

\formal{\lean{HPC.Hagopian.Hagopian84Decomp} (the statement), \lean{HPC.Hagopian.hagopian84_decomp} (the proof).}

The route is Hagopian's 1984 paper \emph{Atriodic homogeneous continua}, which simplifies the decomposition of his
1976 paper and works for atriodic homogeneous continua in general. Since the argument is long and consists mostly of
careful bookkeeping with chains of open sets, this chapter gives the definitions and statements in full and the proofs
in outline; the complete arguments are in the Lean files \lean{HPC/Hagopian/H84*.lean} and
\lean{HPC/Hagopian/Hagopian84*.lean}. Throughout, $M$ is as in the standing hypotheses of Chapter~\ref{ch:homog},
and ``minimal admissible'' and ``layer'' are as in Chapter~\ref{ch:typelambda}.

\section{Rigidity of layers}

The following facts express that type $\lambda$ continua in $M$ fit together like arcs in a one-dimensional
manifold: where two of them overlap essentially, their layers agree.

\begin{lemma}\label{lem:rigid}
Let $(G,g)$ be minimal admissible in $M$.
\begin{enumerate}[label=(\alph*)]
\item \emph{(Staying inside.)} A subcontinuum $C$ of $M$ that meets $G$ and misses the end layers of $G$ lies in $G$.
\item \emph{(No arm at an interior layer.)} If $0<u<1$ and $A$ is a continuum with $\emptyset\ne A\cap G\subseteq\lay gu$, then
$A\subseteq\lay gu$.
\end{enumerate}
\end{lemma}

\begin{proof}
(a) is Step~1 of the proof of Theorem~\ref{thm:layer-invariant}, with $C$ in place of $h(Z)$. (b) Otherwise $\lay
g{[0,u]}$, $\lay g{[u,1]}$ and $A\cup\lay gu$ are three continua through the core $\lay gu$ that pairwise meet only in it,
and none lies in it; Corollary~\ref{cor:nested-or-arc}(b).
\end{proof}

\formal{\lean{HPC.Hagopian.subset_of_meets_avoids}, \lean{HPC.Hagopian.no_arm_at_interior_layer}.}

\begin{theorem}[Continuity of the decomposition; Hagopian 1976, Lemma~5]\label{thm:h76-lemma5}
If $(E,k)$ is minimal admissible in $M$, the decomposition of $E$ into layers is continuous: the map $s\mapsto\lay ks$ is
continuous for the Hausdorff metric.
\end{theorem}

\begin{proof}[Proof in outline]
Upper semicontinuity holds for every admissible map. Lower semicontinuity at an interior layer comes from
Theorem~\ref{thm:h76-lemma3}: if layers $\lay k{s_n}$ accumulated only on part of $\lay ks$, small homeomorphisms
taking points of $\lay ks$ to nearby points of $\lay k{s_n}$ (Proposition~\ref{prop:effros-nbhd}) would map $\lay ks$
onto $\lay k{s_n}$. At an end layer, Hagopian passes to a larger type $\lambda$ continuum in which it becomes an
interior layer: an \emph{extension} of $E$ away from the other end set (Hagopian 1976, Lemma~4: the extensions of
$E$ away from an end set are linearly ordered by inclusion and have no largest element; the formal proof supplies the
argument for ``no largest element'', which Hagopian leaves unargued). The layers of $E$ are layers of the extension.
\end{proof}

\formal{\lean{HPC.Hagopian.IsExtensionAway}, \lean{HPC.Hagopian.hagopian76_lemma4_chain},
\lean{HPC.Hagopian.hagopian76_lemma4}, \lean{HPC.Hagopian.hagopian76_lemma5}.}

\begin{corollary}[Restriction]\label{cor:restrict}
If $(L,d)$ is minimal admissible and $a<b$, then $\lay d{[a,b]}$ is of type $\lambda$, with a minimal admissible map
whose layers are the layers $\lay d{a+t(b-a)}$; its end sets are $\lay da$ and $\lay db$.
\end{corollary}

\formal{\lean{HPC.Hagopian.minimalAdmissible_restrict}.}

A subcontinuum $H$ of a type $\lambda$ continuum $E$ is \emph{essential} in $E$ if it meets more than one layer of $E$.

\begin{theorem}[Hagopian 1984, Lemma~6]\label{thm:h84-lemma6}
If $H$ is an essential subcontinuum of $(E,k)$, then $H=\lay k{[a,b]}$ for some $a<b$, $H$ is of type $\lambda$, and every
layer of the decomposition of $H$ is a layer of $k$.
\end{theorem}

\begin{proof}
By the interval lemma, $H\supseteq\lay k{(a,b)}$ where $[a,b]=k(H)$. By continuity of the decomposition
(Theorem~\ref{thm:h76-lemma5}) the layers over $a$ and $b$ are limits of layers over $(a,b)$, so $H=\lay k{[a,b]}$;
Corollary~\ref{cor:restrict} gives the minimal admissible map, and uniqueness (Theorem~\ref{thm:unique-decomp}) the last
statement.
\end{proof}

\formal{\lean{HPC.Hagopian.IsEssential84}, \lean{HPC.Hagopian.hagopian84_lemma6}.}

This replaces Thomas's Theorem~8, which Hagopian cites.

\begin{theorem}[Rigidity; Hagopian 1984, Lemma~7]\label{thm:layers-agree}
Let $(E,k)$ and $(G,d)$ be minimal admissible in $M$, and suppose $E\cap G$ meets two layers of $E$. Then either $E$ lies in a
single layer of $G$, or $E\cap G=\lay k{[a,b]}=\lay d{[a',b']}$ with $a<b$ and $a'<b'$, every layer of $k$ meeting $G$ is a layer
of $d$, and every layer of $d$ meeting $E$ is a layer of $k$. In particular (Hagopian's Lemma~7), if neither of $E,G$
contains the other and $G$ meets a layer $\lay kr$, then $\lay kr$ is a layer of $d$.
\end{theorem}

\begin{proof}[Proof in outline]
$K=E\cap G$ is a continuum (Corollary~\ref{cor:nested-or-arc}(a)). If $K$ lies in one layer of $G$ and $E\not\subseteq G$, then
$E$ is an arm at that layer and Lemma~\ref{lem:rigid}(b) applies (end layers are handled by passing to extensions,
Lemma~4 of 1976, in which they become interior). If $K$ meets two layers of each, Theorem~\ref{thm:h84-lemma6} applied in
$E$ and in $G$ shows that $K$ is a union of layers of both, and that the two decompositions of $K$ coincide by
uniqueness. The remaining case ($K$ meets two layers of $E$ and lies in one layer $Y$ of $G$) is excluded by passing to a
two-sided extension of $G$ in which $Y$ is interior, and applying Lemma~\ref{lem:rigid}(b) again. This route differs from
Hagopian's printed proof.
\end{proof}

\formal{\lean{HPC.Hagopian.layers_agree}, \lean{HPC.Hagopian.hagopian84_lemma7_aux}.}

\begin{remark}
Hagopian 1976, Lemma~6, and 1984, Lemma~11 (with $N=M$): the closure of the union of the extensions of $E$ away from an end
set is indecomposable (\formal{\lean{HPC.Hagopian.hagopian76_lemma6_X}, \lean{HPC.Hagopian.hagopian84_lemma11}.}) In
the 1976 paper this union plays the role that the decomposition $\mathfrak D$ below plays in 1984.
\end{remark}

\section{Chains through which a continuum runs straight}

We work with chains of five open subsets of $M$, indexed $L_0,\dots,L_4$ (Hagopian's $L_1,\dots,L_5$).

\begin{definition}\label{def:chain5}
A \emph{5-chain} is a sequence $L=(L_0,\dots,L_4)$ of open subsets of $M$ with $L_i\cap L_j\ne\emptyset$ iff $|i-j|\le1$. Its
\emph{middle} is $L'=L_1\cup L_2\cup L_3$. It is
\begin{itemize}
\item \emph{normal} if $\cl{L_i}\cap\cl{L_j}=\emptyset$ for $|i-j|>1$;
\item \emph{free} if $\fr\bigcup_iL_i\subseteq\fr(L_0\cup L_4)\setminus\cl{L'}$ (frontiers and closures in $M$).
\end{itemize}
A minimal admissible $(E,k)$ \emph{runs straight through} $L$ if (1) $E\subseteq\bigcup_iL_i$; (2) $\lay k0\subseteq L_0\setminus L_1$
and $\lay k1\subseteq L_4\setminus L_3$; (3) a layer meeting $\fr L_i$ lies in $\fr L_i$; (4) if $r<t$ and $\lay kr\cup\lay kt\subseteq
L_i$, then $\lay k{[r,t]}\subseteq L_i$. The chain is \emph{regular} if for each component $K$ of $L'$, $\cl K$ (a
\emph{strand}) is of type $\lambda$ and runs straight through $L$. An \emph{ordered refinement} of $L$ is a 5-chain $P$ with
$P_i\subseteq L_i$.
\end{definition}

\formal{\lean{HPC.Hagopian.IsChain5}, \lean{HPC.Hagopian.midUnion}, \lean{HPC.Hagopian.IsNormal5},
\lean{HPC.Hagopian.IsFree5}, \lean{HPC.Hagopian.RunsStraight}, \lean{HPC.Hagopian.IsRegular5},
\lean{HPC.Hagopian.IsOrderedRefinement5}.}

The chains are built from Effros neighbourhoods of pieces of $E$: for $S\subseteq M$ and $\delta>0$ let $\Eff(S,\delta)$ be the
set of points $h(x)$, $x\in S$, $h$ a $\delta$-homeomorphism; it is open by Proposition~\ref{prop:effros-nbhd}.

\begin{lemma}[Folds]\label{lem:fold}
Let $V_0,\dots,V_4$ be open subsets of $[0,1]$ covering it, with $V_i\cap V_j=\emptyset$ for $|i-j|>1$, $0\in V_0\setminus V_1$ and
$1\in V_4\setminus V_3$. If some $V_i$ is not an interval, there are $r<s<t$ and adjacent indices $p,q$ with $s\in V_p$ and
$r,t\in V_q\setminus V_p$ (a \emph{fold}).
\end{lemma}

\begin{proof}
Sweep from left to right. Suppose there is no fold. If $a\in V_p$, no point left of $a$ lies in a later set and no point $\ge
a$ in an earlier one, then $V_p$ is an interval, and the first point after $a$ outside $V_p$ lies in $V_{p+1}$, with the
same two properties. Starting with $a=0$, $p=0$, every $V_p$ is an interval. (If $V_p$ had a gap, the last point of the
complement before a point of $V_p$ and the first after it would lie in neighbours of $V_p$, and the two properties force both
into the same neighbour: a fold.)
\end{proof}

\formal{\lean{HPC.Hagopian.exists_fold}.}

\begin{theorem}[Hagopian 1984, Lemmas~12 and~13]\label{thm:h84-lemma13}
Let $(E,k)$ be minimal admissible in $M$, $\varepsilon>0$, $A$ a closed set missing $E$, and $B$ an open set containing $\lay k0\cup\lay
k1$. Then $E$ runs straight through a free normal regular 5-chain $P$ in $M\setminus A$ with $\cl{P_0\cup P_4}\subseteq B$, such that
for each component $K$ of the middle $P'$ there is an $\varepsilon$-homeomorphism $h$ for which $\cl K$ is an essential
subcontinuum of $h(E)$.
\end{theorem}

\begin{proof}[Proof in outline]
\emph{Lemma~13, the chain.} Choose $0=x_0<x_1<\dots<x_{15}=1$ with $\lay k{[0,x_3]}$ and $\lay k{[x_{12},1]}$ inside $B$, and let
$L_i=\Eff(\lay k{[x_{3i},x_{3i+3}]},\delta)$. For small $\delta$ these form a normal chain in $M\setminus A$ (pieces of non-adjacent
links are disjoint compact sets), and every $\delta$-translate $h(E)$ satisfies (1) and (2). Condition (3) uses the Effros
theorem with the room $\delta-\max_vd(v,h(v))$; condition (4) uses Lemma~\ref{lem:fold}: a non-interval among the parameter
sets $\set{s: h(\lay ks)\subseteq L_i}$ gives a fold, and folds are excluded by atriodicity
(Lemma~\ref{lem:rigid}(a)) and the gap bounds $\dist(\lay k{[0,x_j]},\lay k{[x_{j+1},1]})$. Two translates that meet have
their layers matched by Theorem~\ref{thm:h76-lemma3}. Regularity: a component $K$ of the middle contains a point $h(v)$,
$v\in E$; $\cl K\subseteq h(E)$ by Lemma~\ref{lem:rigid}(a), and it is essential in $h(E)$; by
Theorem~\ref{thm:h84-lemma6} it is of type $\lambda$ and runs straight.

\emph{Lemma~12, freeness.} Put $A'=L_0\setminus\cl{L_1}$ and $B'=L_4\setminus\cl{L_3}$. Moore's theorem on separating closed sets
(\vthree{Theorem~1.9}) applied in $M\setminus(A'\cup B')$ splits off the points that cannot be joined to $E$, and removing
them from the middle links gives an ordered refinement $P$ that is free, normal and regular, through which $E$ still runs
straight, and whose middle components are components of $L'$. The hypothesis of Moore's theorem comes from atriodicity
directly.
\end{proof}

\formal{\lean{HPC.Hagopian.WsetOf}, \lean{HPC.Hagopian.chainLinks84}, \lean{HPC.Hagopian.chainLinks84_chain},
\lean{HPC.Hagopian.chainLinks84_runsStraight}, \lean{HPC.Hagopian.chainLinks84_regular},
\lean{HPC.Hagopian.hagopian84_lemma12}, \lean{HPC.Hagopian.hagopian84_lemma13}.}

\section{The layers of a chain}

For a regular 5-chain $P$ let $\Omega(P)$ be the set of all layers of all strands of $P$, and $\Delta(P)$ the set of those
that miss $\cl{P_0\cup P_4}$.

\formal{\lean{HPC.Hagopian.Omega5}, \lean{HPC.Hagopian.Delta5}.}

\begin{proposition}\label{prop:delta}
Let $P$ be a regular 5-chain.
\begin{enumerate}[label=(\alph*)]
\item $\Delta(P)$ is a partition of $\bigcup_iP_i\setminus\cl{P_0\cup P_4}$.
\item If $P$ is free, there is $\eta>0$ with $d(z,y)\ge\eta$ for $z\in\bigcup\Delta(P)$ and $y\notin P'$.
\item \emph{(Lemma~14.)} If $P$ is free, $\Delta(P)$ is continuous: members through nearby points are Hausdorff close.
\item \emph{(Lemma~15.)} If $P$ is free and normal, there is $\varepsilon>0$ such that every $\varepsilon$-homeomorphism maps each member of
$\Delta(P)$ onto a member of $\Omega(P)$ inside $P'$.
\item \emph{(The arc.)} If $P$ is normal, every strand has a closed interval of layers $\lay c{[a,b]}$, $a<b$, all in $\Delta(P)$.
\end{enumerate}
\end{proposition}

\begin{proof}[Proof in outline]
(a) A layer through a point outside $\cl{P_i}$ misses $\cl{P_i}$ by (3); two members that meet come from the same strand,
where the decomposition is unique. (b) $\cl{\bigcup\Delta(P)}\subseteq P'$ by freeness, then compactness. (c) For $\varepsilon'=\min(\varepsilon,\eta)$,
an $\varepsilon'$-homeomorphism carrying $z$ to $z'$ maps the member through $z$ onto the member through $z'$ by
Theorem~\ref{thm:h76-lemma3} (its hypotheses are checked with $\eta$). (d) The image $h(Z)$ lies in one component $K'$ of
the middle; $h(\cl K)\cap\cl{K'}$ meets two layers of $h(\cl K)$ by Lemma~\ref{lem:rigid}(b), and
Theorem~\ref{thm:layers-agree} makes $h(Z)$ a layer of $\cl{K'}$. (e) The strand meets $P_2$, and the layer through such a
point lies in $P_2$, hence misses $\cl{P_0\cup P_4}$ by normality; the parameters of layers inside $P_2$ form an open set.
\end{proof}

\formal{\lean{HPC.Hagopian.delta5_partition}, \lean{HPC.Hagopian.delta5_dist_pos},
\lean{HPC.Hagopian.hagopian84_lemma14}, \lean{HPC.Hagopian.hagopian84_lemma15},
\lean{HPC.Hagopian.exists_interval_in_delta}.}

\begin{theorem}[Lemma~16$'$]\label{thm:lemma16p}
Let $P$ be a free normal regular 5-chain and $(F,d)$ minimal admissible with $\lay d0\cap P'=\emptyset$. Every layer of $F$ that meets
$P'$ is a layer of the strand of $P$ through any of its points in $P'$.
\end{theorem}

\begin{proof}[Proof in outline]
By Theorem~\ref{thm:layers-agree}, either the layer of $F$ is a layer of the strand, or the strand lies in a single layer of
$F$. The second alternative is impossible: following the layers of $F$ from the end layer $\lay d0$ (outside $P'$) to the
strand, a limit point argument in $\cl{P_2}$ produces an arm at an interior layer.
\end{proof}

\formal{\lean{HPC.Hagopian.strand_layer_mem_delta}, \lean{HPC.Hagopian.not_strand_subset_extension_layer},
\lean{HPC.Hagopian.strand_layer_of_extension}; Hagopian's Lemma~16 as printed is
\lean{HPC.Hagopian.hagopian84_lemma16}.}

\section{The decomposition}

Fix a type $\lambda$ continuum in $M$ (Theorem~\ref{thm:exists-lambda}) and a free normal regular chain $P$ through which it runs
straight (Theorem~\ref{thm:h84-lemma13} with $A=\emptyset$, $B=M$). The members of $\Delta(P)$ cover $\bigcup_iP_i$ except near the
end links; the rest of $M$ is covered by further chains.

\begin{proposition}[Bridges and the chains $P(Z)$]\label{prop:bridges}
Let $Z\in\Omega(P)\setminus\Delta(P)$. Then $Z$ lies in $\cl{P_0}$ or in $\cl{P_4}$ and misses $\cl{P_2}$. There is a type $\lambda$
continuum $J(Z)$ (a \emph{bridge}) containing $Z$ as a layer, missing $\cl{P_2}$, with both end layers in $\Delta(P)$. By
Theorem~\ref{thm:h84-lemma13} with $A=\cl{P_2}$ and $B=\bigcup\Delta(P)$, $J(Z)$ runs straight through a free normal regular chain
$P(Z)$ in $M\setminus\cl{P_2}$ with $\cl{P_0(Z)\cup P_4(Z)}\subseteq\bigcup\Delta(P)$, and then $Z\in\Delta(P(Z))$.
\end{proposition}

\formal{\lean{HPC.Hagopian.omega_not_delta_side}, \lean{HPC.Hagopian.exists_bridge84},
\lean{HPC.Hagopian.mem_delta_of_bridge_chain}, \lean{HPC.Hagopian.exists_chainZ84}.}

\begin{proposition}[The finite cover]\label{prop:finite-cover}
There are finitely many $Z_1,\dots,Z_n\in\Omega(P)\setminus\Delta(P)$ with
$M=\bigcup\Delta(P)\cup\bigcup_{i=1}^n\bigcup\Delta(P(Z_i))$.
\end{proposition}

\begin{proof}
$C=\cl{\bigcup\Delta(P)}\setminus\bigcup\Delta(P)$ is compact. Each point of $C$ lies in $P'\cap\cl{P_0\cup P_4}$, so the layer $Z$
through it of its strand is in $\Omega(P)\setminus\Delta(P)$, and the point lies in the open set $\bigcup\Delta(P(Z))$; finitely many of
these cover $C$. The union $U$ of $\bigcup\Delta(P)$ with them is open; it is closed because
$\cl{\bigcup\Delta(P(Z))}\subseteq\bigcup\Delta(P(Z))\cup\cl{P_0(Z)\cup P_4(Z)}$ (freeness) and $\cl{P_0(Z)\cup P_4(Z)}\subseteq\bigcup\Delta(P)$. As
$M$ is connected, $U=M$.
\end{proof}

\formal{\lean{HPC.Hagopian.exists_finite_cover84}, \lean{HPC.Hagopian.hagopian84_cover},
\lean{HPC.Hagopian.CoverData84}, \lean{HPC.Hagopian.exists_coverData84}.}

Hagopian argues instead that the union of the whole (uncountable) family is closed; the formal proof uses the finite form
above, and deduces the printed statement from it (\formal{\lean{HPC.Hagopian.sUnion_cover84_eq_univ}.})

Put $P(Z_0)=P$ and let $\mathfrak D=\bigcup_{i=0}^n\Delta(P(Z_i))$.

\begin{theorem}\label{thm:frakD}
\leavevmode
\begin{enumerate}[label=(\alph*)]
\item \emph{(Transfer.)} If $Z\in\Omega(P)$ contains a point $z$ of the middle of some $Q=P(Z_i)$, then $Z$ is the layer through $z$ of
the strand of $Q$ through $z$; if moreover $z\notin\cl{Q_0\cup Q_4}$, then $Z\in\Delta(Q)$.
\item $\mathfrak D$ is a partition of $M$ into proper subcontinua.
\item \emph{(Claims (5) and (6).)} $\Omega(P)\subseteq\mathfrak D$, and there is $\delta>0$ such that every $\delta$-homeomorphism maps members of
$\mathfrak D$ onto members of $\mathfrak D$.
\item $\mathfrak D$ is upper semicontinuous, and some compact $C\subseteq M$ has a continuous map onto $[0,1]$ whose fibres are the traces of
members of $\mathfrak D$ (an arc of members).
\end{enumerate}
\end{theorem}

\begin{proof}[Proof in outline]
(a) $Z=c^{-1}(r)$ for the strand map $c$; the strand meets $P_2$ in a layer $c^{-1}(s)$, which misses $\bigcup Q\ni z$ (as $Q$ lies in
$M\setminus\cl{P_2}$), so $s\ne r$; the piece of the strand between these two layers is of type $\lambda$
(Corollary~\ref{cor:restrict}) with one end layer missing $Q'$ and the other $Z$; apply Theorem~\ref{thm:lemma16p} to $Q$.
(b) Every point is covered (Proposition~\ref{prop:finite-cover}); two members that meet are equal, by (a) and
Proposition~\ref{prop:delta}(a), and in the last case by Theorem~\ref{thm:h76-lemma3}. (c) A member of $\Omega(P)\setminus\Delta(P)$ lies in
$\cl{P_0}$ or $\cl{P_4}$, so a point of it lies in some $\bigcup\Delta(P(Z_i))$; then (a). For (6), take $\delta$ below the constants of
Proposition~\ref{prop:delta}(b),(d) for each of the finitely many chains, and use (a) again. (d) Upper semicontinuity from
Proposition~\ref{prop:delta}(c); the arc from Proposition~\ref{prop:delta}(e).
\end{proof}

\formal{\lean{HPC.Hagopian.omega_layer_of_mem_mid}, \lean{HPC.Hagopian.omega_mem_delta_of_mem_mid},
\lean{HPC.Hagopian.CoverData84.members}, \lean{HPC.Hagopian.CoverData84.members_disjoint},
\lean{HPC.Hagopian.CoverData84.omega_subset_members}, \lean{HPC.Hagopian.CoverData84.invariant},
\lean{HPC.Hagopian.CoverData84.isJonesDecomp}.}

\section{The quotient}

Call an equivalence relation on $M$ a \emph{Jones decomposition} if it is upper semicontinuous, its classes are connected and different
from $M$, there is $\delta>0$ such that every $\delta$-homeomorphism respects it, and it has an arc of classes as in
Theorem~\ref{thm:frakD}(d). By Theorem~\ref{thm:frakD}, $\mathfrak D$ is one.

\formal{\lean{HPC.Hagopian.IsJonesDecomp84}.}

\begin{proof}[Proof of Theorem~\ref{thm:h84}]
Let $q\colon M\to Y=M/\mathfrak D$ be the quotient map. $Y$ is compact and metrizable (\vthree{Theorem~6.5}), $q$ is continuous and onto
with connected fibres, none equal to $M$.

\emph{$Y$ is homogeneous.} Let $\mathcal G$ be the set of finite compositions of $\delta$-homeomorphisms. Each $g\in\mathcal G$ respects $\mathfrak
D$ and so induces a homeomorphism of $Y$ (\vthree{Proposition~6.8}). The orbit of a point $x$ under $\mathcal G$ is open (it contains
$\Eff(y,\delta)$ around each of its points $y$, Proposition~\ref{prop:effros-nbhd}), and the orbits partition $M$, so each orbit is also
closed. As $M$ is connected, $\mathcal G$ acts transitively, and so does the group of induced homeomorphisms of $Y$.

\emph{$Y$ contains an arc.} The map $C\to[0,1]$ of Theorem~\ref{thm:frakD}(d) factors through $q(C)$ as a continuous bijection of a
compact space onto $[0,1]$.
\end{proof}

\formal{\lean{HPC.Hagopian.exists_comp_small_homeo}, \lean{HPC.Hagopian.hagopian84Decomp_of_package},
\lean{HPC.Hagopian.hagopian84_decomp}.}

\begin{remark}
In Hagopian's text the chain of points used for homogeneity, the positivity of several constants, the closedness of the union
$\bigcup\mathfrak B$, the choice of $\varepsilon$ in Lemma~15 and the existence of the arc in $M/\mathfrak D$ are stated without argument;
the formal proof supplies them as indicated. Several printed lemmas (7, 12, 15, 16) are proved by routes different from the
printed ones, and Lemma~16 is used in the stronger form Lemma~16$'$.
\end{remark}

\section*{Exercises}

\begin{exercise}
Prove Lemma~\ref{lem:fold} in detail, and show by an example that the conclusion can fail if $V_0$ and $V_2$ are allowed to meet.
\end{exercise}

\begin{exercise}
Let $Y$ be a compact metric space and $\mathcal G$ a set of homeomorphisms of $Y$ closed under composition such that the $\mathcal
G$-orbit of every point contains a neighbourhood of the point. Show that if $Y$ is connected, $\mathcal G$ acts transitively.
\end{exercise}

%======================================================================
\chapter{No arc, and hereditary indecomposability}\label{ch:noarc}
%======================================================================

This chapter proves that a homogeneous indecomposable plane continuum contains no arc (Bing 1960 for the
indecomposable case; F.~B.~Jones, \emph{Use of a new technique in homogeneous continua}, 1975), and deduces from
it and Chapter~\ref{ch:decomp} Hagopian's theorem that such a continuum is hereditarily indecomposable.

Jones's paper proves the no-arc theorem from six properties (a)--(f) of a homogeneous indecomposable plane
continuum containing an arc. The formal proof keeps his overall idea (study the arc components, which are lines,
and use the Effros theorem to transport pieces of a line along itself) but replaces two of his steps, whose
printed proofs are defective, and replaces his final planar argument by a different one.

\begin{itemize}
\item Jones's property (f) (``no folding in a thin neighbourhood'') is proved from (a)--(c) and the Effros theorem
alone. That argument cannot be right: the dyadic solenoid satisfies everything it uses, and (f) fails there. The
formal proof replaces (f) by two properties, (R) and (O) below, which do hold in that generality.
\item Jones's argument that every point of a line is a limit of the rest of the line has a gap (the Effros
homeomorphism might reverse the line); it is replaced by a short recurrence argument.
\item The final planar contradiction is a ``trapping'' argument in a collar of an arc, in the spirit of the
Poincar\'e--Bendixson theorem, using the Jordan curve theorem.
\end{itemize}

\section{The setting}

\begin{definition}\label{def:jones-setting}
A \emph{Jones setting} is a compact metric space $Y$ that is homogeneous, hereditarily unicoherent and atriodic and
contains an arc.
\end{definition}

\formal{\lean{HPC.Hagopian.JonesSetting}, \lean{HPC.Hagopian.jonesSetting_of_plane}.}

By Theorem~\ref{thm:hagopian-lemma1}, a homogeneous indecomposable plane continuum that contains an arc is a Jones
setting. Sections~\ref{sec:lines}--\ref{sec:recurrence} work in an arbitrary Jones setting $Y$; planarity enters
only in Section~\ref{sec:trap}. (The dyadic solenoid is a Jones setting, so these sections are not vacuous.)

\section{Arcs and lines}\label{sec:lines}

\begin{proposition}[Jones (c)]\label{prop:jones-c}
Let $A,B$ be arcs in $Y$ that meet. Then $A\cup B$ is an arc. If $A$ and $B$ have a common end point and meet in another
point, one of them contains the other.
\end{proposition}

\begin{proof}
$A\cap B$ is a continuum (hereditary unicoherence), hence a point or a sub-arc of each. At each end point of $A\cap B$
count the arms of $A\cup B$ there: the sub-arc $A\cap B$ (if nondegenerate), a tail of $A$, a tail of $B$. Three arms that
pairwise meet only at that point form a triod (\vthree{Lemma~4.5}), which is excluded. So there is at most one tail at
each end (at most two in all if $A\cap B$ is a point), and $A\cup B$ is an arc. For the second statement, $A\cap B$ is a sub-arc
from the common end point $a$; its far end is an end point of $A$ or of $B$, since otherwise there would be three arms
there.
\end{proof}

\formal{\lean{HPC.Hagopian.JonesSetting.inter_isPreconnected}, \lean{HPC.Hagopian.JonesSetting.no_three_arms},
\lean{HPC.Hagopian.JonesSetting.union_arc}, \lean{HPC.Hagopian.JonesSetting.nested}.}

In particular two arcs with the same end points are equal. By homogeneity every point is an interior point of an arc
(map an interior point of a given arc to it).

\formal{\lean{HPC.Hagopian.JonesSetting.arc_unique}, \lean{HPC.Hagopian.JonesSetting.exists_arc_through}.}

The \emph{arc component} of $x$ is the union of $\set x$ and all arcs containing $x$. A \emph{line} is a continuous
injection $\ell\colon\R\to Y$ whose image is an arc component.

\formal{\lean{HPC.Hagopian.arcComp}, \lean{HPC.Hagopian.IsLine}.}

\begin{theorem}\label{thm:lines}
\leavevmode
\begin{enumerate}[label=(\alph*)]
\item Every arc component of $Y$ is the image of a line.
\item If $\ell$ is a line and $s<t$, then $\ell[s,t]$ is an arc, and every arc between two points of $\ell$ is of this form.
\item If $\ell,\ell'$ are lines and $h$ is a homeomorphism of $Y$ with $h(\ell(0))\in\ell'(\R)$, then $h\circ\ell=\ell'\circ\psi$ for a
continuous strictly monotone bijection $\psi$ of $\R$.
\end{enumerate}
\end{theorem}

\begin{proof}
(a) Fix an arc $[a,c]$ with $x$ in its interior. The arcs from $x$ containing $[x,c]$ form a chain under inclusion
(Proposition~\ref{prop:jones-c}). There is no largest one: if $[x,z]$ were largest, an arc with $z$ in its interior would
extend it beyond $z$. The union $R^+$ of the chain is separable; a countable dense subset is covered by countably many
arcs of the chain, and these are cofinal (if all lay in one arc $[x,y]$, $R^+$ would lie in the closed set $[x,y]$, which
would be largest). Parametrize $[0,\infty)$ piecewise along an increasing cofinal sequence. Do the same for $R^-$ through
$a$. $R^+\cap R^-=\set x$, since an arc from $x$ containing $[x,a]$ and $[x,c]$ would make them nested, and $R^+\cup R^-$ is
the arc component by Proposition~\ref{prop:jones-c}. Gluing the two parametrizations gives a continuous injection of $\R$.

(b) The first statement holds for any continuous injection; the second because two arcs with the same end points are
equal.

(c) $h$ maps arc components onto arc components, so it maps $\ell(\R)$ onto $\ell'(\R)$; $\psi=\ell'^{-1}\circ h\circ\ell$ is a bijection of $\R$
that maps compact intervals onto compact intervals by (b), hence is continuous and strictly monotone.
\end{proof}

\formal{\lean{HPC.Hagopian.JonesSetting.rayArcs_total}, \lean{HPC.Hagopian.JonesSetting.extend},
\lean{HPC.Hagopian.JonesSetting.exists_ray_seq}, \lean{HPC.Hagopian.JonesSetting.exists_line},
\lean{HPC.Hagopian.JonesSetting.line_arc}, \lean{HPC.Hagopian.JonesSetting.arc_in_line},
\lean{HPC.Hagopian.JonesSetting.homeo_on_line}.}

\section{No small reversal}\label{sec:reversal}

\begin{lemma}[(R0)]\label{lem:R0}
Let $A=[a,p]$ and $C=[p,c]$ be arcs and $\rho=\min(\dist(a,C),\dist(c,A))$. If $h$ is a homeomorphism with $h(p)=p$ that moves every
point by less than $\rho$, then $h(A)\cap C=\set p$.
\end{lemma}

\begin{proof}
Suppose $h(A)$ meets $C\setminus\set p$. The arcs $h(A)$ and $C$ start at $p$ and meet elsewhere, so they are nested
(Proposition~\ref{prop:jones-c}). If $C\subseteq h(A)$, then $c=h(y)$ with $y\in A$ and $d(c,y)<\rho\le\dist(c,A)$; if $h(A)\subseteq C$,
then $h(a)\in C$ and $d(a,h(a))<\rho\le\dist(a,C)$. Both are impossible.
\end{proof}

\formal{\lean{HPC.Hagopian.JonesSetting.no_reversal_arc}.}

\begin{theorem}[(R) and (O)]\label{thm:RO}
There is $\varepsilon_0>0$ with the following properties.
\begin{enumerate}[label=(\alph*)]
\item \emph{(R)} If $h$ is an $\varepsilon_0$-homeomorphism fixing $\ell(s)$ for a line $\ell$, then $h$ maps $\ell(t)$, $t>s$, to points
$\ell(t')$ with $t'>s$.
\item \emph{(O)} If $h$ is an $\varepsilon_0$-homeomorphism with $h(\ell(0))\in\ell(\R)$, then $h\circ\ell=\ell\circ\psi$ with $\psi$ strictly
increasing, continuous and onto.
\end{enumerate}
\end{theorem}

\begin{proof}
(a) For each $q\in Y$ take an arc with $q$ in its interior and its $\rho_q$ as in Lemma~\ref{lem:R0}, and let $\delta_q$ be the
Effros bound for $\rho_q/3$. If $d(b,q)<\delta_q$, let $g$ be a $\rho_q/3$-homeomorphism with $g(b)=q$. If $h$ fixes $b$, moves points
less than $\rho_q/3$ and reverses the line through $b$, then $ghg^{-1}$ fixes $q$, moves points less than $\rho_q$, and reverses
the line through $q$; by Theorem~\ref{thm:lines}(c) reversing the line at a fixed point means reversing every arc there,
which contradicts Lemma~\ref{lem:R0}. Cover $Y$ by finitely many balls $B(q,\delta_q)$ and take $\varepsilon_0=\min\rho_q/3$.

(b) By Theorem~\ref{thm:lines}(c), $h\circ\ell=\ell\circ\psi$ with $\psi$ a continuous monotone bijection. If $\psi$ decreased, $\psi(s)-s$ would be
continuous, strictly decreasing and onto $\R$, so $\psi(s^*)=s^*$ for some $s^*$, and $h$ would fix $\ell(s^*)$ and reverse $\ell$
there, against (a).
\end{proof}

\formal{\lean{HPC.Hagopian.JonesSetting.local_no_reversal}, \lean{HPC.Hagopian.JonesSetting.exists_no_reversal},
\lean{HPC.Hagopian.JonesSetting.orientation}.}

\begin{remark}
Jones's (f) asserts, for a point $p$ and an open set $D$ with $p\notin\cl D$, an $\varepsilon$ such that no arc $abc$ with $a,c\in D$ has
the $\varepsilon$-neighbourhood of $ab$ containing $p$ and $bc$. His proof uses only the Effros theorem and (c). In the dyadic
solenoid, where all of these hold, (f) fails: a long initial piece $\lambda[0,S]$ of a line is $\varepsilon/2$-dense, and an arc $abc$
along the line beyond it is a counterexample. Theorem~\ref{thm:RO} is what the argument actually needs.
\end{remark}

\section{Recurrence}\label{sec:recurrence}

For a line $\ell$, its \emph{forward limit set} is $\omega(\ell)=\set{z: z=\lim\ell(u_n)\text{ for some }u_n\to\infty}$.

\formal{\lean{HPC.Hagopian.omegaLim}, \lean{HPC.Hagopian.mem_omegaLim_iff}.}

\begin{theorem}\label{thm:recurrent}
\leavevmode
\begin{enumerate}[label=(\alph*)]
\item $\omega(\ell)$ contains the arc component of each of its points.
\item Every line lies in its own forward limit set: $\ell(\R)\subseteq\omega(\ell)$. The same holds backwards.
\end{enumerate}
\end{theorem}

\begin{proof}
(a) Let $z=\lim\ell(u_n)$, $u_n\to\infty$, and $\lambda$ the line through $z$ with $\lambda(0)=z$. The Effros theorem gives
$\varepsilon_n$-homeomorphisms $h_n$, $\varepsilon_n\to0$, with $h_n(z)=\ell(u_n)$; by Theorem~\ref{thm:lines}(c), $h_n(\lambda(s))=\ell(\psi_n(s))$
with $\psi_n(0)=u_n$. Fix $s$. If $\psi_n(s)\le C$ along a subsequence, then for every $L$, eventually $\ell[C,C+L]\subseteq h_n(\lambda[0,s])$ (the
parameter interval between $\psi_n(s)$ and $u_n$ eventually covers $[C,C+L]$), which lies within $\varepsilon_n$ of $\lambda[0,s]$. So
$\ell[C,\infty)\subseteq\lambda[0,s]$, a compact arc, which is impossible (the arc components of $\ell$ and $\lambda$ are equal or disjoint, and
$\ell$ is injective). Hence $\psi_n(s)\to\infty$ and $\lambda(s)\in\omega(\ell)$.

(b) $K=\omega(\ell)$ is nonempty and compact. Suppose $\ell(0)\notin K$ and let $\varepsilon=\frac12\dist(\ell(0),K)$. Take $p\in K$, $\delta$ the Effros
bound for $\varepsilon$, $t$ with $d(\ell(t),p)<\delta$, and an $\varepsilon$-homeomorphism $g$ with $g(\ell(t))=p$. $g$ maps the line $\ell$ onto the
line through $p$, which lies in $K$ by (a); so $g(\ell(0))\in K$, while $d(\ell(0),g(\ell(0)))<\varepsilon<\dist(\ell(0),K)$. Apply this to the shifted
lines $s\mapsto\ell(s+c)$ and the reversed line $s\mapsto\ell(-s)$.
\end{proof}

\formal{\lean{HPC.Hagopian.JonesSetting.omegaLim_saturated}, \lean{HPC.Hagopian.JonesSetting.zero_mem_omegaLim},
\lean{HPC.Hagopian.JonesSetting.line_recurrent}, \lean{HPC.Hagopian.JonesSetting.jones75_d}.}

\section{Passes through a box}\label{sec:passes}

From now on $M$ is a homogeneous indecomposable plane continuum containing an arc, $\ell$ is a line in $M$, and
$\gamma=\ell|_{[0,1]}$. By the collar theorem (\vfive{Theorem~9.19}) there is a \emph{collar} $F$ of $\gamma$: an embedding of
$[0,1]\times[-1,1]$ into the plane with $F(t,0)=\gamma(t)$ and $F((0,1)\times(-1,1))$ open. On $F([0,1]\times[-1,1])$ we use the collar coordinates
$(X,Y)=F^{-1}$. The \emph{flipped} collar $\bar F(x,y)=F(x,-y)$ is again a collar; everything below that is stated ``above'' the arc
holds ``below'' by flipping.

\formal{\lean{HPC.Hagopian.IsCollar}, \lean{HPC.Hagopian.exists_collar}, \lean{HPC.Hagopian.collarCoord},
\lean{HPC.Hagopian.flipC}.}

Fix $\kappa=1/32$. The \emph{middle band} is $X\in[\frac14+\kappa,\frac34-\kappa]=[\frac9{32},\frac{23}{32}]$; it contains the column
$[\frac38,\frac58]$. For $\eta\le1$ let $Q_\eta=F([\frac14,\frac34]\times[-\eta,\eta])$ be the \emph{box} of height $\eta$.

\begin{definition}\label{def:pass}
A \emph{pass} of $Q_\eta$ is a piece $\ell[a,b]$, $a<b$, with $\ell(a)$ on the left edge $F(\set{\frac14}\times[-\eta,\eta])$, $\ell(b)$ on the right
edge $F(\set{\frac34}\times[-\eta,\eta])$, and $\ell((a,b))$ in the open box $F((\frac14,\frac34)\times(-\eta,\eta))$. A pass has property
\emph{(Mon)} if $X(\ell(v))>X(\ell(v'))-2\kappa$ whenever $a\le v'<v\le b$. It is \emph{above} (below) the arc if it lies, end points
included, in $Y>0$ ($Y<0$).
\end{definition}

\formal{\lean{HPC.Hagopian.IsPass}, \lean{HPC.Hagopian.PassMon}, \lean{HPC.Hagopian.PassAbove},
\lean{HPC.Hagopian.PassBelow}.}

On a pass with (Mon), if $X(\ell v)=c$, $X(\ell v')=c'$ and $c'-c\ge2\kappa$, then $v<v'$ (the \emph{order corollary}).

\begin{lemma}[Passes]\label{lem:passes}
\leavevmode
\begin{enumerate}[label=(P\arabic*)]
\item Distinct passes of $Q_\eta$ have disjoint closed parameter intervals.
\item Only finitely many passes meet a bounded parameter interval; so among the passes with a given property starting after
a given time, if any, there is a first one.
\item $\ell[\frac14,\frac34]$ is a pass (the \emph{base pass}).
\item Every other pass is above or below the arc.
\end{enumerate}
For $\eta'\le\eta$, every pass of $Q_{\eta'}$ is a pass of $Q_\eta$.
\end{lemma}

\begin{proof}
(P1) If $a<a'<b$, then $\ell(a')$ would be on the left edge and in the open box; if $a'=b$, on both edges; if $a=a'$ and $b<b'$,
$\ell(b)$ would be on the right edge and in the open box. (P2) Passes inside $[-N,N]$ have length at least the modulus of uniform
continuity of $\ell$ for the distance between the edges, and at most two straddle $\pm N$. (P3) is clear. (P4) A pass other than the
base pass misses the base arc (injectivity of $\ell$ and (P1)), and its preimage under $F$ is connected.
\end{proof}

\formal{\lean{HPC.Hagopian.IsPass.disjoint}, \lean{HPC.Hagopian.IsPass.finite}, \lean{HPC.Hagopian.exists_first_pass},
\lean{HPC.Hagopian.isPass_base}, \lean{HPC.Hagopian.IsPass.above_or_below}, \lean{HPC.Hagopian.IsPass.mono}.}

\begin{lemma}[The shadow lemma]\label{lem:shadow}
For every $\eta'\in(0,1]$ there is $\delta>0$ such that: if $x\in[\frac9{32},\frac{23}{32}]$ and $d(\ell(u),\ell(x))<\delta$, then $u$ lies in the open
parameter interval of a pass of $Q_{\eta'}$ with (Mon).
\end{lemma}

\begin{proof}
Take $\varepsilon<\varepsilon_0$ (Theorem~\ref{thm:RO}) such that every point within $\varepsilon$ of $\gamma(s)$, $s\in[\kappa,1-\kappa]$, lies in the open collar with
$\abs{X-s}<\kappa$ and $\abs Y<\eta'$ (uniform continuity of $F^{-1}$ near the compact arc), and let $\delta$ be the Effros bound for $\varepsilon$.
\begin{enumerate}
\item The Effros theorem gives $h$ with $h(\ell(x))=\ell(u)$ moving points less than $\varepsilon$; by (O), applied to the shifted line, $h\circ\ell=\ell\circ\psi$
with $\psi$ increasing, continuous and onto, $\psi(x)=u$. Let $\theta=\psi^{-1}$.
\item Put $s_1=\frac3{16}$, $s_2=\frac{13}{16}$. For $v\in[\psi(s_1),\psi(s_2)]$ the point $\ell(v)=h(\gamma(\theta(v)))$ lies in the open collar with
$\abs{X(\ell v)-\theta(v)}<\kappa$ and $\abs Y<\eta'$. So $X<\frac14$ at $\psi(s_1)$ and $X>\frac34$ at $\psi(s_2)$.
\item Let $a=\max\set{v\in[\psi(s_1),u]:X(\ell v)\le\frac14}$ and $b=\min\set{v\in[u,\psi(s_2)]:X(\ell v)\ge\frac34}$. On $(a,u]$, $\frac14<X<\theta(v)+\kappa\le
x+\kappa\le\frac34$; on $[u,b)$, $\frac34>X>\theta(v)-\kappa\ge x-\kappa\ge\frac14$. Hence $a<u<b$, $X(\ell a)=\frac14$, $X(\ell b)=\frac34$, and
$\ell[a,b]$ is a pass of $Q_{\eta'}$.
\item (Mon): for $v'<v$, $X(\ell v)>\theta(v)-\kappa>\theta(v')-\kappa>X(\ell v')-2\kappa$.\qedhere
\end{enumerate}
\end{proof}

\formal{\lean{HPC.Hagopian.shadow_pass}.}

The shadow lemma holds only in the middle band: near an edge the line may turn back by up to $2\kappa$.

\begin{corollary}[Two scales]\label{cor:scales}
Put $\eta=\frac12$.
\begin{enumerate}[label=(\alph*)]
\item \emph{(Flow box.)} There is $\eta_2\in(0,\eta)$ such that every $u$ with $\ell(u)\in F([\frac9{32},\frac{23}{32}]\times[-\eta_2,\eta_2])$ lies in a pass of
$Q_\eta$ with (Mon); by (P1) every pass of $Q_\eta$ through such a point has (Mon).
\item A pass of $Q_\eta$ is \emph{thin} if it lies in $Q_{\eta_2}$. Every thin pass has (Mon) (it meets $X=\frac12$).
\item \emph{(Scale two.)} There is $\delta_2>0$ such that if $d(\ell(u),\ell(\frac12))<\delta_2$, then $u$ lies in the open interval of a thin pass.
\item \emph{(Late thin passes.)} For every $T$ there is a thin pass starting after $T$.
\end{enumerate}
\end{corollary}

\begin{proof}
(a) Take $\delta_1$ from Lemma~\ref{lem:shadow} for $\eta'=\eta$ and $\eta_2$ with $d(F(x,y),F(x,0))<\delta_1$ for $\abs y\le\eta_2$; apply the lemma at
$x=X(\ell u)$. (c) Lemma~\ref{lem:shadow} with $\eta'=\eta_2/2$, nesting and (P1). (d) Let $T'$ be the end of the pass containing $T$, if any,
and $T'=T$ otherwise. By recurrence (Theorem~\ref{thm:recurrent}) there is $u>T'$ with $d(\ell(u),\ell(\frac12))<\delta_2$; the thin pass $[a,b]\ni u$
has $a>T$, since otherwise it would contain $T$ and end at $T'<u$.
\end{proof}

\formal{\lean{HPC.Hagopian.flow_box}, \lean{HPC.Hagopian.FBProp}, \lean{HPC.Hagopian.PassThin},
\lean{HPC.Hagopian.scale_two}, \lean{HPC.Hagopian.exists_late_thin_pass}.}

\section{The trap}\label{sec:trap}

Let $P_m=\ell[a_m,b_m]$ be a thin pass above the arc with $a_m>\frac34$.

\begin{lemma}[The region under a thin pass]\label{lem:region}
$I_m=F^{-1}(P_m)$ is an arc from $r=(\frac14,r_2)$ to $s=(\frac34,s_2)$ in $[\frac14,\frac34]\times(0,\eta_2]$, meeting the vertical sides only in $r,s$.
There is a connected bounded open set $D_m\subseteq(\frac14,\frac34)\times(0,\eta_2)$ with
\[
\fr D_m=I_m\cup[\tfrac14,\tfrac34]\times\set0\cup\set{\tfrac14}\times[0,r_2]\cup\set{\tfrac34}\times[0,s_2].
\]
Put $R_m=F(D_m)$. For $c\in(\frac14,\frac34)$ let $w_m(c)$ be the least height of a point of $I_m$ over $c$ and $\tau_m(c)=\set c\times(0,w_m(c))$; let
$h_m>0$ be below all heights of $I_m$. Then $\tau_m(c)\subseteq D_m$, $(\frac14,\frac34)\times(0,h_m)\subseteq D_m$, and $\tau_m(c)$ is a cross-cut of $D_m$: no
connected subset of $D_m\setminus\tau_m(c)$ has closure meeting both open vertical sides.
\end{lemma}

\begin{proof}[Proof in outline]
$D_m$ is the bounded complementary domain of the simple closed curve formed by $I_m$ and the three sides, given by the Jordan curve
theorem (\vfive{Theorem~4.8}). For the cross-cut, the Jordan curve formed by the part of $I_m$ from $r$ to $(c,w_m(c))$, the left side, the
bottom from $\frac14$ to $c$ and $\tau_m(c)$ has the points of $D_m$ near the left side inside and the points near the right side outside.
\end{proof}

\formal{\lean{HPC.Hagopian.rect_domain_aux}, \lean{HPC.Hagopian.rect_domain_low}, \lean{HPC.Hagopian.seg_subset_rect_domain},
\lean{HPC.Hagopian.rect_domain_crosscut}, \lean{HPC.Hagopian.thinPass_arc}.}

\begin{lemma}[Passes entering the region]\label{lem:R1R2}
Let $P'=\ell[a',b']$ be a pass other than the base pass and $P_m$ whose interior meets $R_m$. Then (R1) its interior lies in $R_m$, its end
points lie on the open sides $F(\set{\frac14}\times(0,r_2))$ and $F(\set{\frac34}\times(0,s_2))$, and $P'$ is thin; and (R2) its interior meets $F(\tau_m(c))$
for every $c\in(\frac14,\frac34)$.
\end{lemma}

\begin{proof}
(R1) By (P4) $P'$ is above the arc; the preimage of its interior is connected, meets $D_m$, and misses $I_m$ (by (P1) and injectivity), the
bottom and the sides, so it lies in $D_m$. Its end points are limits of the interior, so they lie in $\fr D_m$, on the sides and not at the
corners. (R2) is the cross-cut property of Lemma~\ref{lem:region}.
\end{proof}

\formal{\lean{HPC.Hagopian.region_pass}, \lean{HPC.Hagopian.region_cross}.}

\begin{theorem}[The trapping lemma]\label{thm:trap}
Suppose that every thin pass above the arc starting after $\frac34$ starts at or after $a_m$. Then no pass starting after $b_m$ has its
interior in $R_m$.
\end{theorem}

\begin{proof}
\emph{Lemma A: no $v\in(\frac34,a_m)$ has $\ell(v)\in R_m$ with $X(\ell v)\in[\frac9{32},\frac{23}{32}]$.} Otherwise $R_m\subseteq Q_{\eta_2}$ and the flow box put $v$
in a pass $P'\ne$ base, $P_m$; by (R1) $P'$ is thin with interior in $R_m$, and by (P1) its interval lies in $(\frac34,a_m)$, against the
hypothesis.

\emph{Notation.} $w_L=w_m(\frac38)$, $y_m=w_m(\frac12)$, $w_R=w_m(\frac58)$; $T_L=F(\tau_m(\frac38))$, $\sigma=F(\set{\frac12}\times[0,y_m])$, $T_R=F(\tau_m(\frac58))$;
$x_m=F(\frac12,y_m)=\ell(t_m)$ and $z_R=F(\frac58,w_R)=\ell(t_R)$, with $t_m,t_R\in(a_m,b_m)$.

\emph{Step 1: $t_m<t_R$}, by (Mon) for $P_m$ and the order corollary ($\frac58-\frac12\ge2\kappa$).

\emph{Step 2: $J=\sigma\cup\ell[\frac12,t_m]$ is a simple closed curve.} $\sigma\cap\ell[\frac12,t_m]=\set{\ell(\frac12),x_m}$: the base piece $\ell[\frac12,\frac34]$
meets $\sigma$ only at $\ell(\frac12)$; the open segment $\sigma^\circ$ lies in $R_m$ and Lemma A excludes points of $\ell((\frac34,a_m))$ on it; points of $P_m$
on $X=\frac12$ have height at least $y_m$. By the Jordan curve theorem, $\R^2\setminus J=W\cup W'$ with $W,W'$ open, connected, disjoint, and
$\fr W=\fr W'=J$. Note \textbf{(J1)}: for $v>t_m$, $\ell(v)\in J$ iff $\ell(v)\in\sigma^\circ$.

\emph{Step 3.} Put $h=h_m$, $B_L=F([\frac38,\frac12)\times(0,h))$, $B_R=F((\frac12,\frac58]\times(0,h))$. The four connected sets $B_L,T_L,B_R,T_R$ lie in $R_m$
and miss $J$ (they miss $\sigma$ as $X\ne\frac12$, the base arc as $Y>0$, $P_m$ as they lie below its lowest points, and $\ell((\frac34,a_m))$ by
Lemma A). $B_L$ meets $T_L$ and $B_R$ meets $T_R$; let $W_L$ and $W_R$ be the components of $\R^2\setminus J$ containing $B_L\cup T_L$ and $B_R\cup T_R$.

\emph{Step 4: $W_L\ne W_R$.} Let $p=F(\frac12,h/2)\in\sigma^\circ$ and $S$ a small open square about $(\frac12,h/2)$ inside $[\frac38,\frac58]\times(0,h)$. $F(S)$ is
open and $F(S)\setminus J\subseteq B_L\cup B_R$. As $p$ is in the frontier of both components of $\R^2\setminus J$, each of them meets $B_L\cup B_R$; if both
$B_L$ and $B_R$ lay in one component, the other would miss $B_L\cup B_R$.

\emph{Step 5: $z_R\in W_R$.} $z_R\notin J$ by (J1), and $z_R\in\cl{T_R}$.

\emph{Step 6.} Suppose $P=\ell[a,b]$ has $a>b_m$ and interior in $R_m$. By (R2) at $c=\frac38$ there is $v_L\in(a,b)$ with $\ell(v_L)\in T_L\subseteq W_L$.
The path $\ell[t_R,v_L]$ joins $W_R$ to $W_L$, so it meets $J$; let $s^*$ be the largest $v\in[t_R,v_L]$ with $\ell(v)\in J$. By (J1),
$\ell(s^*)\in\sigma^\circ$, and $\ell((s^*,v_L])\subseteq W_L$. Also $s^*>b_m$, because $\ell((t_m,b_m])$ misses $\sigma$. By the flow box at $X=\frac12$, $s^*$
lies in the open interval of a pass $P^*=\ell[a^*,b^*]$ with (Mon); by (R1) its interior lies in $R_m$ and by (P1) $a^*>b_m$. If $s^*\in(a,v_L)$ then
$P^*=P$, and (Mon) for $P$ would put its point on $X=\frac38$ after its point on $X=\frac12$, which the order corollary forbids; so $b^*<a$. By
(R2) at $c=\frac58$ there is $v_R\in(a^*,b^*)$ with $\ell(v_R)\in T_R\subseteq W_R$, and (Mon) for $P^*$ gives $v_R>s^*$. So $v_R\in(s^*,v_L)$ and
$\ell(v_R)\in W_L$, contradicting Step~4.
\end{proof}

\formal{\lean{HPC.Hagopian.region_lemmaA}, \lean{HPC.Hagopian.two_sides}, \lean{HPC.Hagopian.trap_jordan},
\lean{HPC.Hagopian.trap_pass}; the Jordan curve theorem is \lean{HPC.Jones.jordan_curve_theorem}.}

\section{The contradiction}

\begin{theorem}[Bing 1960; Jones 1975]\label{thm:noarc}
No homogeneous indecomposable plane continuum contains an arc.
\end{theorem}

\begin{proof}
Suppose $M$ does, and let $\ell$, $F$ be as above.

\emph{C1.} By Corollary~\ref{cor:scales}(d) and (P2) let $P_1=\ell[a_1,b_1]$ be the first thin pass with $a_1>\frac34$; flipping the collar if
necessary, it is above the arc. The hypothesis of Theorem~\ref{thm:trap} holds for $R_1$: a violating pass would be thin (R1) and start in
$(\frac34,a_1)$.

\emph{C2.} Take $h_1=h_{m_1}$ and $\rho_1$ with $B(\ell(\frac12),\rho_1)\subseteq F((\frac38,\frac58)\times(-h_1,h_1))$. Recurrence gives $u_1>b_1$ with
$d(\ell(u_1),\ell(\frac12))<\min(\delta_2,\rho_1)$.

\emph{C3.} Write $\ell(u_1)=F(x,y)$. If $y=0$, then $u_1=x<b_1$ by injectivity. If $y>0$, $\ell(u_1)\in R_1$ lies on a thin pass (scale two) whose
interior is in $R_1$ (R1) and which starts after $b_1$ (P1), against Theorem~\ref{thm:trap}. So $y<0$, and some thin pass below the arc starts after
$b_1$; let $P_2=\ell[a_2,b_2]$ be the first.

\emph{C4.} The trapping lemma applies to $P_2$ (with the flipped collar): a thin pass below the arc starting in $(\frac34,a_2)$ would, by (P1),
start in $(\frac34,a_1)$, against the minimality of $P_1$, or in $(b_1,a_2)$, against the minimality of $P_2$.

\emph{C5.} Take $h=\min(h_1,h_2)$ and $\rho$ with $B(\ell(\frac12),\rho)\subseteq F((\frac38,\frac58)\times(-h,h))$; recurrence gives $u>b_2$ with
$d(\ell(u),\ell(\frac12))<\min(\delta_2,\rho)$. $Y(\ell u)=0$ is excluded by injectivity; $Y>0$ puts $\ell(u)$ in $R_1$ on a thin pass after $b_1$ (C1);
$Y<0$ puts it in $R_2$ on a thin pass after $b_2$ (C4). This is a contradiction.
\end{proof}

\formal{\lean{HPC.Hagopian.jones75_final_aux}, \lean{HPC.Hagopian.jones75_contra},
\lean{HPC.Hagopian.NoArcIndecHomPlane}, \lean{HPC.Hagopian.jones75_noArc}.}

\section{Hereditary indecomposability}

\begin{lemma}\label{lem:monotone-image}
A monotone image of an indecomposable continuum is indecomposable: if $q\colon X\to Y$ is a continuous surjection of a continuum onto a metric space
with connected fibres and $X$ is indecomposable, so is $Y$.
\end{lemma}

\begin{proof}
$q$ is closed with connected fibres, so preimages of subcontinua are subcontinua (as in Lemma~\ref{lem:adm-basic}(a)). If $Y=H\cup K$ with $H,K$
proper subcontinua, then $X=q^{-1}(H)\cup q^{-1}(K)$ with both proper.
\end{proof}

\formal{\lean{HPC.Hagopian.indecomposable_of_monotone_image}.}

\begin{theorem}[Hagopian 1976, Theorem~1]\label{thm:HI}
Every homogeneous indecomposable plane continuum is hereditarily indecomposable.
\end{theorem}

\begin{proof}
This is Hagopian's 1984 proof (his Corollary~3), following a plan of Jones. Suppose $M$ is not hereditarily indecomposable, and let $q\colon M\to Y$ be
as in Theorem~\ref{thm:h84}. Its fibres are proper subcontinua of $M$, so they do not separate the plane (Theorem~\ref{thm:nosep}). By Moore's
theorem on monotone decompositions (\vfive{Corollary~9.16}), $Y$ is homeomorphic to a compact plane set $Q$. $Q$ is connected (an image of $M$),
homogeneous (as $Y$ is), indecomposable (Lemma~\ref{lem:monotone-image}) and contains an arc. This contradicts Theorem~\ref{thm:noarc}.
\end{proof}

\formal{\lean{HPC.Hagopian.exists_plane_quotient_of_monotone}, \lean{HPC.Hagopian.hagopian76_theorem1_of_jones},
\lean{HPC.Hagopian.hagopian76_theorem1_jones}, \lean{HPC.Hagopian.Hagopian76Thm1}, \lean{HPC.Hagopian.hagopian76_theorem1}.}

\begin{corollary}\label{cor:bing60}
\leavevmode
\begin{enumerate}[label=(\alph*)]
\item \emph{(Bing 1960, for plane continua.)} A homogeneous plane continuum that contains an arc is a simple closed curve.
\item \emph{(Hagopian 1976, Theorem~2.)} A homogeneous plane continuum that does not separate the plane is hereditarily indecomposable.
\end{enumerate}
\end{corollary}

These consequences are recorded in the formal development; they are not used in this volume.

\formal{\lean{HPC.Hagopian.bing60_plane_uncond}, \lean{HPC.Hagopian.hagopian76_theorem2_uncond}.}

\section*{Exercises}

\begin{exercise}
Show that in a compact metric space the forward limit set of a continuous map $\ell\colon\R\to Y$ is nonempty and compact, and that it does not
change if $\ell$ is replaced by $s\mapsto\ell(s+c)$.
\end{exercise}

\begin{exercise}
Show that the dyadic solenoid is homogeneous and that its arc components are lines. Which part of the proof of Theorem~\ref{thm:noarc} fails for
it?
\end{exercise}

%======================================================================
\chapter{Span zero}\label{ch:span}
%======================================================================

The \emph{span} of a continuum, introduced by Lelek, measures how far apart two points can stay while each of them runs
over the whole continuum (or one runs over the part traced by the other). Snake-like continua have span zero, and the
main theorem of Hoehn and Oversteegen (Chapters~\ref{ch:snake}--\ref{ch:ladders}) is a converse for hereditarily
indecomposable continua. This chapter and the next prove:

\begin{theorem}[Theorem~C of Hoehn--Oversteegen]\label{thm:theoremC}
Every homogeneous indecomposable plane continuum has span zero.
\end{theorem}

\formal{\lean{HPC.Span.TheoremC}, \lean{HPC.Span.theoremC}.}

The proof follows Oversteegen and Tymchatyn, \emph{Plane strips and the span of continua (I)} (1982): Theorem~3.2 (the strip
theorem, Chapter~\ref{ch:strips}) $\Rightarrow$ Corollary~3.5 $\Rightarrow$ Theorem~3.6 $\Rightarrow$ Theorem~C. Several
steps of the printed proofs are unargued or, as printed, impossible; the formal proof supplies them, and proves the strip
theorem in a slightly weaker form that suffices.

\section{Span}

Let $X$ be a metric space with projections $\pr1,\pr2\colon X\times X\to X$, and $\Diag$ the diagonal. Only positivity of span matters,
so we use topological forms.

\begin{definition}\label{def:span}
A set $Z\subseteq X\times X$ is \emph{off the diagonal} if $Z\cap\Diag=\emptyset$. A set $A\subseteq X$ has \emph{positive surjective span} if there is a
continuum $Z\subseteq A\times A$ off the diagonal with $\pr1(Z)=\pr2(Z)=A$, and \emph{positive surjective semi-span} if there is one with $\pr1(Z)=A$.
$X$ has \emph{span zero} (\emph{semi-span zero}) if no subcontinuum of $X$ has positive surjective span (semi-span). Following
Hoehn--Oversteegen, $X$ has \emph{span zero in their sense} if every subcontinuum $Z\subseteq X\times X$ with $\pr1(Z)\subseteq\pr2(Z)$ meets the diagonal.
\end{definition}

\formal{\lean{HPC.Span.OffDiag}, \lean{HPC.Span.SurjSpanPos}, \lean{HPC.Span.SurjSemiSpanPos},
\lean{HPC.Span.SpanZero}, \lean{HPC.Span.SemiSpanZero}, \lean{HPC.Span.HOSpanZero}.}

\begin{proposition}\label{prop:span-basic}
\leavevmode
\begin{enumerate}[label=(\alph*)]
\item For compact $X$, Lelek's numerical definitions agree with these: $\sigma_0(X)>0$ (respectively $\sigma(X)>0$) iff some subcontinuum has
positive surjective semi-span (span).
\item Positive surjective span (semi-span) of $A$ passes to $f(A)$ for any continuous $f$ injective on $A$.
\item Semi-span zero is the same as span zero in the sense of Hoehn--Oversteegen; and that is equivalent to: for all maps $f,g\colon C\to X$
of a continuum with $f(C)\subseteq g(C)$ there is $c$ with $f(c)=g(c)$.
\end{enumerate}
\end{proposition}

\begin{proof}
(a) Given a connected $C\subseteq A\times A$ with $d\ge\alpha>0$ on $C$ and $\pr1(C)=A$, the closure $\cl C$ is a continuum off the diagonal with
$\pr1(\cl C)=\cl A$; conversely, on a continuum off the diagonal $d$ has a positive minimum. (b) Use $(f\times f)(Z)$. (c) Use the swap
$(a,b)\mapsto(b,a)$: if $Z$ is off the diagonal with $\pr1(Z)\subseteq\pr2(Z)=A$, its swap lies in $A\times A$ and has first projection $A$; and
conversely. For maps, take $Z=\set{(f(c),g(c))}$, and conversely $C=Z$ with the two projections.
\end{proof}

\formal{\lean{HPC.Span.otSemiSpanPos_iff}, \lean{HPC.Span.otSpanPos_iff}, \lean{HPC.Span.SurjSpanPos.image},
\lean{HPC.Span.SurjSemiSpanPos.image}, \lean{HPC.Span.semiSpanZero_iff_hoSpanZero}, \lean{HPC.Span.hoSpanZero_iff_maps}.}

\section{Confluence and Davis's theorem}

\begin{theorem}[Cook 1967, Theorem~4]\label{thm:cook}
Let $f$ be a continuous map of a compact metric continuum $X$ onto a hereditarily indecomposable continuum $S$. Then $f$ is
\emph{confluent}: for every subcontinuum $Q\subseteq S$ and every component $C$ of $f^{-1}(Q)$, $f(C)=Q$.
\end{theorem}

\begin{proof}
If $C=X$ then $f(C)=S\supseteq Q$. Otherwise let $q\in Q$ and suppose $q\notin f(C)$. Then $f^{-1}(q)$ is compact and misses $C$, so a closed
$\delta$-neighbourhood $N$ of $C$ misses $f^{-1}(q)$; we may take $N\ne X$. Let $K$ be the component of a point $x\in C$ in $N$. By boundary
bumping (\vthree{Theorem~1.10}) $K$ meets the frontier of $N$, so $K\not\subseteq C$. Then $f(K)$ meets $Q$ (at $f(x)$) and is not inside $Q$
(otherwise $K$ would be a connected subset of $f^{-1}(Q)$ larger than $C$). Two subcontinua of the hereditarily indecomposable $S$ that meet are
nested, so $Q\subseteq f(K)$, and $q=f(z)$ with $z\in K\subseteq N$; but $N$ misses $f^{-1}(q)$.
\end{proof}

\formal{\lean{HPC.Span.nested_of_hi}, \lean{HPC.Span.cook67_theorem4}.}

This uses one closed neighbourhood of $C$ instead of Cook's decreasing sequence $C_n$, whose existence Cook does not argue.

\begin{lemma}[Brouwer reduction]\label{lem:brouwer}
Let $Z$ be a continuum and $P$ a property of subcontinua of $Z$ that holds for $Z$ and passes to intersections of nonempty chains. Then $Z$
contains a subcontinuum $Z_0$ with $P$ that is minimal: no proper subcontinuum of $Z_0$ has $P$. Moreover, the intersection of a nonempty chain
of continua is a continuum, and continuous images commute with such intersections.
\end{lemma}

\begin{proof}
The intersection of a nonempty chain of continua is compact, nonempty by the finite intersection property, and connected as a directed
intersection of continua. Zorn's lemma gives $Z_0$. For a continuous $\varphi$ and a chain $(W_j)$ of compact sets, $\varphi(\bigcap W_j)=\bigcap\varphi(W_j)$:
if $y\in\bigcap\varphi(W_j)$, the sets $W_j\cap\varphi^{-1}(y)$ form a chain of nonempty compacta.
\end{proof}

\formal{\lean{HPC.Span.isContinuum_sInter_chain}, \lean{HPC.Span.brouwer_reduction}.}

\begin{theorem}[Davis's theorem for hereditarily indecomposable continua]\label{thm:davis}
If $A$ is hereditarily indecomposable and has positive surjective semi-span, some subcontinuum $B\subseteq A$ has positive surjective span.
\end{theorem}

\begin{proof}
Let $Z\subseteq A\times A$ be a continuum off the diagonal with $\pr1(Z)=A$. The property $P(W)$: ``$\pr2(W)\subseteq\pr1(W)$'' holds for $Z$ and passes to
intersections of chains (Lemma~\ref{lem:brouwer}); let $Z_0\subseteq Z$ be minimal with it, and $B=\pr1(Z_0)$, a hereditarily indecomposable continuum.
Suppose $Q=\pr2(Z_0)\ne B$. Take $w\in Z_0$ with $\pr1(w)\in Q$ and let $Z'$ be the component of $w$ in $Z_0\cap\pr1^{-1}(Q)$. By Cook's theorem (applied to
$\pr1\colon Z_0\to B$), $\pr1(Z')=Q$; so $\pr2(Z')\subseteq Q=\pr1(Z')$ and $Z'\ne Z_0$, against minimality. Hence $\pr1(Z_0)=\pr2(Z_0)=B$.
\end{proof}

\formal{\lean{HPC.Span.exists_surjSpanPos_of_surjSemiSpanPos}, \lean{HPC.Span.semiSpanZero_of_spanZero}.}

\section{Theorem 3.6 from Corollary 3.5}

\begin{definition}
\emph{Corollary~3.5} is the statement: if $M$ is an indecomposable homogeneous plane continuum, every proper subcontinuum of $M$ has zero surjective
semi-span (that is, does not have positive surjective semi-span).
\end{definition}

\formal{\lean{HPC.Span.OTCor35}.}

\begin{lemma}\label{lem:dense-proper}
A nondegenerate indecomposable continuum has, for every $\delta>0$, a $\delta$-dense proper subcontinuum.
\end{lemma}

\begin{proof}
The proper subcontinua through a point $x$ are directed by union (Theorem~\ref{thm:composants}(b)) and their union, the composant, is dense
(Theorem~\ref{thm:composants}(a)). Finitely many of them meet each ball of a finite $\delta/2$-net, and their union is one proper subcontinuum.
\end{proof}

\formal{\lean{HPC.Span.union_proper_of_indec}, \lean{HPC.Span.exists_dense_proper_subcontinuum}.}

\begin{theorem}[Oversteegen--Tymchatyn, Theorem~3.6, indecomposable case]\label{thm:ot36}
Let $X$ be a homogeneous, indecomposable, hereditarily indecomposable continuum none of whose proper subcontinua has positive surjective
semi-span. Then $X$ does not have positive surjective semi-span.
\end{theorem}

\begin{proof}
Suppose $Z\subseteq X\times X$ is a continuum off the diagonal with $\pr1(Z)=X$; write $f=\pr1$, $g=\pr2$ on $Z$, and let $\alpha>0$ be the minimum of $d(f,g)$ on $Z$.

\emph{Reduction.} We may assume that $g(Z)\ne X$, or that no proper subcontinuum of $Z$ is mapped onto $X$ by $g$: if $g(Z)=X$, take by
Lemma~\ref{lem:brouwer} a minimal $Z_0\subseteq Z$ with $g(Z_0)=X$; if $f(Z_0)=X$ replace $Z$ by $Z_0$, and otherwise by the swap of $Z_0$ (whose
first projection is onto and second is not).

\emph{Effros.} Let $\delta$ be the Effros bound (\vthree{Theorem~7.15}) for $\alpha$, and $K$ a $\delta$-dense proper subcontinuum (Lemma~\ref{lem:dense-proper};
$X$ is nondegenerate since $Z$ exists).

\emph{Cook.} By Theorem~\ref{thm:cook}, a component $C$ of $f^{-1}(K)$ has $f(C)=K$; $C\ne Z$ since $K\ne X$.

\emph{The copy.} Take $y\in g(C)$, $x\in K$ with $d(x,y)<\delta$, and an $\alpha$-homeomorphism $h$ with $h(x)=y$. The subcontinua $h(K)$ and $g(C)$ meet at
$y$, so they are nested; let $Z^*$ be the larger. It is proper: $h(K)\ne X$, and $g(C)\ne X$ because either $g(Z)\ne X$, or $g$ is irreducible onto $X$
and $C\subsetneq Z$.

\emph{The contradiction.} $W=\set{(h(f(c)),g(c)):c\in C}$ is a continuum with projections $h(K)$ and $g(C)$; it is off the diagonal, since
$h(f(c))=g(c)$ would give $d(f(c),g(c))=d(f(c),h(f(c)))<\alpha$. If $g(C)\subseteq h(K)=Z^*$, $W$ witnesses positive surjective semi-span of $Z^*$; otherwise its
swap does. This contradicts the hypothesis.
\end{proof}

\formal{\lean{HPC.Span.exists_reduced_witness}, \lean{HPC.Span.not_surjSemiSpanPos_univ}.}

\begin{corollary}\label{cor:C-from-35}
Corollary~3.5 implies Theorem~C.
\end{corollary}

\begin{proof}
Let $M$ be a homogeneous indecomposable plane continuum. It is hereditarily indecomposable (Theorem~\ref{thm:HI}). By Corollary~3.5 and
Theorem~\ref{thm:ot36} it has semi-span zero, which is span zero in the sense of Hoehn--Oversteegen
(Proposition~\ref{prop:span-basic}(c)).
\end{proof}

\formal{\lean{HPC.Span.ot82_theorem36_of_cor35}, \lean{HPC.Span.theoremC_of_cor35}.}

\section{Corollary 3.5 from the strip theorem}

The strip theorem is proved in the next chapter in the following form.

\begin{definition}[The strip theorem, strip form]\label{def:otthm32}
Let $Y$ be a continuum, $X\subseteq Y$ a subcontinuum, $x\in X$ and $y\in Y$. Let $(K_i)_{i\in I}$ be pairwise disjoint continua in the strip
$[0,1]\times\R$, each meeting both lines $x=0$ and $x=1$, and $h_i\colon Y\to K_i$ continuous maps with $h_i(x)$ on the line $x=0$ and
$\pr1(h_i(z))<\pr1(h_i(y))$ for all $z\in X$ (here $\pr1$ is the first coordinate). Then only countably many of the sets $h_i(X)$ have positive surjective
span.
\end{definition}

\formal{\lean{HPC.Span.OTThm32}.}

Oversteegen and Tymchatyn's Theorem~3.2 does not have the continua $K_i$; it only asks that the images $h_i(Y)$ be pairwise disjoint. Its printed
proof relies on two ``we may suppose'' steps that are not available together, and on three images one of which separates the other two, which is
not argued. The strip form is what Corollary~3.5 uses.

\begin{lemma}[Re-embedding]\label{lem:reembed}
\leavevmode
\begin{enumerate}[label=(E\arabic*)]
\item For a nonempty compact $K\subseteq\R^2$ there is a homeomorphism $\Phi$ of the plane with $\Phi(K)\subseteq\set{x\ge0}$ meeting the line $x=0$ in exactly one
point: if $(u_0,v_0)$ has least first coordinate in $K$, take $\Phi(u,v)=(u-u_0+(v-v_0)^2,v)$.
\item \emph{(Finger move.)} If $U$ is a connected open set and $p,q\in U$, there is a homeomorphism of the plane fixing every point off $U$ and
taking $p$ to $q$.
\item If $A\subseteq\set{x\ge0}$ is a continuum that does not separate the plane and meets $x=0$ in at most one point, then $\set{x>0}\setminus A$ is connected.
\end{enumerate}
\end{lemma}

\begin{proof}
(E2) If $B(c,r)\subseteq U$ and $\nrm{z-c}<r/2$, the map $T(w)=w+\beta(w)(z-c)$ with $\beta(w)=\max(0,1-\nrm{w-c}/r)$ is a homeomorphism (the perturbation is
Lipschitz with constant less than $\frac12$, so $T$ is injective and, by the contraction principle, onto), $T(c)=z$, and $T$ is the identity off
$B(c,r)$. The set of points reachable from $p$ by compositions of such maps is open and closed in $U$.

(E3) Let $F$ be the part of $\set{x>0}$ above, below, or to the right of a large box containing $A$; it is connected and misses $A$. Every component
$C$ of $\set{x>0}\setminus A$ meets $F$: otherwise $C$ is bounded with frontier in $A\cup J$, $J$ the segment of the line $x=0$ next to $A$, so $A\cup J$
separates the plane; but $A\cap J$ has at most one point and neither $A$ nor $J$ separates, against Janiszewski's theorem (\vfour{Theorem~2.7}).
\end{proof}

\formal{\lean{HPC.Span.exists_homeomorph_unique_leftmost}, \lean{HPC.Span.exists_homeomorph_ball_step},
\lean{HPC.Span.exists_homeomorph_push}, \lean{HPC.Span.isConnected_rightHalf_diff}.}

\begin{lemma}[Crossing continua]\label{lem:crossing}
\leavevmode
\begin{enumerate}[label=(\alph*)]
\item \emph{(The cut-wire argument.)} If $P$ is a continuum with points on both sides of the strip $S=[a,b]\times\R$, some component of $P\cap S$ meets
both lines $x=a$ and $x=b$.
\item If $M$ is a nondegenerate indecomposable plane continuum with points left of $x=a$ and right of $x=b$, there is an uncountable family of
pairwise disjoint subcontinua of $M\cap S$ each meeting both lines.
\end{enumerate}
\end{lemma}

\begin{proof}
(a) Otherwise, by \vthree{Theorem~1.9} applied in the compact set $P\cap S$, there is a relatively clopen $F_1$ containing $P\cap\set{x=a}$ and missing
$P\cap\set{x=b}$. With $F_2=(P\cap S)\setminus F_1$, the sets $(P\cap\set{x\le a})\cup F_1$ and $(P\cap\set{x\ge b})\cup F_2$ are disjoint, closed, nonempty and
cover $P$.

(b) Each composant $\kappa$ is dense, so it contains a point $p$ left of $x=a$, and $\cmp Mp=\kappa$ contains a point $q$ right of $x=b$; a proper subcontinuum
$P\ni p,q$ lies in $\kappa$, and (a) gives $K_\kappa\subseteq P\cap S$. Distinct composants are disjoint and uncountably many
(Theorem~\ref{thm:composants}).
\end{proof}

\formal{\lean{HPC.Span.exists_crossing_subcontinuum}, \lean{HPC.Span.exists_uncountable_crossing}.}

\begin{theorem}[Oversteegen--Tymchatyn, Corollary~3.5]\label{thm:ot35}
The strip theorem implies Corollary~3.5.
\end{theorem}

\begin{proof}
Let $M$ be an indecomposable homogeneous plane continuum and $A_0\ne M$ a subcontinuum with positive surjective semi-span. $M$ is hereditarily
indecomposable (Theorem~\ref{thm:HI}).

\emph{1. Davis.} By Theorem~\ref{thm:davis} some $A_1\subseteq A_0$ has positive surjective span.

\emph{2. Re-embed.} By (E1) we may assume $M\subseteq\set{x\ge0}$ with a unique point $x_0$ on $x=0$ (positivity of span and all hypotheses are carried
by $\Phi$). By homogeneity move $A_1$ to a continuum $A\ni x_0$; $A\ne M$ and $A$ has positive surjective span.

\emph{3. Finger move.} $A$ does not separate the plane (Theorem~\ref{thm:nosep}), so by (E3) $U=\set{x>0}\setminus A$ is open and connected. Take
$p\in M\setminus A$ and $q$ to the right of $A$, and by (E2) a homeomorphism $T$ fixing everything off $U$ with $T(p)=q$. Replacing $M$ by $T(M)$, we may
assume in addition that some point of $M$ lies to the right of $A$. (These printed ``without loss of generality'' steps are not a matter of
scaling; this is how the formal proof achieves them.)

\emph{4. Constants.} Let $s<m$ be the largest first coordinates on $A$ and on $M$, $\eta=(m-s)/8$ and $r=s+4\eta$. Let $\delta$ be the Effros bound for $\eta$,
and choose $\varepsilon\in(0,r)$ with $M\cap\set{x\le\varepsilon}\subseteq B(x_0,\delta)$.

\emph{5. The family.} Lemma~\ref{lem:crossing}(b) gives an uncountable family $\mathcal K$ of disjoint subcontinua crossing $[\varepsilon,r]\times\R$.

\emph{6. Copies of $A$.} For $K\in\mathcal K$ pick $x_K\in K$ on $x=\varepsilon$ and an $\eta$-homeomorphism $h_K$ with $h_K(x_0)=x_K$. $h_K(A)$ and $K$ meet, so
they are nested; $K\subseteq h_K(A)$ is impossible, since $K$ reaches $x=r$ while $h_K(A)\subseteq\set{x\le s+\eta}$. So $h_K(A)\subseteq K$.

\emph{7. The continuum $B$ and the point $y$.} The continua $B_K=h_K^{-1}(K)$ all contain $A$, so they are pairwise nested, and $B=\bigcap_KB_K$ is a
continuum containing $A$. The sets $B_K\cap\set{x\ge r-\eta}$ form a chain of nonempty compacta, so their intersection contains a point $y\in B$. For $z\in
A$: $\pr1h_K(z)<s+\eta<r-2\eta<\pr1h_K(y)$.

\emph{8.} Rescale the strip $[\varepsilon,r]\times\R$ to $[0,1]\times\R$ and apply the strip theorem with $Y=B$, $X=A$, $x=x_0$, $y$, the continua $K\in\mathcal K$ and the
maps $h_K|_B$. Every $h_K(A)$ has positive surjective span (Proposition~\ref{prop:span-basic}(b)), so $\mathcal K$ would be countable.
\end{proof}

\formal{\lean{HPC.Span.cor35_core}, \lean{HPC.Span.ot82_cor35_of_thm32}, \lean{HPC.Span.theoremC_of_thm32}.}

\begin{remark}
In the printed proof of Corollary~3.5, the normalization $\pr1(A)=[0,\frac12]$, the existence of $\mathcal K$, and the continuum $B$ (``it is easy to see'') are
unargued, the supposition $K\cap(\set\varepsilon\times\R)=\set{x_K}$ is unjustified (and unnecessary), and the inclusion ``$[0,\frac58]\subseteq\pr1(h_K(B))$'' is
impossible as printed; step~7 gives what is needed, one point $y$ for all $K$.
\end{remark}

\section*{Exercises}

\begin{exercise}
Show that an arc has span zero, and that a simple closed curve has positive surjective span.
\end{exercise}

\begin{exercise}
Show that a continuous image of a continuum with span zero need not have span zero, but that span zero is a topological invariant.
\end{exercise}

%======================================================================
\chapter{The strip theorem}\label{ch:strips}
%======================================================================

This chapter proves the strip theorem of Definition~\ref{def:otthm32}, completing the proof of Theorem~C. The idea of
Oversteegen and Tymchatyn is that a continuum lying in a thin strip between two paths cannot have large surjective span: two
points running over it while staying far apart would have to pass each other inside the strip. The formal proof replaces their
Lemmas~2.3, 2.4, Remark~2.5 and the zero-set continua of their Theorem~2.7 by a common-zero lemma on a rectangle and a direct
parity argument, and uses no Jordan curve theorem: the crossing parity of closed polygons from Volume~4 (which need not be
simple) is enough.

Throughout, the parity $\parity(P,z)\in\set{0,1}$ of a closed polygon $P$ at a point $z$ off $P$ is the number of edges crossed by
the eastward ray from $z$, counted by the half-open rule of \vfour{Definition~3.4}, modulo~$2$. It is locally constant off $P$
(\vfour{Theorem~3.7}) and $0$ far away. $L_0$ denotes the line $x=0$.

\formal{\lean{HPC.Plane.parity}, \lean{HPC.Plane.polygonOf}.}

\section{Parity tools}

\begin{lemma}\label{lem:parity-tools}
Let $P$ be a closed polygon with vertices $v_0,\dots,v_n$.
\begin{enumerate}[label=(\alph*)]
\item If $P$ lies in a compact convex set $C$ and $z\notin C$, then $\parity(P,z)=0$.
\item \emph{(Perturbation.)} If $w_0,\dots,w_n$ satisfy $\nrm{w_k-v_k}\le\rho$ and $\rho<\dist(z,P)$, then the polygon $P'$ with vertices $w_k$ has
$\parity(P',z)=\parity(P,z)$.
\item \emph{(Flip across a vertical edge.)} If edge $i$ is a nondegenerate vertical segment through $b=(c,b_2)$ with $b_2$ strictly between the heights of its
ends, and no other edge comes within $r$ of $b$, then for $0<\xi<r$ the parities at $(c+\xi,b_2)$ and $(c-\xi,b_2)$ differ.
\end{enumerate}
\end{lemma}

\begin{proof}
(a) Separate $z$ from $C$ by an open half-plane $H\ni z$; $H$ is connected, misses $P$, and contains far points, where the parity is $0$.

(b) Move one vertex at a time. Replacing $v_k$ by $w_k$ changes the parity by the parities of the two closed triangles $[v_{k-1},v_k,w_k]$ and
$[v_k,v_{k+1},w_k]$ (the edge $[v_k,w_k]$ is counted twice); each lies in the closed $\rho$-neighbourhood of an edge of $P$, a compact convex set
not containing $z$, so by (a) both parities vanish at $z$.

(c) The two eastward rays are at the same height; an edge is crossed by exactly one of them iff it straddles the height $b_2$ at a point
with first coordinate in $(c-\xi,c+\xi]$. Edge $i$ qualifies, and no other edge does.
\end{proof}

\formal{\lean{HPC.Span.parity_eq_zero_of_not_mem_convex}, \lean{HPC.Span.parity_eq_of_perturb},
\lean{HPC.Span.parity_flip_vertical_edge}.}

\section{A common zero on a rectangle}

\begin{lemma}\label{lem:common-zero}
Let $R=[0,1]\times[p,q]$ and $\varphi,\psi\colon R\to\R$ continuous with: $\varphi\ge0$ on the left side $\set0\times[p,q]$, $\varphi\le0$ on the right side, $\varphi\le0$ on
the bottom row $[0,1]\times\set p$, and $\psi\ge0$ on the left side, $\psi\le0$ on the right side, $\psi\le0$ on the top row $[0,1]\times\set q$. Then $\varphi$ and
$\psi$ have a common zero.
\end{lemma}

\begin{proof}
Suppose not. Then $\Phi=\varphi+i\psi$ is continuous and nowhere zero on the convex set $R$, so it has a continuous logarithm, and $\theta=\Im\log\Phi$ is a
continuous argument. At $(0,p)$, $\varphi=0$ ($\ge0$ from the left side, $\le0$ from the bottom), so $\psi>0$ and $\theta\equiv\pi/2$; at $(0,q)$, $\psi=0$, so
$\varphi>0$ and $\theta\equiv0$ (mod $2\pi$). On each side $\Phi$ lies in a closed quarter- or half-plane, so by connectedness $\theta$ stays in one interval of
length less than $2\pi$:
\begin{itemize}
\item left side (first quadrant): $\theta\in[2\pi k_1,2\pi k_1+\frac\pi2]$, so $\theta(0,p)=2\pi k_1+\frac\pi2$, $\theta(0,q)=2\pi k_1$;
\item top row ($\psi\le0$): $\theta\in[2\pi k_2-\pi,2\pi k_2]$, and at $(0,q)$ this forces $k_2=k_1$;
\item right side (third quadrant): $\theta\in[2\pi k_3-\pi,2\pi k_3-\frac\pi2]$, and at $(1,q)$ this forces $k_3=k_1$;
\item bottom row ($\varphi\le0$): $\theta\in[2\pi k_4+\frac\pi2,2\pi k_4+\frac{3\pi}2]$, and at $(1,p)$ this forces $k_4=k_1-1$.
\end{itemize}
Then $\theta(0,p)=2\pi k_1+\frac\pi2$ would lie in $[2\pi k_1-\frac{3\pi}2,2\pi k_1-\frac\pi2]$, which it does not.
\end{proof}

\formal{\lean{HPC.Span.common_zero_rect}, using \lean{HPC.Moore.exists_log_of_convex}.}

\section{The strip lemma}

\emph{Data.} Let $s\ge1$ and $a_0,\dots,a_s$, $b_0,\dots,b_s$ points of $\set{x\ge0}$ with $a_0,b_0\in L_0$. Let $f$ and $g$ be the piecewise linear paths on
$[0,1]$ through the $a_j$ and the $b_j$ at the times $j/s$. Put $B_t=[f(t),g(t)]$, so that $B_0=[a_0,b_0]\subseteq L_0$, and let $S_t$ be the closed polygon
$B_0\cup f[0,t]\cup B_t\cup g[0,t]$ (traversed: $B_0$, along $f$, across $B_t$, back along $g$), $U_t$ the set of points off $S_t$ of parity $1$, and $\varepsilon\ge\max_j\nrm{a_j-b_j}$,
so that $\nrm{f(t)-g(t)}\le\varepsilon$.

\formal{\lean{HPC.Span.plPath}, \lean{HPC.Span.plPath_piece}, \lean{HPC.Span.stripPoly}.}

\begin{lemma}\label{lem:strip-basic}
\leavevmode
\begin{enumerate}[label=(\alph*)]
\item $U_0=\emptyset$, and $U_t\subseteq\set{x>0}$ for all $t$.
\item \emph{(Oversteegen--Tymchatyn, Lemma~2.2.)} If $z_0\notin S_{t_0}$, then for $(t,z)$ near $(t_0,z_0)$, $z\notin S_t$ and $\parity(S_t,z)=\parity(S_{t_0},z_0)$.
\end{enumerate}
\end{lemma}

\begin{proof}
(a) $S_0$ is the segment $B_0$ traversed twice, so its edges cancel. On $\set{x<0}$ the parity vanishes by Lemma~\ref{lem:parity-tools}(a) (the convex hull of
the vertices lies in $\set{x\ge0}$), and a point of $L_0$ off $S_t$ has a disc off $S_t$ meeting $\set{x<0}$. (b) By uniform continuity the vertices of $S_t$ are
within $\rho<\dist(z_0,S_{t_0})/4$ of those of $S_{t_0}$ for $t$ near $t_0$; apply Lemma~\ref{lem:parity-tools}(b) and local constancy.
\end{proof}

\formal{\lean{HPC.Span.parity_strip_eventually}.}

\emph{The sign function.} Let $V=U_1\cup(B_0\setminus(f[0,1]\cup g[0,1]\cup B_1))$ and $h\colon[0,1]\to V$ continuous. Define on $[0,1]^2$
\[
\varphi_h(t,u)=\begin{cases}-\dist(h(u),B_t)&\text{if }h(u)\in U_t\text{ or }h(u)\in L_0,\\ +\dist(h(u),B_t)&\text{otherwise.}\end{cases}
\]

\formal{\lean{HPC.Span.StripInner}, \lean{HPC.Span.stripSign}.}

\begin{lemma}\label{lem:sign}
$\varphi_h$ is continuous; $\varphi_h(0,u)\ge0$, $\varphi_h(1,u)\le0$, and $\varphi_h(t,u_0)\le0$ for all $t$ if $h(u_0)\in L_0$; and $\varphi_h(t,u)=0$ iff $h(u)\in B_t$.
\end{lemma}

\begin{proof}
Let $\delta_0>0$ be the distance from the compact set $h([0,1])\subseteq V$ to $f[0,1]\cup g[0,1]\cup B_1$. Continuity at $(t_0,u_0)$: if $h(u_0)\in B_{t_0}$, then
$\varphi_h(t_0,u_0)=0$ and $\abs{\varphi_h}$ is continuous. If $h(u_0)\notin B_{t_0}$ and $h(u_0)\notin L_0$, then $h(u_0)\in U_1$ is off $S_{t_0}$, and by
Lemma~\ref{lem:strip-basic}(b) the sign is locally constant. If $h(u_0)\notin B_{t_0}$ and $h(u_0)\in L_0$, then $t_0>0$, and near $h(u_0)$ the only edge of $S_t$
($t$ near $t_0$) is $B_0$; by Lemma~\ref{lem:parity-tools}(c) and (a), nearby points right of $L_0$ have parity $1$, so they lie in $U_t$ and the sign is $-$ near
$(t_0,u_0)$. The boundary signs follow from $U_0=\emptyset$, $h(u)\in L_0\Rightarrow h(u)\in B_0$, and $h(u)\in V$.
\end{proof}

\begin{theorem}[The strip lemma]\label{thm:strip-lemma}
Let $W\subseteq\set{x\ge0}$ be a continuum meeting $L_0$ with $W\cap L_0\subseteq B_0\setminus\set{a_0,b_0}$, disjoint from $f[0,1]$, $g[0,1]$ and $B_1$. If $Z\subseteq W\times W$ is a
continuum with $\pr1(Z)=\pr2(Z)=W$, then $Z$ contains a pair at distance at most $\varepsilon$. In particular the surjective span of $W$ is at most $\varepsilon$.
\end{theorem}

\begin{proof}
\emph{1. $W\setminus L_0\subseteq U_1$.} Let $C$ be a component of $W\setminus L_0$. By boundary bumping $\cl C$ meets $W\cap L_0$ at a point $b$ interior to $B_0$ and away from
the other edges of $S_1$. By Lemma~\ref{lem:parity-tools}(c),(a) the parity is $1$ on a half-disc right of $L_0$ near $b$, which $C$ meets; $C$ is connected and off
$S_1$, so $C\subseteq U_1$.

\emph{2.} Suppose $\alpha=\min_Zd>\varepsilon$, and take $\eta>0$ with $2\eta<\alpha-\varepsilon$ and $\eta$ less than the distance from $W$ to $f[0,1]\cup g[0,1]\cup B_1$.

\emph{3. Two paths.} Pick $z_A\in Z$ with $\pr1(z_A)\in L_0$ and $z_B\in Z$ with $\pr2(z_B)\in L_0$, and an $\eta$-chain $z_A=z_0,z_1,\dots,z_m=z_B$ in $Z$. Let $h$ and $k$ be the
piecewise linear paths through the points $\pr1(z_j)$ and $\pr2(z_j)$. Each segment joins points of $W$ less than $\eta$ apart, so it misses $f$, $g$, $B_1$ and stays in
$\set{x\ge0}$; it either lies in $B_0\setminus\set{a_0,b_0}$ or its part off $L_0$ lies in $U_1$ (by step 1). So $h,k$ map into $V$, $h(0)\in L_0$ and $k(1)\in L_0$, and
$d(h(u),k(u))>\alpha-2\eta>\varepsilon$ for all $u$ ($h(u)-k(u)$ is a convex combination of two consecutive differences).

\emph{4.} Apply Lemma~\ref{lem:common-zero} on $[0,1]\times[0,1]$ (coordinates $(t,u)$) to $\varphi=\varphi_h$ and $\psi=\varphi_k$, whose boundary signs are those of
Lemma~\ref{lem:sign} (the bottom row $u=0$ for $h$, the top row $u=1$ for $k$). At a common zero $(t,u)$, $h(u),k(u)\in B_t$, so $d(h(u),k(u))\le\varepsilon$, contradicting
step 3.
\end{proof}

\formal{\lean{HPC.Span.strip_lemma}.}

\section{Disjoint crossing continua are ordered}

\begin{proposition}\label{prop:ordered}
Let $K,K'$ be disjoint continua in the strip $S=[a,b]\times\R$, each meeting both lines $L_a$ and $L_b$. Then the traces $K\cap L_a$ and $K'\cap L_a$ are separated
by height: one lies strictly below the other.
\end{proposition}

\begin{proof}
Otherwise, after exchanging $K$ and $K'$, there are $p,q\in K\cap L_a$ and $w\in K'\cap L_a$ with $p_2<w_2<q_2$. Let $W$ be the component of $S\setminus K$ containing $K'$;
it is relatively open in $S$, and contains a point $e\in K'\cap L_b$. Small balls about points of $W$, intersected with $S$, are convex and miss $K$, so $W$ is
polygonally connected: there is a polygonal path $\beta$ in $W$ from $w$ to $e$ (not necessarily simple). Close it outside the strip: from $e$ to $(b+1,e_2)$, down to
$(b+1,-N)$, across to $(a-1,-N)$, up to $(a-1,w_2)$, back to $w$, with $N$ larger than all heights involved. $K$ misses this closed polygon $P$, so $\parity(P,p)=\parity(P,q)$.
The segments from $p$ and $q$ to $(a-\frac12,p_2)$ and $(a-\frac12,q_2)$ miss $P$. On the line $x=a-2$ the parity is $0$. Moving from $(a-2,p_2)$ to $(a-\frac12,p_2)$
one crosses only the vertical edge $x=a-1$, whose height range $[-N,w_2]$ contains $p_2$; moving from $(a-2,q_2)$ one crosses nothing, as $q_2>w_2$. So the parities at
$p$ and $q$ are $1$ and $0$: a contradiction.
\end{proof}

\formal{\lean{HPC.Span.crossing_traces_separated}.}

\section{Proof of the strip theorem}

\begin{theorem}\label{thm:ot32}
The strip theorem (Definition~\ref{def:otthm32}) holds.
\end{theorem}

\begin{proof}
Let $J$ be the set of $i$ with $W_i=h_i(X)$ of positive surjective span, and suppose $J$ is uncountable.

\emph{1. One scale.} For $i\in J$ let $\alpha_i$ be the minimum of $d$ on a witness $Z_i$, and $g_i=\pr1h_i(y)-\max_{z\in X}\pr1h_i(z)>0$. For some $n\ge1$ the set $J_1$ of
$i$ with $\alpha_i,g_i\ge\frac1{n+1}$ is uncountable.

\emph{2. Clustering.} Put $\rho=\frac1{4(n+1)}$. Some $m$ is a modulus of uniform continuity of $h_i$ for $\rho$ for uncountably many $i$; on a finite $\frac1m$-net
of $Y$ the maps take values in a separable metric space, so an uncountable $J_3\subseteq J_1$ has $\nrm{h_i(c)-h_j(c)}<3\rho$ for all $i,j\in J_3$ and all $c\in Y$.

\emph{3. Order.} By Proposition~\ref{prop:ordered}, the traces of the $K_i$ on $L_0$ are linearly ordered; take $i_1,i_2,i_3\in J_3$ with the trace of $K_{i_2}$ strictly between
$h_{i_1}(x)$ and $h_{i_3}(x)$.

\emph{4. The strip.} Let $\rho_0$ be a third of the least distance between $K_{i_1},K_{i_2},K_{i_3}$, and take an $\eta'$-chain $x=c_0,\dots,c_s=y$ in $Y$ ($s\ge1$ because
$y\notin X$) fine enough that consecutive images under $h_{i_1}$ and $h_{i_3}$ are less than $\rho_0$ apart. Put $a_j=h_{i_1}(c_j)$, $b_j=h_{i_3}(c_j)$, $\varepsilon=3\rho$.

\emph{5. Hypotheses of the strip lemma for $W=W_{i_2}$.} $\nrm{a_j-b_j}<3\rho$; $a_0,b_0\in L_0$; everything lies in $[0,1]\times\R$; $h_{i_2}(x)\in W\cap L_0$, and $W\cap L_0$ lies
strictly inside $B_0$ by step 3; $f[0,1]$ lies within $\rho_0$ of $K_{i_1}$, so it misses $W\subseteq K_{i_2}$, and likewise $g$; the first coordinates on $B_1=[h_{i_1}(y),h_{i_3}(y)]$
exceed $\pr1h_{i_2}(y)-3\rho\ge\max_W\pr1+\frac1{n+1}-3\rho>\max_W\pr1$, so $B_1$ misses $W$.

\emph{6.} Theorem~\ref{thm:strip-lemma} gives a pair in $Z_{i_2}$ at distance at most $3\rho<\frac1{n+1}\le\alpha_{i_2}$: a contradiction.
\end{proof}

\formal{\lean{HPC.Span.exists_not_countable_of_cover}, \lean{HPC.Span.exists_between_three}, \lean{HPC.Span.ot82_theorem32}.}

\begin{proof}[Proof of Theorem~\ref{thm:theoremC}]
Theorem~\ref{thm:ot32}, Theorem~\ref{thm:ot35} and Corollary~\ref{cor:C-from-35}.
\end{proof}

\formal{\lean{HPC.Span.theoremC}.}

\section*{Exercises}

\begin{exercise}
Show by an example that Lemma~\ref{lem:common-zero} fails if the sign condition for $\varphi$ on the bottom row is dropped.
\end{exercise}

\begin{exercise}
Show that the closed polygon $P$ in the proof of Proposition~\ref{prop:ordered} need not be simple, and explain why this does not matter.
\end{exercise}

%======================================================================
\chapter{Snake-likeness: the plan of Hoehn and Oversteegen}\label{ch:snake}
%======================================================================

The last step of the indecomposable case is a theorem of Hoehn and Oversteegen (2016, Theorem~1): a hereditarily indecomposable
continuum of span zero is snake-like. Together with Theorems~\ref{thm:HI} and~\ref{thm:theoremC} it shows that a nondegenerate homogeneous
indecomposable plane continuum is a pseudo-arc. This chapter states the theorem in the form proved here (for plane continua), reduces it to
a lifting theorem (Hoehn--Oversteegen's Theorem~20), and carries out the reduction. Chapters~\ref{ch:folds}--\ref{ch:ladders} prove
Theorem~20.

The strategy is this. Snake-likeness asks for $\varepsilon$-maps onto $\I$ (\vthree{Theorem~5.10}). Span zero says, roughly, that two points
of $X$ that start together cannot get far apart while one of them sweeps out all of $X$. Hoehn and Oversteegen turn this into a separation
statement (Theorem~4): if a graph $G$ and a path $\gamma$ both approximate $X$, the set of pairs $(x,t)$ with $x$ close to $\gamma(t)$
separates the bottom of $G\times\I$ from the top. Theorem~20 then says that a hereditarily indecomposable space mapped into $G$ can be lifted, up to
a small error, into any neighbourhood of such a separator. The $t$-coordinate of the lift is the required $\varepsilon$-map.

\section{Statements}

\begin{definition}\label{def:plgraph}
Let $E$ be a real normed space.
\begin{enumerate}[label=(\alph*)]
\item A \emph{PL graph} in $E$ is a finite union of nondegenerate segments, two different ones meeting at most in common end points.
\item Let $Y$ be a topological space and $M,P,Q\subseteq Y$. \emph{$M$ separates $P$ from $Q$ in $Y$} if there are open sets $U_P\supseteq P$ and
$U_Q\supseteq Q$ that cover $Y\setminus M$ and whose intersection misses $Y\setminus M$.
\item A \emph{fan at $y$ in $A$ of diameter $<r$} is a union $\set y\cup\bigcup_k[y,z_k]$ of finitely many segments meeting pairwise only in $y$,
contained in $A$, which is a neighbourhood of $y$ in $A$ and has diameter $<r$.
\end{enumerate}
\end{definition}

\formal{\lean{HPC.HO16.IsPLGraph}, \lean{HPC.HO16.SeparatesIn}, \lean{HPC.HO16.IsFanAt}.}

A PL graph is compact, and every point of a PL graph has fans of every diameter (cut the segments through $y$ short). If $R_0\supseteq P$ and $R_1\supseteq
Q$ are disjoint closed sets covering $Y\setminus M$, then $M$ separates $P$ from $Q$ (take $U_P=Y\setminus R_1$, $U_Q=Y\setminus R_0$); this is how
separation is always produced below.

\formal{\lean{HPC.HO16.IsPLGraph.isCompact}, \lean{HPC.HO16.IsPLGraph.exists_fan}, \lean{HPC.HO16.separatesIn_of_closed}.}

Throughout Chapters~\ref{ch:snake}--\ref{ch:ladders}, $G$ is a PL graph, $G^\star=G\times\I\subseteq E\times\R$, and $\pi\colon E\times\R\to E$ is the
first projection.

\begin{theorem}[Hoehn--Oversteegen, Theorem~20, forward direction]\label{thm:ho20}
Let $X$ be a compact metric space that is hereditarily indecomposable, $G$ a PL graph in a normed space $E$, and $f\colon X\to G$ continuous.
Let $M\subseteq G\times(0,1)$ separate $G\times\set0$ from $G\times\set1$ in $G^\star$, and let $U\supseteq M$ be open in $E\times\R$. Then for every $\varepsilon>0$
there is a continuous $h\colon X\to U\cap G^\star$ with $\nrm{f(x)-\pi h(x)}<\varepsilon$ for all $x$.
\end{theorem}

\formal{The statement is the proposition \lean{HPC.HO16.HOThm20}; it is proved as \lean{HPC.HO16.ho16_theorem20} (Chapter~\ref{ch:ladders}).}

Hoehn and Oversteegen prove an equivalence: Theorem~20 characterizes hereditary indecomposability. Only the direction stated is needed and
proved. ``Hereditarily indecomposable'' for a compactum (not necessarily connected) means that every subcontinuum is indecomposable.

\begin{theorem}[Hoehn--Oversteegen, Theorem~1, for plane continua]\label{thm:ho1}
A hereditarily indecomposable plane continuum of span zero is snake-like.
\end{theorem}

\formal{\lean{HPC.HO16.ho16_theorem1_plane}, which takes Theorem~\ref{thm:ho20} as a hypothesis \lean{h20 : HOThm20}; span zero is
\lean{HPC.Span.HOSpanZero} (Definition~\ref{def:span}).}

\begin{theorem}\label{thm:indec-pseudoarc}
A nondegenerate homogeneous indecomposable plane continuum $M$ is a pseudo-arc, hence homeomorphic to the plane pseudo-arc $M_0$. Conversely
a pseudo-arc is homogeneous and indecomposable. So a nondegenerate plane continuum is homogeneous and indecomposable iff it is a pseudo-arc.
\end{theorem}

\begin{proof}
$M$ is hereditarily indecomposable (Theorem~\ref{thm:HI}) and has span zero (Theorem~\ref{thm:theoremC}), so it is snake-like by
Theorem~\ref{thm:ho1}; it is nondegenerate by hypothesis, so it is a pseudo-arc (\vone{Definition~2.14}). Uniqueness
(\vtwo{Corollary~8.6}) gives the homeomorphism with $M_0$. Conversely, a pseudo-arc is homogeneous (\vtwo{Theorem~8.4}) and,
being a hereditarily indecomposable continuum, indecomposable.
\end{proof}

\formal{\lean{HPC.HO16.isPseudoArc_of_indec_homogeneous} and \lean{HPC.HO16.homeomorph_pseudoArc_of_indec_homogeneous} (with
\lean{h20} as hypothesis); unconditionally \lean{HPC.HO16.ho16_indec_homogeneous_pseudoArc}; the converse
\lean{HPC.HO16.ho16_pseudoArc_indec_homogeneous}; both directions \lean{HPC.HO16.ho16_indec_homogeneous_iff}. Non-vacuity of the
hypotheses: \lean{HPC.HO16.exists_indec_homogeneous_plane}.}

\begin{remark}\label{rem:ho1-point}
Hoehn and Oversteegen restate Theorem~1 on their p.~211 as: ``a continuum $X$ is homeomorphic to the pseudo-arc iff $X$ is hereditarily
indecomposable and has span zero'', and Theorem~$1'$ on p.~212 with surjective semi-span in place of span. Read literally, both fail for a one-point
continuum, which is hereditarily indecomposable with span zero but is not a pseudo-arc. The nondegeneracy hypothesis is needed.
\end{remark}

\formal{\lean{HPC.HO16.ho16_theorem1_printed_false}, \lean{HPC.HO16.ho16_theorem1prime_printed_false}.}

\section{Theorem 4: span zero gives separation}

\begin{theorem}[Hoehn--Oversteegen, Theorem~4, path form]\label{thm:ho4}
Let $E$ be a proper metric space, $X\subseteq E$ compact of span zero (in the sense of Hoehn--Oversteegen), and $\varepsilon>0$. There is $\delta>0$
such that: for every compact $G\subseteq B_\delta(X)$ and every path $\gamma\colon\I\to E$ with $\gamma\I\subseteq B_\delta(X)$ and $X\subseteq
B_\delta(\gamma\I)$, the set
\[
M_\gamma=\set{(x,t)\in G\times(0,1) : d(x,\gamma(t))<\varepsilon}
\]
separates $G\times\set0$ from $G\times\set1$ in $G^\star$.
\end{theorem}

Here $B_\delta(A)$ is the open $\delta$-neighbourhood of $A$. Compared with Hoehn and Oversteegen: the arc $I$ with end points $p,q$ is replaced by a path
on $\I$, and $I\setminus\set{p,q}$ by $(0,1)$; of $G$ only $G\subseteq B_\delta(X)$ is used; $X$ need not be connected; and $E$ is any proper metric
space rather than the Hilbert cube.

\begin{proof}
Suppose the theorem fails for $\varepsilon$. Then for each $n$ there are $G_n$ and $\gamma_n$ as in the statement with $\delta=1/(n+1)$ such that $M_n=M_{\gamma_n}$
does not separate.

\emph{Step 1: a continuum far from the path.} Let $D_n=\set{(x,t)\in G_n\times\I : d(x,\gamma_n(t))\ge\varepsilon}$, a compact set. Some
subcontinuum $Z_n\subseteq D_n$ meets both rows $t=0$ and $t=1$. Otherwise, by the cut-wire theorem (\vthree{Theorem~1.9}) applied in $D_n$, $D_n=D^0\cup
D^1$ with $D^0,D^1$ disjoint and closed, the row-$0$ points of $D_n$ in $D^0$ and the row-$1$ points in $D^1$. Then $R_0=D^0\cup(G_n\times\set0)$ and
$R_1=D^1\cup(G_n\times\set1)$ are disjoint closed sets covering $G_n^\star\setminus M_n$ (a point of $G_n\times(0,1)$ outside $M_n$ lies in $D_n$; a point
of an end row lies in $R_0$ or $R_1$), so $M_n$ separates, a contradiction. Since $Z_n$ is connected and meets both rows, its second
projection is $\I$.

\emph{Step 2.} $W_n=\set{(x,\gamma_n(t)) : (x,t)\in Z_n}\subseteq E\times E$ is a continuum on which $d\ge\varepsilon$, with first coordinates in
$B_{1/(n+1)}(X)$ and second coordinates exactly $\gamma_n\I$. Every point of $X$ is within $1/(n+1)$ of a second coordinate.

\emph{Step 3: a limit.} All $W_n$ lie in $K\times K$, where $K=\cl{B_1(X)}$ is compact because $E$ is proper. Nonempty compact subsets of a compact
metric space form a compact space in the Hausdorff metric, so a subsequence converges to a compactum $W$. $W$ is connected (a limit of continua),
$d\ge\varepsilon$ on $W$, the first coordinates of $W$ lie in $X$, and the second coordinates form exactly $X$: they lie in $X$ since
$\gamma_n\I\subseteq B_{1/(n+1)}(X)$, and every $x\in X$ is a limit of second coordinates of $W_n$.

\emph{Step 4.} $W$ is a subcontinuum of $X\times X$ with $\pr1(W)\subseteq\pr2(W)=X$ that misses the diagonal. This contradicts span zero.
\end{proof}

\formal{\lean{HPC.HO16.ho16_theorem4}; Step~1 is \lean{HPC.HO16.thm4_continuum_of_not_separates}, Steps~2--4
\lean{HPC.HO16.thm4_limit_contra}, using \lean{HPC.HO16.subset_of_tendsto_of_eventually} for limits of compacta.}

\section{Plane approximations}

\begin{lemma}\label{lem:hi-interior}
A hereditarily indecomposable plane continuum $X$ has empty interior.
\end{lemma}

\begin{proof}
If $X$ contains a closed square $Q$, then $Q$ is the union of its left and right closed halves, two proper subcontinua that meet; so $Q$ is a
decomposable subcontinuum of $X$.
\end{proof}

\formal{\lean{HPC.HO16.interior_eq_empty_of_hi}.}

\begin{lemma}[Grid graphs]\label{lem:grid}
Let $X\subseteq\R^2$ be compact with empty interior and $\delta>0$. There are a PL graph $G\subseteq B_\delta(X)$ and a map $f\colon X\to G$, continuous,
with $d(f(x),x)<\delta$ for all $x\in X$.
\end{lemma}

\begin{proof}
Use the grid of mesh $\eta=\delta/2$: squares $Q_{ij}=[i\eta,(i+1)\eta]\times[j\eta,(j+1)\eta]$, $i,j\in\Z$. Let $\mathcal Q$ be the finite set of squares
meeting $X$. For each square choose a point $c_Q$ in its interior but not in $X$ (possible since $X$ has empty interior; for squares missing $X$ any
interior point will do).

The radial projection $r_Q$ from $c_Q$ onto the frontier of $Q=[p_1,q_1]\times[p_2,q_2]$ is $r_Q(x)=c_Q+\mu(x-c_Q)^{-1}(x-c_Q)$, where
\[
\mu(v)=\max\Bigl(\frac{v_1}{q_1-c_1},\ \frac{-v_1}{c_1-p_1},\ \frac{v_2}{q_2-c_2},\ \frac{-v_2}{c_2-p_2}\Bigr)
\]
is the gauge of $Q$ centred at $c_Q$. The gauge is continuous, positively homogeneous, positive for $v\neq0$; $\mu(x-c_Q)\le1$ iff $x\in Q$, with
equality iff $x$ lies on the frontier. So $r_Q$ is continuous off $c_Q$, maps into the frontier of $Q$, and is the identity on it.

Put $f=r_Q$ on $X\cap Q$ for $Q\in\mathcal Q$. Two grid squares meet only in their frontiers, where both projections are the identity. As finitely many
closed pieces cover $X$, $f$ is well defined and continuous. Let $G$ be the union of the edges of the squares of $\mathcal Q$; two distinct grid edges
meet at most in a common end, so $G$ is a PL graph. Each square has diameter at most $\eta$ in the sup-norm and meets $X$, which gives $G\subseteq
B_\delta(X)$ and $d(f(x),x)\le\eta<\delta$.
\end{proof}

\formal{\lean{HPC.HO16.exists_plGraph_approx}, with \lean{HPC.HO16.radProj}, \lean{HPC.HO16.rectGauge}, \lean{HPC.HO16.radProj_of_onBdry},
\lean{HPC.HO16.radProj_mem}, \lean{HPC.HO16.gridEdge_inter_subset} (in \lean{HPC/HO16/Grid.lean}).}

The plane has the sup-norm in the formalization; distances and diameters are with respect to it. Nothing in this chapter depends on the
choice of norm.

\begin{lemma}\label{lem:path-approx}
Let $X$ be a compact connected metric subset of the plane and $\delta>0$. There is a path $\gamma\colon\I\to\R^2$ with $\gamma\I\subseteq B_\delta(X)$ and
$X\subseteq B_\delta(\gamma\I)$.
\end{lemma}

\begin{proof}
Take a finite $\delta/2$-net of $X$ and join its points in turn by $\delta/2$-chains in $X$ (possible since $X$ is connected). Let $\gamma$ be the
polygonal path through the resulting list of points. Each of its segments has both ends in $X$ and length $<\delta/2$.
\end{proof}

\formal{\lean{HPC.HO16.exists_path_approx}, with \lean{HPC.HO16.exists_chain_through}.}

\section{The proof of Theorem 1 from Theorem 20}

\begin{proof}[Proof of Theorem~\ref{thm:ho1}]
Let $X\subseteq\R^2$ be a hereditarily indecomposable continuum of span zero and $\varepsilon>0$. By \vthree{Theorem~5.10} it suffices to find a map
$\psi\colon X\to\R$ such that $\psi(x)=\psi(x')$ implies $d(x,x')<\varepsilon$.
\begin{enumerate}
\item Theorem~\ref{thm:ho4} for $X$ and $\varepsilon/6$ gives $\delta_0$. Put $\delta=\min(\delta_0,\varepsilon/6)$.
\item By Lemma~\ref{lem:hi-interior}, $X$ has empty interior. Lemma~\ref{lem:grid} gives a PL graph $G\subseteq B_\delta(X)$ and $f\colon X\to G$ with
$d(f(x),x)<\delta$; Lemma~\ref{lem:path-approx} gives a path $\gamma$ with $\gamma\I\subseteq B_\delta(X)$ and $X\subseteq B_\delta(\gamma\I)$.
\item By Theorem~\ref{thm:ho4}, $M'=\set{(y,t)\in G\times(0,1) : d(y,\gamma(t))<\varepsilon/6}$ separates $G\times\set0$ from $G\times\set1$ in $G^\star$.
\item The set $U=\set{(y,t) : 0<t<1,\ d(y,\gamma(t))<\varepsilon/6}$ is open and contains $M'$.
\item Theorem~\ref{thm:ho20}, for $X$, $f$, $M'$, $U$ and $\varepsilon/6$, gives a continuous $h\colon X\to U\cap G^\star$ with $d(f(x),\pi h(x))<\varepsilon/6$.
\item Let $\psi(x)$ be the second coordinate of $h(x)$. If $\psi(x)=t$, then, as $h(x)\in U$,
\[
d(x,\gamma(t))\le d(x,f(x))+d(f(x),\pi h(x))+d(\pi h(x),\gamma(t))<\delta+\tfrac\varepsilon6+\tfrac\varepsilon6\le\tfrac\varepsilon2 .
\]
So $\psi(x)=\psi(x')=t$ gives $d(x,x')<\varepsilon$.\qedhere
\end{enumerate}
\end{proof}

\formal{\lean{HPC.HO16.ho16_theorem1_plane}; the $\varepsilon$-map criterion is \lean{HPC.Plane.isSnakeLike_univ_of_epsMap}.}

Only the plane is needed for the classification, and only there are the approximations of Lemmas~\ref{lem:grid} and~\ref{lem:path-approx} carried
out. Theorem~\ref{thm:ho4} and Theorem~\ref{thm:ho20} are proved in general.

\section*{Exercises}

\begin{exercise}
Show that in Step~1 of the proof of Theorem~\ref{thm:ho4} the end rows themselves are never used: the cut-wire theorem is applied to $D_n$.
Why is it nevertheless important that $M_\gamma$ is defined with $t\in(0,1)$ rather than $t\in\I$?
\end{exercise}

\begin{exercise}
Show that an arc has span zero and that a circle does not. (For the circle, use a half turn.)
\end{exercise}

\begin{exercise}
Let $X=[0,1]\subseteq\R$, $G=X$ and $\gamma(t)=t$. Describe $M_\gamma$ and check directly that it separates $G\times\set0$ from $G\times\set1$
when $\varepsilon>0$.
\end{exercise}

\begin{exercise}
In Lemma~\ref{lem:grid}, why can $f$ not be taken to be the nearest-point retraction onto the union of the grid edges?
\end{exercise}

\begin{exercise}
Prove that the radial projection $r_Q$ is the identity on the frontier of $Q$ directly from the formula for $\mu$.
\end{exercise}

%======================================================================
\chapter{Simple folds}\label{ch:folds}
%======================================================================

A simple fold of a graph $G$ is a graph $F$ obtained by cutting $G$ open along two finite sets and gluing in an extra sheet, so that the projection
$\varphi\colon F\to G$ runs forward over one part of $G$, back over a middle part, and forward again. Hoehn and Oversteegen's Theorem~8 says that a
map of a hereditarily indecomposable compactum into $G$ lifts, up to $\varepsilon$, through every simple fold. This is the only place where hereditary
indecomposability enters the proof of Theorem~\ref{thm:ho20}; it enters through a lemma of Krasinkiewicz and Minc.

Hoehn and Oversteegen build a fold as a quotient space. Here a fold over a PL graph $G\subseteq E$ is a PL graph in $E\times\R$, made of three sheets that
are graphs of height functions, and $\varphi=\pi$ is the first projection. Folds of folds then live in $(E\times\R)\times\R$, and so on; this is why
everything is done in an arbitrary normed space $E$.

\section{The lemma of Krasinkiewicz and Minc}

\begin{lemma}[Krasinkiewicz--Minc 1977, Lemma~2, for compacta]\label{lem:km}
Let $X$ be a hereditarily indecomposable compact metric space, $A,B\subseteq X$ disjoint closed sets, and $U\supseteq A$, $V\supseteq B$ open. Then there are closed
sets $X_1,X_2,X_3$ with
\[
X=X_1\cup X_2\cup X_3,\quad A\subseteq X_1,\quad B\subseteq X_3,\quad X_1\cap X_3=\emptyset,\quad X_1\cap X_2\subseteq V,\quad X_2\cap X_3\subseteq U.
\]
\end{lemma}

In Krasinkiewicz and Minc's words, $(X,A,B)$ is \emph{crooked} between $U$ and $V$. They state the lemma for continua; Hoehn and Oversteegen remark
that the proof works for compacta. Their proof uses a further lemma (a decomposable continuum irreducible between $A$ and $B$); here the decomposable
continuum $C_1\cup C_2$ of Step~3 is enough.

\begin{proof}
\emph{Step 1.} Take open $U'\supseteq A$, $V'\supseteq B$ with $\cl{U'}\subseteq U$, $\cl{V'}\subseteq V$ and $\cl{U'}\cap\cl{V'}=\emptyset$.

\emph{Step 2.} Let $A^*$ be the union of the components of $X\setminus V'$ that meet $A$, and $B^*$ the union of the components of $X\setminus U'$ that meet $B$.
Both are closed: in a compact space, the union of the components of a closed set $K$ that meet a closed $C\subseteq K$ is closed, since a component of $K$
missing $C$ lies in a set clopen in $K$ that misses $C$ (\vthree{Theorem~1.9}).

\emph{Step 3: $A^*\cap B^*=\emptyset$.} Let $p$ be in both, $C_1$ its component in $X\setminus V'$ and $C_2$ its component in $X\setminus U'$. $C_1$ meets $A$ and
misses $B\subseteq V'$; $C_2$ meets $B$ and misses $A$. So $C_1\cup C_2$ is a continuum (both contain $p$), and $C_1$, $C_2$ are proper subcontinua whose union
is all of it: $C_1\cup C_2$ is decomposable, which is impossible.

\emph{Step 4.} A continuum in $X\setminus V'$ meeting $A^*$ lies in a component meeting $A$, so it misses $B^*$. By the cut-wire theorem
(\vthree{Theorem~1.9}), $X\setminus V'=D_1\cup D_2$ with $D_1,D_2$ disjoint and closed, $A^*\subseteq D_1$ and $B^*\setminus V'\subseteq D_2$. Then $B^*\cap D_1=\emptyset$.

\emph{Step 5.} In the same way $X\setminus U'=E_1\cup E_2$ with $E_1,E_2$ disjoint and closed, $B^*\subseteq E_1$ and $D_1\setminus U'\subseteq E_2$.

\emph{Step 6.} Put
\[
X_1=D_1\cup(E_2\cap\cl{V'}),\quad X_2=(D_2\cap E_2)\cup(D_2\cap\cl{U'})\cup(E_2\cap\cl{V'}),\quad X_3=E_1\cup(D_2\cap\cl{U'}).
\]
They cover $X$: a point outside $U'\cup V'$ lies in some $D_i\cap E_j$, and $U'\subseteq D_1\cup(D_2\cap U')$, $V'\subseteq E_1\cup(E_2\cap V')$. The
intersection $X_1\cap X_2$ lies in $\cl{V'}$ and $X_2\cap X_3$ in $\cl{U'}$, using $D_1\cap D_2=E_1\cap E_2=\cl{U'}\cap\cl{V'}=\emptyset$. Finally
$X_1\cap X_3=D_1\cap E_1=(D_1\setminus U')\cap E_1\subseteq E_2\cap E_1=\emptyset$, the other three intersections being empty for the same reasons.
\end{proof}

\formal{\lean{HPC.HO16.km77_lemma2}, with \lean{HPC.HO16.isClosed_biUnion_connectedComponentIn} for Step~2.}

\section{Theorem 8 before folds}

Let $G_1,G_3\subseteq E$ be closed, $G=G_1\cup G_3$, $G_2=G_1\cap G_3$, and suppose
\begin{equation}\label{eq:F3}
\cl{G_1\setminus G_2}\cap\cl{G_3\setminus G_2}=\emptyset .
\end{equation}
Write $B_1=\cl{G_1\setminus G_2}\cap G_2$ and $B_3=\cl{G_3\setminus G_2}\cap G_2$: the points where the first, respectively the third, part is attached to the middle.

\begin{theorem}[Hoehn--Oversteegen, Theorem~8, core]\label{thm:ho8-core}
In this situation assume that $B_1$ and $B_3$ are finite and that every point of $B_1\cup B_3$ has fans in $G_2$ of every diameter. Let $X$ be a hereditarily
indecomposable compact metric space, $f\colon X\to G$ continuous and $\varepsilon>0$. Then there are a continuous $h\colon X\to E$ and closed sets
$X_1,X_2,X_3$ covering $X$ with $X_1\cap X_3=\emptyset$, $\nrm{f(x)-h(x)}<\varepsilon$,
\[
h(X_i)\subseteq G_i\ (i=1,2,3),\qquad h(X_1\cap X_2)\subseteq B_3,\qquad h(X_2\cap X_3)\subseteq B_1 .
\]
\end{theorem}

\begin{proof}
\emph{Fans.} $B_1\cap B_3=\emptyset$ by \eqref{eq:F3}. Choose $r\le\varepsilon$ smaller than half the distance between any two points of the finite set $B_1\cup
B_3$, and for $y\in B_1\cup B_3$ a fan $Y_y$ at $y$ in $G_2$ of diameter $<r$. The fans are pairwise disjoint, closed, and star-shaped about their centres:
$[y,p]\subseteq Y_y$ for $p\in Y_y$. Let $c(p)=y$ for $p\in Y_y$.

\emph{Neighbourhoods of the two sides.} Put $C_1=\cl{G_1\setminus G_2}$, $C_3=\cl{G_3\setminus G_2}$, $K_1=f^{-1}(C_1)$, $K_3=f^{-1}(C_3)$, disjoint closed sets.
There is an open $N_1\supseteq C_1$ with $G\cap N_1\subseteq C_1\cup\bigcup_{y\in B_1}Y_y$: a point of $C_1$ outside $G_2$ lies outside the closed set $G_3$, so
near it $G$ is inside $G_1\setminus G_2$; a point $y\in B_1$ is not in $C_3$, and $Y_y$ is a neighbourhood of $y$ in $G_2$. Likewise $N_3$. Let $U_1\supseteq K_1$,
$U_3\supseteq K_3$ be disjoint open sets inside $f^{-1}(N_1)$ and $f^{-1}(N_3)$.

\emph{Krasinkiewicz--Minc.} Lemma~\ref{lem:km} with $A=K_1$, $B=K_3$, $U=U_1$, $V=U_3$ gives closed $X_1\supseteq K_1$, $X_2$, $X_3\supseteq K_3$ with $X_1\cap
X_3=\emptyset$, $X_1\cap X_2\subseteq U_3$ and $X_2\cap X_3\subseteq U_1$.

\emph{The correction.} Let $W_1=U_1\setminus X_1$ and $W_3=U_3\setminus X_3$, disjoint open sets. A point $x\in W_1$ is not in $K_1$, so $f(x)\in N_1\cap G$ but
$f(x)\notin C_1$: hence $f(x)\in Y_y\setminus\set y$ for a unique $y\in B_1$. Since the fans are disjoint and closed, $x\mapsto c(f(x))$ is locally constant on $W_1$.
Let $u_1\colon X\to\I$ be an Urysohn function with $u_1=0$ on $X_2\cap X_3$ and $u_1=1$ off $U_1$; similarly $u_3$, with $u_3=0$ on $X_1\cap X_2$ and $u_3=1$ off
$U_3$. Define
\[
h(x)=\begin{cases} f(x)+(1-u_1(x))\bigl(c(f(x))-f(x)\bigr), & x\in W_1,\\ f(x)+(1-u_3(x))\bigl(c(f(x))-f(x)\bigr), & x\in W_3,\\ f(x) & \text{otherwise.}\end{cases}
\]
On $W_1$, $h(x)$ lies on the segment from $f(x)$ to its fan centre, so inside the fan, within $r\le\varepsilon$ of $f(x)$; it equals the centre where $u_1=0$.

\emph{Continuity.} The correction term on $W_1$ is continuous there and bounded by $r$. At a point $x\notin W_1$ either $x\notin U_1$, where $u_1(x)=1$ and the
term tends to $0$; or $x\in X_1\cap U_1$. In the latter case $x\notin X_2\cup X_3$ (as $X_1\cap X_2\subseteq U_3$ misses $U_1$, and $X_1\cap X_3=\emptyset$), so a
neighbourhood of $x$ lies in $X_1$ and misses $W_1$. The same holds for $W_3$.

\emph{Checks.} On $X_2\cap X_3\subseteq W_1$ we have $u_1=0$, so $h$ is a fan centre in $B_1$; likewise $h(X_1\cap X_2)\subseteq B_3$. If $x\in X_1$ and $h(x)\ne f(x)$,
then $h(x)$ is in a fan, inside $G_2\subseteq G_1$; if $h(x)=f(x)$ and $x\notin U_3$, then $x\notin K_3$, so $f(x)\notin G_3\setminus G_2$ and $f(x)\in G_1$; and
$X_1\cap U_3\subseteq W_3$ unless $x\in X_3$, which is impossible. The cases of $X_3$ and $X_2$ are similar (off $U_1\cup U_3$, $f(x)\notin C_1\cup C_3$, so $f(x)\in G_2$).
\end{proof}

\formal{\lean{HPC.HO16.ho16_theorem8_core}, with \lean{HPC.HO16.exists_center_fun}, \lean{HPC.HO16.continuousOn_center_comp},
\lean{HPC.HO16.continuous_correction} and \lean{HPC.HO16.exists_open_nbhd_side}.}

Hoehn and Oversteegen obtain the map near the attaching points from Tietze extensions on each clopen piece $f^{-1}((y,z])\cap W_1$; they do not say that
the fans must be disjoint, which their construction needs. The linear push towards the fan centre used here avoids the pieces altogether.

\section{PL graphs}

The combinatorics of folds and stairwells rests on a library for PL graphs. A \emph{decomposition} of a PL graph $G$ is a finite set of segments as in
Definition~\ref{def:plgraph}(a); its end points are the \emph{vertices}.

\begin{proposition}\label{prop:pl-library}
Let $G$ be a PL graph in $E$.
\begin{enumerate}[label=(\alph*)]
\item \emph{(Unions.)} A finite union of segments is a PL graph. The union of two PL graphs is a PL graph.
\item \emph{(Refinement.)} For PL graphs $A_1,\dots,A_k\subseteq G$ and a finite $P\subseteq G$ there is a decomposition of $G$ in which every $A_i$ is a union of
edges, an edge not in $A_i$ meets $A_i$ only in vertices, and every point of $P$ is a vertex.
\item \emph{(Frontiers.)} For a PL graph $A\subseteq G$, $\cl{G\setminus A}$ is a PL graph, and the relative frontier $\rfr A=A\cap\cl{G\setminus A}$ is finite.
\item \emph{(Components.)} $G$ has finitely many components, each a PL graph and open in $G$.
\item \emph{(Projections.)} If $F\subseteq E\times\R$ is a PL graph on which $\pi$ is injective, then $\pi(F)$ is a PL graph.
\item \emph{(Common subdivisions.)} Any two decompositions of $G$ have a common subdivision.
\end{enumerate}
\end{proposition}

\begin{proof}[Proof sketch]
(a) The intersection of two segments is empty or a segment (possibly a point). Let $V$ be the finite set of all end points of the given segments and of
their pairwise intersections; subdivide each segment at the points of $V$ on it and keep each piece once. Two pieces sharing a point that is not an end of
both would lie in a common intersection $[p,q]$ with $p,q\in V$, and both are cut at $V\cap[p,q]$, so they coincide. (b)--(f) follow from (a) by taking the
vertices of a suitable common subdivision: frontiers of subgraphs consist of vertices, and a component is the union of a class of edges under ``sharing a
vertex''.
\end{proof}

\formal{\lean{HPC.HO16.isPLGraph_biUnion}, \lean{HPC.HO16.IsPLGraph.union}, \lean{HPC.HO16.IsPLGraph.exists_refinement},
\lean{HPC.HO16.IsPLGraph.closure_diff}, \lean{HPC.HO16.IsPLGraph.relFrontier_finite}, \lean{HPC.HO16.IsPLGraph.connectedComponentIn},
\lean{HPC.HO16.IsPLGraph.finite_components}, \lean{HPC.HO16.IsPLGraph.image_fst}, \lean{HPC.HO16.Subdiv.exists_common_subdivision}; the relative
frontier is \lean{HPC.HO16.relFrontier}.}

For $A\subseteq G$ the relative frontier in general is $\rfr A=\cl A\cap\cl{G\setminus A}$; for closed $A$ it is $A\cap\cl{G\setminus A}$.

\section{Simple folds}

\begin{definition}\label{def:fold}
\begin{enumerate}[label=(\alph*)]
\item A \emph{fold triple} on a PL graph $G$ is a triple $(G_1,G_2,G_3)$ of nonempty PL graphs with $G_1\cup G_3=G$, $G_2=G_1\cap G_3$ and
$\cl{G_1\setminus G_2}\cap\cl{G_3\setminus G_2}=\emptyset$.
\item A \emph{simple fold} on $G$ along a fold triple is a triple $(F_1,F_2,F_3)$ of PL graphs in $E\times\R$ such that $\pi$ is injective on each $F_i$ with
$\pi(F_i)=G_i$, and
\[
\pi(F_1\cap F_2)=\rfr G_1,\qquad \pi(F_2\cap F_3)=\rfr G_3,\qquad F_1\cap F_3=\emptyset .
\]
The fold itself is $F=F_1\cup F_2\cup F_3$, with projection $\varphi=\pi|_F$.
\end{enumerate}
\end{definition}

\formal{\lean{HPC.HO16.IsFoldTriple}, \lean{HPC.HO16.IsSimpleFoldOf}.}

Hoehn and Oversteegen ask the $G_i$ to be \emph{regular}: closed, with finitely many components, all nondegenerate. For subsets of a PL graph this amounts
to being a PL graph. Their conditions (F1)--(F3) are part (a); (F4) and (F5) are part (b).

\begin{theorem}[Folds exist]\label{thm:fold-exists}
Every fold triple has a simple fold.
\end{theorem}

\begin{proof}
Put $B_3=\rfr G_1$ and $B_1=\rfr G_3$. Both are finite subsets of $G_2$ (Proposition~\ref{prop:pl-library}(c)), and they are disjoint, since
$\rfr G_1\subseteq\cl{G_3\setminus G_2}$ and $\rfr G_3\subseteq\cl{G_1\setminus G_2}$.

Take a decomposition of $G$ in which $G_1,G_2,G_3$ are unions of edges and all points of $B_1\cup B_3$ are vertices
(Proposition~\ref{prop:pl-library}(b)), and subdivide every edge at its midpoint. Now no edge has both ends in $B_3$, nor both in $B_1$. Heights at vertices:
$h_1(v)=1$ for $v\in B_3$ and $0$ otherwise; $h_3(v)=1$ for $v\in B_1$ and $2$ otherwise. Let $F_1$ be the union, over the edges $[a,b]\subseteq G_1$, of the
segments $[(a,h_1(a)),(b,h_1(b))]$; $F_2$ the union over the edges of $G_2$ of $[(a,1),(b,1)]$; and $F_3$ the union over the edges of $G_3$ of
$[(a,h_3(a)),(b,h_3(b))]$.

Lifts of different edges meet only over common vertices, where the heights agree; so each $F_i$ is a PL graph on which $\pi$ is injective with image $G_i$.
Over a point $x$ of an edge of $G_2$, the $F_1$-height is $\le1$, with equality only if $x$ is a vertex in $B_3$ (an interior point of an edge has height $<1$,
as at most one end is in $B_3$); the $F_2$-height is $1$; and the $F_3$-height is $\ge1$, with equality only at vertices in $B_1$. Hence $\pi(F_1\cap F_2)=B_3$,
$\pi(F_2\cap F_3)=B_1$ and $F_1\cap F_3=\emptyset$.
\end{proof}

\formal{\lean{HPC.HO16.exists_simpleFold}, with the lifted graphs \lean{HPC.HO16.liftG} and the midpoint subdivision
\lean{HPC.HO16.exists_midpoint_subdivision}.}

\begin{lemma}[Hoehn--Oversteegen, Lemma~6]\label{lem:ho6}
Let $(F_1,F_2,F_3)$ be a simple fold on $G$ along $(G_1,G_2,G_3)$, and $F=F_1\cup F_2\cup F_3$. Relative frontiers of subsets of $G$ are taken in $G$, those of subsets
of $F$ in $F$.
\begin{enumerate}[label=(\arabic*)]
\item $\rfr G_2=\rfr G_1\cup\rfr G_3$, and $\rfr G_1\cap\rfr G_3=\emptyset$.
\item $\rfr(G_1\setminus G_2)=\rfr G_3$ and $\rfr(G_3\setminus G_2)=\rfr G_1$.
\item[(4)] $F_1\cap F_2$ and $F_2\cap F_3$ are finite.
\item[(5)] $\rfr F_1=F_1\cap F_2$, $\rfr F_3=F_2\cap F_3$, and $\rfr F_2=(F_1\cap F_2)\cup(F_2\cap F_3)$.
\item[(6)] $F_1\cap F_2$ separates $F_1\setminus F_2$ from $(F_2\cup F_3)\setminus F_1$ in $F$, and $F_2\cap F_3$ separates $F_3\setminus F_2$ from $(F_1\cup F_2)\setminus F_3$.
\item[(7)] At a point $p\in F_1\setminus F_2$, $\varphi$ maps neighbourhoods of $p$ in $F_1$ onto neighbourhoods of $\varphi(p)$ in $G$; likewise at $p\in F_2\setminus(F_1\cup F_3)$
for $F_2$ and at $p\in F_3\setminus F_2$ for $F_3$.
\end{enumerate}
\end{lemma}

(Item (3) of Hoehn and Oversteegen, that the $F_i$ are regular, is part of the definition.)

\begin{proof}
(1), (2): $G\setminus G_1=G_3\setminus G_2$ and $G\setminus G_3=G_1\setminus G_2$, and $\cl{G_i\setminus G_2}\subseteq G_i$. So $\rfr G_1=G_2\cap\cl{G_3\setminus G_2}$,
$\rfr G_3=G_2\cap\cl{G_1\setminus G_2}$, which are disjoint by (F3), and $\rfr G_2$ is their union; and $\cl{G\setminus(G_1\setminus G_2)}=G_3$.
(4): $\pi$ is injective on $F_1$ and maps $F_1\cap F_2$ into the finite set $\rfr G_1$.
(5): $F_1\cap(F_2\cup F_3)=F_1\cap F_2$ is finite. A point of $F_1\setminus F_2$ has positive distance from the compact set $F_2\cup F_3$. A point of $F_1\cap F_2$ is a limit of
points of $F_2\setminus F_1$, as a PL graph has no isolated points.
(6): take the open sets $(F_2\cup F_3)^c$ and $F_1^c$.
(7): for $p\in F_1\setminus F_2$, $\varphi(p)\notin\rfr G_1$ (the $F_1$-point over a point of $\rfr G_1$ is in $F_2$), so $\varphi(p)$ has a neighbourhood in $G$ inside $G_1$,
and $\varphi|_{F_1}$ is a homeomorphism onto $G_1$.
\end{proof}

\formal{\lean{HPC.HO16.ho16_lemma6}; Hoehn and Oversteegen leave the proof to the reader.}

\begin{theorem}[Hoehn--Oversteegen, Theorem~8, (1)$\Rightarrow$(2)]\label{thm:ho8}
Let $X$ be a hereditarily indecomposable compact metric space, $(F_1,F_2,F_3)$ a simple fold on $G$, $f\colon X\to G$ continuous and $\varepsilon>0$. There is a
continuous $g\colon X\to F$ with $\nrm{f(x)-\varphi(g(x))}<\varepsilon$ for all $x$.
\end{theorem}

\begin{proof}
Theorem~\ref{thm:ho8-core} applies: $B_1=\cl{G_1\setminus G_2}\cap G_2=\rfr G_3$ and $B_3=\rfr G_1$ are finite by Lemma~\ref{lem:ho6}, and points of the PL graph $G_2$ have fans
of every diameter. It gives $h$ and $X_1,X_2,X_3$. The inverse $\psi_i$ of $\varphi|_{F_i}$ is continuous (a continuous injection on a compactum). Put $g=\psi_i\circ h$ on
$X_i$. On $X_1\cap X_2$, $h(x)\in\rfr G_1=\varphi(F_1\cap F_2)$, so the $F_1$-point and the $F_2$-point over $h(x)$ coincide; likewise on $X_2\cap X_3$; and $X_1\cap
X_3=\emptyset$. So $g$ is well defined and continuous, and $\varphi\circ g=h$.
\end{proof}

\formal{\lean{HPC.HO16.ho16_theorem8}, with \lean{HPC.HO16.continuousOn_invFunOn_fst} for $\psi_i$.}

\section{Proposition 7: a component that covers}

\begin{proposition}[Hoehn--Oversteegen, Proposition~7]\label{prop:ho7}
Let $(F_1,F_2,F_3)$ be a simple fold on a connected $G$, with $\rfr G_1\neq\emptyset\neq\rfr G_3$. Then some component $C$ of $F$ has $\varphi(C)$ meeting both $\rfr G_1$ and
$\rfr G_3$. For every such component, $\varphi(C)=G$, and $(F_1\cap C,F_2\cap C,F_3\cap C)$ is a simple fold on $G$ along $(\varphi(F_1\cap C),\varphi(F_2\cap C),\varphi(F_3\cap C))$.
\end{proposition}

The new triple may be smaller than the old one: $\varphi(F_i\cap C)$ can be a proper subset of $G_i$.

\begin{proof}
For a component $K$ of $G_2$ write $K_i$ for its copy $\psi_i(K)$ in $F_i$; these are connected.

\emph{Existence.} Some component $K$ of $G_2$ meets both $\cl{G_1\setminus G_2}$ and $\cl{G_3\setminus G_2}$: otherwise, grouping the finitely many components of $G_2$
according to which closure they meet splits $G$ into two disjoint nonempty closed sets. Then $K$ meets $\rfr G_3$ and $\rfr G_1$ (Lemma~\ref{lem:ho6}(2)), and the
component of $F$ containing $K_2$ will do.

\emph{Lemma (L).} Let $C$ be a component of $F$ with $\varphi(C)$ meeting $\rfr G_1$ and $\rfr G_3$, and $K$ a component of $G_2$. (a) If $K\cap\rfr G_1=\emptyset$, then $K$ is a
component of $G_3$, and $K_2\cup K_3$ is clopen in $F$ and disjoint from $C$; symmetrically if $K\cap\rfr G_3=\emptyset$. (b) Hence $K_2\subseteq C$ implies $K_1\cup K_3\subseteq C$:
then $K$ meets both frontiers, $K_1$ meets $K_2$ over $\rfr G_1$ and $K_3$ meets $K_2$ over $\rfr G_3$. (c) $K$ cannot miss both frontiers, since then it would be clopen in
$G$, so $K=G$ and $\rfr G_1=\emptyset$.

\emph{$\varphi(C)$ is open in $G$, hence equal to $G$} ($\varphi(C)$ is closed and $G$ connected). At a point $p\in C$ outside $\rfr F_2$ use Lemma~\ref{lem:ho6}(7) and the
openness of components of a PL graph. At $p\in F_1\cap F_2$, let $K$ be the component of $G_2$ containing $\varphi(p)$. If $K$ missed $\rfr G_3$, then $K$ would be a
component of $G_1$, $K_1\cup K_2$ would be clopen in $F$ and contain $p$, so $C\subseteq K_1\cup K_2$ and $\varphi(C)\subseteq K$ would miss $\rfr G_3$. So by (L)(b) all three copies
of $K$ lie in $C$. In particular the $F_3$-point $q$ over $\varphi(p)$ lies in $C$, and $q\in F_3\setminus F_2$ because $\varphi(p)\in\rfr G_1$ is not in $\rfr G_3$. As $C$ is
open in $F$, Lemma~\ref{lem:ho6}(7) at $q$ shows that $\varphi(C)$ is a neighbourhood of $\varphi(p)$ in $G$. The case $p\in F_2\cap F_3$ is symmetric.

\emph{The restricted fold.} Put $G_i'=\varphi(F_i\cap C)$. Then $G_1'\cup G_3'=G$; $G_2'=G_1'\cap G_3'$ by (L)(b),(c); $G_1'$ and $G_3'$ are nonempty; the PL clauses
follow from Proposition~\ref{prop:pl-library}. For (F3) use that the component of $G_2$ through a point is locally constant near points of $G_2'$; and for (F5):
a point of $G_3'\setminus G_2'$ lying in $G_2$ lies in a component $K$ which by (L)(a) is a component of $G_1$ disjoint from $G_1'$, so near $G_1'$ the set $G_3'\setminus G_2'$ lies in
$G_3\setminus G_2$.
\end{proof}

\formal{\lean{HPC.HO16.ho16_prop7}, with \lean{HPC.HO16.exists_comp_meets_both}, \lean{HPC.HO16.image_comp_eq}, \lean{HPC.HO16.lift_subset_of_meets_both}.}

Hoehn and Oversteegen prove openness of $\varphi(C)$ at a glue point by boundary bumping (``let $C'$ be the closure of the component of $C\setminus\varphi^{-1}(\rfr G_3)$
containing $p$''), and they need, but do not argue, that $\varphi(C')\subseteq G_2$. The argument above avoids boundary bumping. Their sentence ``it is straightforward
to check \dots\ that $F'$ is a simple fold'' rests on Lemma~(L), which they do not state.

\begin{example}
Two test cases on intervals. (i) $G=[0,4]$, $G_1=[0,3]$, $G_3=[\frac12,1]\cup[2,4]$, $G_2=[\frac12,1]\cup[2,3]$: the copies of $[\frac12,1]$ in $F_2$ and $F_3$ form a
circle component of $F$; the other component restricts to the triple $([0,3],[2,3],[2,4])$. (ii) $G=[0,3]$, $G_1=[0,1]\cup[2,3]$, $G_3=[\frac12,3]$,
$G_2=[\frac12,1]\cup[2,3]$: the glue point over $2\in\rfr G_1$ is not in the covering component $C$, and the restriction is $([0,1],[\frac12,1],[\frac12,3])$, with
$\rfr G_1'=\set1$.
\end{example}

\section*{Exercises}

\begin{exercise}
Show that Lemma~\ref{lem:km} fails for $X=\I$, $A=\set0$, $B=\set1$, $U=[0,\frac13)$, $V=(\frac23,1]$.
\end{exercise}

\begin{exercise}
Show that Lemma~\ref{lem:km} characterizes hereditary indecomposability among compacta: if every triple $(X,A,B)$ is crooked between every $U$ and $V$, then
$X$ is hereditarily indecomposable. (Krasinkiewicz--Minc.)
\end{exercise}

\begin{exercise}
Draw the simple fold of Theorem~\ref{thm:fold-exists} for $G=[0,3]$, $G_1=[0,2]$, $G_3=[1,3]$. Which component of $F$ maps onto $G$?
\end{exercise}

\begin{exercise}
Why is the midpoint subdivision needed in the proof of Theorem~\ref{thm:fold-exists}? Give a decomposition without it for which $F_1\cap F_2$ is too large.
\end{exercise}

\begin{exercise}
Show that the map $g$ of Theorem~\ref{thm:ho8} cannot in general be chosen with $\varphi\circ g=f$, even for $X=\I$ and $f$ the identity of $G=\I$.
\end{exercise}

%======================================================================
\chapter{Stairwells, and Theorem 20 from Theorem 15}\label{ch:stairwells}
%======================================================================

This chapter reduces Theorem~\ref{thm:ho20} to a statement about separators alone (Hoehn--Oversteegen's Theorem~15, proved in
Chapter~\ref{ch:ladders}). A separator $M$ of the bottom from the top of $G^\star$ is approximated from inside a neighbourhood $U$ by a
\emph{stairwell}: finitely many sheets $S_1,\dots,S_k$, each the graph of a function over a subgraph of $G$, joined end to end, with $k$ odd. Think of a
staircase that climbs from left to right over $G$, turns back, climbs over a smaller part, turns again, and so on. If $k=1$ the single sheet is a section
of $G^\star\to G$, and $f$ lifts into it exactly. If $k>1$, one simple fold of $G$ unwinds one turn of the staircase: on the fold, the preimage of the
stairwell contains a stairwell whose ``pit'' has moved up one level. Theorem~\ref{thm:ho8} lifts $f$ to the fold, and after finitely many folds the height
is $1$.

Throughout, $G$ is a PL graph in a normed space $E$, $\pi\colon E\times\R\to E$ the projection, $G^\star=G\times\I$; relative frontiers $\rfr$ are taken in $G$. For a simple fold
$F$ on $G$ with projection $\varphi$ we write $F^\star=F\times\I\subseteq(E\times\R)\times\R$, $\varphi^\star(q,t)=(\varphi(q),t)$, and $\varphi^{\star-1}(S)$ for the preimage of $S\subseteq G^\star$ in
$F^\star$. Hoehn and Oversteegen's ``regular subsets'' of $G$ are read as PL subgraphs throughout.

\section{Sides}

\begin{definition}[Hoehn--Oversteegen, Definition~9]\label{def:sides}
Let $A\subseteq G$ be a PL graph and $B\subseteq\rfr A$.
\begin{enumerate}[label=(\alph*)]
\item $A$ has \emph{consistent complement relative to $B$} if for every component $C$ of $G\setminus A$ either $\rfr C\subseteq B$ or $\rfr C\cap B=\emptyset$.
\item The \emph{$A$ side of $B$}, $\side_B(A)$, is the closure of the union of the components of $G\setminus B$ that meet $A$.
\end{enumerate}
\end{definition}

\formal{\lean{HPC.HO16.HasConsCompl}, \lean{HPC.HO16.sideOf}.}

Everything below uses that $G\setminus K$ has finitely many components for a PL graph $K\subseteq G$ or a finite $K$, each open in $G$ with PL closure and
frontier inside $\rfr K$ (Proposition~\ref{prop:pl-library}).

\formal{\lean{HPC.HO16.IsPLGraph.finite_components_diff}, \lean{HPC.HO16.IsPLGraph.comp_diff_mem_nhdsWithin},
\lean{HPC.HO16.IsPLGraph.closure_comp_diff}, \lean{HPC.HO16.IsPLGraph.relFrontier_comp_diff}.}

\begin{proposition}[the observations after Definition~9]\label{prop:sides-char}
Let $G$ be connected, $A\subseteq G$ a PL graph and $B\subseteq\rfr A$. In (b)--(d) assume that $A$ and $B$ are nonempty and that $A$ has consistent complement relative to $B$.
\begin{enumerate}[label=(\alph*)]
\item If $B=\emptyset$ then $\side_B(A)$ is the union of the components of $G$ meeting $A$; in particular $\side_B(\emptyset)=\emptyset$.
\item For every neighbourhood $V$ of $B$, $\side_B(A)=\side_B(A\cap V)$.
\item $\side_B(A)$ is a PL graph containing $A$, $\rfr\side_B(A)=B$, and $\side_B(A)\cap V=A\cap V$ for some neighbourhood $V$ of $B$.
\item $\side_B(A)$ is the only closed $D\subseteq G$ with $\rfr D=B$ and $D\cap V=A\cap V$ for some neighbourhood $V$ of $B$.
\end{enumerate}
\end{proposition}

\begin{proof}
(b) Let $K$ be a component of $G\setminus B$ meeting $A$. As $G$ is connected and $B\neq\emptyset$, $\cl K$ meets $B$ at some $b$. A component $C$ of $G\setminus A$ inside $K$
has $\rfr C\cap B=\emptyset$: otherwise $\rfr C\subseteq B$, so $C$ is clopen in $K$ and $C=K$, but $K$ meets $A$. There are finitely many such $C$, none with $b$ in its closure;
so the points of $K$ near $b$ are in $A$, and $K$ meets $A\cap V$.

(c) $\rfr\side_B(A)\subseteq B$ since the union of components is open. Conversely $b\in B$ is a limit of points of $A\setminus B$ ($B$ is finite and $A$ has no isolated
points), so $b\in\side_B(A)$; and $b$ is in the closure of a component $C$ of $G\setminus A$, which then has $\rfr C\subseteq B$, so $C$ is a component of $G\setminus B$
missing $A$, and $C\cap\side_B(A)=\emptyset$. For $V$ take a neighbourhood of $B$ missing the closures of the finitely many components $C$ of $G\setminus A$ with
$\rfr C\cap B=\emptyset$.

(d) If $\rfr D=B$, then $D\setminus B$ is clopen in $G\setminus B$, hence a union of components; each of them meets $V\cap V'$ near a point of its closure in $B$, so it
meets $\side_B(A)\setminus B$; and conversely.
\end{proof}

\formal{\lean{HPC.HO16.ho16_obs9_emptyB}, \lean{HPC.HO16.ho16_obs9_nhds}, \lean{HPC.HO16.ho16_obs9_char}, \lean{HPC.HO16.ho16_obs9_unique}; the
common core is \lean{HPC.HO16.IsPLGraph.sideOf_char_aux}.}

\begin{proposition}[Hoehn--Oversteegen, Proposition~10]\label{prop:ho10}
Let $G$ be connected, $A,A'\subseteq G$ nonempty PL graphs, $B\subseteq\rfr A\cap\rfr A'$, both of consistent complement relative to $B$, and $A\cap V=A'\cap V$ for a
neighbourhood $V$ of $B$. Then $\side_B(A)=\side_B(A')$; and every component $C$ of $G\setminus A$ with $\cl C\cap B\neq\emptyset$ is a component of $G\setminus A'$.
\end{proposition}

\begin{proof}
The first part is Proposition~\ref{prop:sides-char}(d). For the second, $\rfr C\subseteq B$ by consistency. Shrink $V$ so that it misses the closures of the finitely
many components of $G\setminus A'$ whose closures miss $B$. Let $C'$ be a component of $G\setminus A'$ meeting $C\cap V$. Then $C'\subseteq C$ (since $\rfr C\subseteq B$ and $C'\cap
B=\emptyset$), and $\rfr C'$ meets $B$, so $\rfr C'\subseteq B$. If $C'\neq C$, a point of $\rfr C'\cap C$ would lie in $B\cap C\subseteq A\cap C=\emptyset$.
\end{proof}

\formal{\lean{HPC.HO16.ho16_prop10}.}

In Hoehn and Oversteegen's proof of the second part, consistency of $A'$ is applied to a component $C'$ that merely meets $C\cap V$; that its frontier meets
$B$ needs the shrinking of $V$ above. The second part is used again in Proposition~\ref{prop:ho19}.

\begin{proposition}[Hoehn--Oversteegen, Proposition~11]\label{prop:ho11}
Let $G$ be connected, $A\subseteq G$ a nonempty PL graph, and $\rfr A=B_1\cup B_2$ with $B_1\cap B_2=\emptyset$, $A$ of consistent complement relative to $B_1$ and $B_2$.
Then $(\side_{B_1}(A),A,\side_{B_2}(A))$ is a fold triple on $G$, with $\rfr\side_{B_1}(A)=B_1$ and $\rfr\side_{B_2}(A)=B_2$. Consistency relative to $B_1$ implies
consistency relative to $B_2$.
\end{proposition}

\begin{proof}[Proof sketch]
Every component of $G\setminus A$ has frontier inside $B_1$ or inside $B_2$; $\side_{B_1}(A)$ is $A$ together with the closures of the components whose frontier lies in
$B_2$, and symmetrically. This gives $\side_{B_1}(A)\cup\side_{B_2}(A)=G$ and $\side_{B_1}(A)\cap\side_{B_2}(A)=A$; the closures of the two differences are separated
because their attaching points lie in the disjoint sets $B_2$ and $B_1$. The frontiers come from Proposition~\ref{prop:sides-char}(c).
\end{proof}

\formal{\lean{HPC.HO16.ho16_prop11}, \lean{HPC.HO16.ho16_prop11_frontier}, \lean{HPC.HO16.ho16_prop11_remark}.}

\section{Straight sets and stairwell structures}

\begin{definition}[Hoehn--Oversteegen, Definition~12]\label{def:straight}
\begin{enumerate}[label=(\alph*)]
\item A point $x$ of $G$ is \emph{ordinary} if some neighbourhood of $x$ in $G$ is the union of two segments $[a,x]$, $[x,b]$ meeting only in $x$. (So $x$ is neither an end point
nor a branch point of $G$.)
\item A list $\langle B_1,\dots,B_n\rangle$ of finite subsets of $G$ is \emph{generic} if all their points are ordinary and the $B_i$ are pairwise disjoint.
\item $S\subseteq G^\star$ is \emph{straight} if $S$ is closed, $\pi$ is injective on $S$, and $\pi(S)$ is a PL graph. Its \emph{end set} is $\mathcal E(S)=S\cap\pi^{-1}(\rfr\pi(S))$.
\end{enumerate}
\end{definition}

\formal{\lean{HPC.HO16.IsOrdinaryPt}, \lean{HPC.HO16.IsGeneric}, \lean{HPC.HO16.IsStraight}, \lean{HPC.HO16.endSet}.}

A straight set need not be a PL graph: it is the graph of a continuous function over the PL graph $\pi(S)$. Ordinariness does not refer to a decomposition:
for any decomposition, $x$ is ordinary iff it is interior to an edge or the end of exactly two edges; all but finitely many points are ordinary. $\pi$ maps
$S\setminus\mathcal E(S)$ openly into $G$.

\formal{\lean{HPC.HO16.isOrdinaryPt_iff_order}, \lean{HPC.HO16.IsPLGraph.finite_not_ordinary}, \lean{HPC.HO16.ho16_obs12}.}

\begin{definition}[Hoehn--Oversteegen, Definition~13]\label{def:stairwell}
A \emph{stairwell structure of height $k\ge1$} for $S\subseteq G^\star$ is a list $\langle S_1,\dots,S_k\rangle$ with finite sets $\alpha_1,\dots,\alpha_{k+1}\subseteq G^\star$ (write
$\beta_i=\alpha_{i+1}$) such that:
\begin{enumerate}[label=(S\arabic*)]
\item $S_1,\dots,S_k$ are nonempty straight sets with union $S$;
\item $\mathcal E(S_i)=\alpha_i\cup\beta_i$ with $\alpha_i\cap\beta_i=\emptyset$, and $\alpha_1=\emptyset=\beta_k$;
\item for $i<k$ there is a neighbourhood $V$ of $\pi(\beta_i)$ with $\pi(S_i)\cap V=\pi(S_{i+1})\cap V$;
\item each $\pi(S_i)$ has consistent complement relative to $\pi(\alpha_i)$ and to $\pi(\beta_i)$;
\item $\langle\pi(\alpha_2),\dots,\pi(\alpha_k)\rangle$ is generic in $G$.
\end{enumerate}
\end{definition}

\formal{\lean{HPC.HO16.IsStairwellStr} (with the sets $\alpha_i$ as one sequence, so that $\beta_i=\alpha_{i+1}$ holds by construction) and its existential form
\lean{HPC.HO16.HasStairwell}.}

So consecutive sheets share their end points and agree near them over $G$; sheet $i$ meets sheet $i+1$ only at $\beta_i$. The sets $\alpha_i$ may be empty for
$2\le i\le k$.

\begin{proposition}[remarks on pp.~193--194]\label{prop:stair-remarks}
Let $\langle S_1,\dots,S_k\rangle$ be a stairwell structure for $S$.
\begin{enumerate}[label=(\alph*)]
\item If $C$ is a component of $G$ meeting every $\pi(S_i)$, intersecting all sets with $C^\star$ gives a stairwell structure for $S\cap C^\star$ relative to the graph $C$.
\item If $G$ is connected and $k=1$, then $\pi$ maps $S$ homeomorphically onto $G$.
\item \emph{(Claim 14.1.)} If $G$ is connected and $k$ odd, every $x\in G$ lies in an odd number of the sets $\pi(S_i)$; in particular $\pi(S)=G$.
\item Without connectedness, (b) and (c) fail: for $G=[0,1]\cup[2,3]$, $S=[0,1]\times\set{\frac12}$ is a stairwell structure of height $1$ (its end set is empty, as $\rfr[0,1]=\emptyset$ in
$G$) and $\pi(S)\neq G$.
\end{enumerate}
\end{proposition}

\begin{proof}
(a) $C$ is clopen in $G$, so frontiers, components of complements and ordinary points relative to $C$ are those relative to $G$. (b) $\mathcal E(S_1)=\emptyset$, so
$\pi(S_1)$ is a nonempty clopen subset of $G$, so it is $G$; and $\pi|_{S_1}$ is a continuous injection of a compactum. (c) By induction along the stairwell, using
Proposition~\ref{prop:ho10} with (S3) and Proposition~\ref{prop:ho11} with $B_1=\pi(\alpha_i)$, $B_2=\pi(\beta_i)$: crossing $\pi(\beta_i)$ changes the count by $0$ or $2$.
\end{proof}

\formal{\lean{HPC.HO16.ho16_obs13_comp}, \lean{HPC.HO16.ho16_height_one}, \lean{HPC.HO16.ho16_claim14_1}, \lean{HPC.HO16.ho16_odd_image},
\lean{HPC.HO16.ho16_p193_remark_false}.}

Hoehn and Oversteegen allow disconnected graphs (``Throughout the remainder of this paper, $G$ will denote a (not necessarily connected) graph''), and state the
remarks (b), (c) on pp.~193--194 without connectedness; (d) shows that it is needed there. Their Proposition~14, to which the remarks refer, does assume it.

\begin{definition}[Hoehn--Oversteegen, Definition~16]\label{def:broken}
A \emph{broken stairwell structure of height $k$ with a pit at level $i_0$} ($1\le i_0\le k$) for $S\subseteq G^\star$ is a list $\langle S_1,\dots,S_k;P_1,P_2\rangle$, finite sets
$\alpha_1,\dots,\alpha_{k+1}$ (with $\beta_i=\alpha_{i+1}$) and a finite set $\gamma$ such that:
\begin{enumerate}[label=(S\arabic*$'$)]
\item $S_1,\dots,S_k,P_1,P_2$ are nonempty straight sets with union $S$;
\item (S2) holds except that $\mathcal E(S_{i_0})=\alpha_{i_0}\cup\beta_{i_0}\cup\gamma$, a disjoint union; and $\mathcal E(P_2)=\mathcal E(P_1)\cup\gamma$ with $\mathcal E(P_1)\cap\gamma=\emptyset$;
\item (S3) holds, and $\pi(P_1)$ and $\pi(P_2)$ agree near $\pi(\mathcal E(P_1))$, $\pi(P_2)$ and $\pi(S_{i_0})$ agree near $\pi(\gamma)$;
\item (S4) holds, $\pi(S_{i_0})$ has consistent complement relative to $\pi(\gamma)$, and $\pi(P_2)$ relative to $\pi(\mathcal E(P_1))$ and to $\pi(\gamma)$;
\item $\langle\pi(\alpha_2),\dots,\pi(\alpha_k),\pi(\mathcal E(P_1)),\pi(\gamma)\rangle$ is generic;
\item $\pi(\alpha_{i_0})\cap\pi(P_1\cup P_2)=\emptyset$.
\end{enumerate}
\end{definition}

\formal{\lean{HPC.HO16.IsBrokenStairwellStr}; the restriction to a component is \lean{HPC.HO16.ho16_obs16_comp}.}

The pit is the pair $P_1,P_2$: a short staircase hanging from sheet $i_0$ at $\gamma$, going back over $\pi(P_2)$ and forward again over $\pi(P_1)$, ending at $\mathcal E(P_1)$.

\begin{proposition}[Hoehn--Oversteegen, Proposition~17]\label{prop:ho17}
If $S$ has a stairwell structure $\langle S_1,\dots,S_k\rangle$ of height $k\ge3$, then $\langle S_3,\dots,S_k;S_1,S_2\rangle$ is a broken stairwell structure for $S$ of height $k-2$ with a pit at level~$1$,
with $\alpha_1'=\emptyset$, $\alpha_i'=\alpha_{i+2}$ for $i\ge2$, and $\gamma'=\alpha_3$.
\end{proposition}

\begin{proof}
$\mathcal E(P_2)=\mathcal E(S_2)=\alpha_2\cup\alpha_3=\mathcal E(P_1)\cup\gamma'$, since $\alpha_1=\emptyset$ gives $\mathcal E(P_1)=\alpha_2$; and $\alpha_2\cap\alpha_3=\emptyset$. The neighbourhoods of (S3$'$) are those of
(S3) at $i=1,2$; (S4$'$) is (S4) at $i=2,3$; the list of (S5$'$) is a permutation of that of (S5); and (S6$'$) holds as $\alpha_1'=\emptyset$.
\end{proof}

\formal{\lean{HPC.HO16.ho16_prop17}, \lean{HPC.HO16.ho16_prop17_exists}. The hypothesis $k\ge3$ is not printed; for $k\le2$ the conclusion has height $\le0$.}

\begin{proposition}[remarks on p.~201]\label{prop:rem17}
Let $\langle S_1,\dots,S_k;P_1,P_2\rangle$ be a broken stairwell structure for $S$. If $\gamma=\emptyset$, then $\langle S_1,\dots,S_k\rangle$ is a stairwell structure for $S_1\cup\dots\cup S_k$. If
$\mathcal E(P_1)=\emptyset$ and $G$ is connected, then $\langle P_1\rangle$ is a stairwell structure of height~$1$ for $P_1$.
\end{proposition}

\formal{\lean{HPC.HO16.ho16_rem17_gamma}, \lean{HPC.HO16.ho16_rem17_P1}.}

\section{Pulling straight sets back to a fold}

\begin{lemma}[Hoehn--Oversteegen, Lemma~18]\label{lem:ho18}
Let $F=F_1\cup F_2\cup F_3$ be a simple fold on $G$, and $S\subseteq G^\star$ straight. Suppose that $\rfr\pi(S)\cap\rfr\varphi(F_2)=\emptyset$, or that $\varphi(F_2)\cap V\subseteq\pi(S)\cap V$ for a
neighbourhood $V$ of $\rfr\pi(S)\cap\rfr\varphi(F_2)$. Then, with relative frontiers of subsets of $F$ taken in $F$:
\begin{enumerate}[label=(\arabic*)]
\item $S'=\varphi^{\star-1}(S)$ is straight in $F^\star$, and $\mathcal E(S')=\varphi^{\star-1}(\mathcal E(S))\setminus(\rfr F_2)^\star$;
\item $S''=S'\cap(F_1\cup F_2)^\star$ is straight in $F^\star$, and $\mathcal E(S'')=\bigl([\varphi^{\star-1}(\mathcal E(S))\cap(F_1\cup F_2)^\star]\setminus(\rfr F_1)^\star\bigr)\cup(S''\cap(\rfr F_3)^\star)$.
\end{enumerate}
\end{lemma}

\begin{proof}
\emph{Claim: if $x\in S'$ lies over $\rfr F_2$, then $\varphi(W)\subseteq\pi(S)$ for a neighbourhood $W$ of $\pi(x)$ in $F$.} Say $\pi(x)\in\rfr F_1=F_1\cap F_2$, so $y=\varphi(\pi(x))\in\rfr G_1\subseteq\rfr G_2$
and $y\in\pi(S)$. If $y$ is interior to $\pi(S)$ take $V\subseteq\pi(S)$; otherwise the hypothesis gives $V$. Shrink $V$ to miss $\cl{G_1\setminus G_2}$, so that $G_1\cap V=G_2\cap V$.
$F_3$ is closed and misses $\pi(x)$, so $W=\varphi^{-1}(V)\cap(F_1\cup F_2)$ is a neighbourhood of $\pi(x)$ in $F$, and $\varphi(W)=G_2\cap V\subseteq\pi(S)$.

(1) $\pi(S')=\varphi^{-1}(\pi(S))$, and $\pi$ is injective on $S'$ because it is on $S$. Off $(\rfr F_2)^\star$, $\varphi$ is a local homeomorphism (Lemma~\ref{lem:ho6}(7)), so
frontier points of $\pi(S')$ correspond to frontier points of $\pi(S)$; over $\rfr F_2$ the claim puts $\pi(x)$ in the interior of $\pi(S')$. $\pi(S')$ is a PL graph plus finitely many
points, and has no isolated points, so it is a PL graph.

(2) The same on $F_1\cup F_2$; in addition a point over $\rfr F_3=F_2\cap F_3$ is an end point, since every neighbourhood of it meets $F_3\setminus F_2$, which misses $\pi(S'')$.
\end{proof}

\formal{\lean{HPC.HO16.ho16_lemma18}, with the preimage \lean{HPC.HO16.foldPre} and $\varphi^\star$ as \lean{HPC.HO16.phiStar}; PL-graph facts for preimages:
\lean{HPC.HO16.IsPLGraph.inter_preimage_fst}, \lean{HPC.HO16.isPLGraph_union_finite}.}

Hoehn and Oversteegen check that the components of $\pi(S')$ are nondegenerate but not that there are finitely many. On p.~202 they print
``$\mathcal E(S'')\subset(F_1\cap F_2)^\star$ since $S''\subset(F_1\cap F_2)^\star$'' for $(F_1\cup F_2)^\star$.

\section{Proposition 19: unwinding one turn}

\begin{proposition}[Hoehn--Oversteegen, Proposition~19]\label{prop:ho19}
Let $G$ be connected and $S\subseteq G^\star$ have a broken stairwell structure $\langle S_1,\dots,S_k;P_1,P_2\rangle$ with a pit at level $i_0$, in which $\gamma\neq\emptyset$ and $\mathcal E(P_1)\neq\emptyset$.
Then there is a simple fold $F$ on $G$ with $F$ connected such that $\varphi^{\star-1}(S)$ contains a set $S'$ with a broken stairwell structure of height $k$ with a pit at level $i_0+1$ if
$i_0<k$, and a stairwell structure of height $k$ if $i_0=k$.
\end{proposition}

\formal{\lean{HPC.HO16.ho16_prop19}, whose statement is the proposition \lean{HPC.HO16.HOProp19}; for the full (possibly disconnected) fold,
\lean{HPC.HO16.ho16_prop19_full}.}

\begin{proof}[Proof outline]
\emph{The fold.} Put $G_1=\pi(P_1)$, $G_2=\pi(P_2)$, $G_3=\side_{\pi(\gamma)}(\pi(S_{i_0}))$. By Proposition~\ref{prop:sides-char}(d) and (S3$'$), $\pi(P_1)=\side_{\pi(\mathcal E(P_1))}(\pi(P_2))$; by
Proposition~\ref{prop:ho10}, $G_3=\side_{\pi(\gamma)}(\pi(P_2))$. So Proposition~\ref{prop:ho11} with $A=\pi(P_2)$, $B_1=\pi(\mathcal E(P_1))$, $B_2=\pi(\gamma)$ shows that $(G_1,G_2,G_3)$ is a fold triple,
and
\begin{equation}\label{eq:19-frontiers}
\rfr G_1=\pi(\mathcal E(P_1)),\qquad \rfr G_3=\pi(\gamma),\qquad \rfr G_2=\pi(\mathcal E(P_1))\cup\pi(\gamma).
\end{equation}
Theorem~\ref{thm:fold-exists} gives a simple fold $F$.

\emph{The sets.} For $i\notin\set{i_0,i_0+1}$: $S_i'=\varphi^{\star-1}(S_i)$, $\alpha_i'=\varphi^{\star-1}(\alpha_i)$. At level $i_0$:
\[
S_{i_0}'=[\varphi^{\star-1}(P_1)\cap F_1^\star]\cup[\varphi^{\star-1}(P_2)\cap F_2^\star]\cup[\varphi^{\star-1}(S_{i_0})\cap F_3^\star],\qquad \alpha_{i_0}'=\varphi^{\star-1}(\alpha_{i_0}),
\]
and $\alpha_{i_0+1}'=\varphi^{\star-1}(\alpha_{i_0+1})\cap F_3^\star$. If $i_0<k$: $S_{i_0+1}'=\varphi^{\star-1}(S_{i_0+1})$, $\gamma'=\varphi^{\star-1}(\alpha_{i_0+1})\cap(F_1\cup F_2)^\star$,
$P_1'=\varphi^{\star-1}(P_2)\cap(F_1\cup F_2)^\star$ and $P_2'=\varphi^{\star-1}(S_{i_0})\cap(F_1\cup F_2)^\star$. In words: the new sheet $i_0$ runs along $P_1$ over $F_1$, back along $P_2$ over $F_2$, and
on along $S_{i_0}$ over $F_3$: the pit has been unwound. The part of sheet $i_0+1$ over $F_1\cup F_2$ becomes the new pit.

\emph{Checks.}
\begin{itemize}
\item For $i\neq i_0$, $\rfr\pi(S_i)=\pi(\alpha_i)\cup\pi(\beta_i)$ misses $\rfr G_2$ by (S5$'$), so Lemma~\ref{lem:ho18}(1) makes $S_i'$ straight with $\mathcal E(S_i')=\varphi^{\star-1}(\mathcal E(S_i))$.
\item At level $i_0$: over $\rfr F_1$ the pieces from $P_1$ and $P_2$ coincide, because $\mathcal E(P_1)\subseteq\mathcal E(P_2)\subseteq P_2$ and $\pi$ is injective on $P_2$; likewise over $\rfr F_3$
with $\gamma\subseteq\mathcal E(P_2)\cap\mathcal E(S_{i_0})$. So $S_{i_0}'$ is straight; $F_1\cup F_2$ is interior to $\pi(S_{i_0}')$, and the end points of $S_{i_0}'$ lie over the interior of $F_3$, as
$\pi(S_{i_0})$ agrees with $G_3$ near $\pi(\gamma)$ (Proposition~\ref{prop:sides-char}(c)).
\item $P_1'$ and $P_2'$: Lemma~\ref{lem:ho18}(2) with $S=P_2$ and $S=S_{i_0}$, and $\mathcal E(P_1')=\varphi^{\star-1}(\gamma)$, $\mathcal E(P_2')=\mathcal E(P_1')\cup\gamma'$, using (S6$'$).
\item Finiteness in (S2): a point of $G$ has at most three preimages in $F$. The neighbourhoods of (S3) pull back, through local homeomorphisms off $\rfr F_2$ and, at a point
$z\in\rfr F_3$, through the two spokes of $F$ at $z$ (one in $F_1\cup F_2$, one in $F_3$), which exist since $\varphi(z)$ is ordinary.
\item (S4): a component $C$ of $F\setminus\pi(S_i')$ maps into a component $\tilde C$ of $G\setminus\pi(S_i)$, with $\varphi(\rfr C)\subseteq\rfr\tilde C$, and consistency for $S_i$ applies; at
level $i_0+1$ one uses Proposition~\ref{prop:ho10} (second part) and the separation of Lemma~\ref{lem:ho6}(6).
\item (S5): if every point of $\rfr G_2$ is ordinary in $G$, a point $z\in F$ is ordinary in $F$ iff $\varphi(z)$ is ordinary in $G$; and preimages of disjoint sets are disjoint.
\item (S6$'$): $\alpha_{i_0+1}'$ lies over the interior of $F_3$.
\end{itemize}

\emph{Connectedness.} By Proposition~\ref{prop:ho7} there is a component $C$ of $F$ whose image meets $\rfr G_1$ and $\rfr G_3$; then $\varphi(C)=G$ and the pieces $F_i\cap C$ form a
connected simple fold. Every set of the new structure meets $C^\star$: for $i\neq i_0$ because $\varphi(C)=G$; $S_{i_0}'$, $P_1'$, $P_2'$ because each contains the points over
$\gamma$ in $(\rfr F_3)^\star$, and $C$ meets $\rfr F_3$ by Lemma~\ref{lem:ho6}(6) ($C$ meets both $F_1\setminus F_2$ and $F_3\setminus\rfr F_3$, since $G\setminus G_3$ and $G\setminus G_1$ are nonempty).
Restrict the structure to $C$ (Proposition~\ref{prop:stair-remarks}(a) and its broken analogue).
\end{proof}

\formal{The checks are spread over \lean{HPC/HO16/Prop19A.lean} to \lean{HPC/HO16/Prop19D.lean}; the fold is \lean{HPC.HO16.ho16_prop19_fold}; the sheets are
\lean{HPC.HO16.p19T}, with a companion definition for the end sets in \lean{HPC/HO16/P19Defs.lean}; ordinary points of folds: \lean{HPC.HO16.IsSimpleFoldOf.isOrdinaryPt_iff}.}

\begin{remark}[A false sentence on p.~208]\label{rem:p208}
In the (S4$'$) check for $P_2'$, Hoehn and Oversteegen take a component $D\subseteq F_1\cup F_2$ of $F\setminus\pi(P_2')$, its image in a component $\tilde D$ of $G\setminus\pi(S_{i_0})$, and assert
that $\rfr\tilde D\cap\pi(\alpha_{i_0})=\emptyset$ by (S6$'$). This is false in a legitimate configuration: take $G=[0,10]$, $k=3$, $i_0=2$, $5<a<6$,
\begin{itemize}
\item $S_1=[4,10]\times\set{0.1}$, $S_2=[(4,0.1),(6,0.5)]$, $S_3=[0,10]\times\set{0.9}$; $\alpha_2=\set{(4,0.1)}$, $\gamma=\set{(6,0.5)}$, all other $\alpha_i$ empty;
\item $P_2=[1,2]\times\set{0.3}\cup[(a,0.3),(6,0.5)]$, and $P_1$ a path over $[1,2]$ from $(1,0.3)$ to $(2,0.3)$ through $(1.5,0.2)$, together with $[(a,0.3),(10,0.2)]$.
\end{itemize}
Then $x=1.5\in\pi(P_2)\setminus\pi(S_2)$ lies in the component $\tilde D=[0,4)$ of $G\setminus\pi(S_2)$, and $\rfr\tilde D=\set4=\pi(\alpha_2)$. In the fold, the copies of $[1,2]$ in $F_1$ and $F_2$ form a
circle component $D$ of $F$ over $\tilde D$. The conclusion of the step survives, by a case split: if $\rfr D=\emptyset$ there is nothing to show; otherwise a point $z\in\rfr D$ has
$\varphi(z)\in\rfr\tilde D\cap\pi(P_1\cup P_2)$, which misses $\pi(\alpha_{i_0})$ by (S6$'$), so consistency puts $\rfr\tilde D$ inside $\pi(\beta_{i_0})$ or inside $\pi(\gamma)$.
\end{remark}

\formal{\lean{HPC.HO16.ho16_p208_false}.}

On p.~205 Hoehn and Oversteegen also print ``$P_1'\cap(\rfr F_3)^\star=\varphi^{\star-1}(\gamma_{i_0})$''; the right-hand side also contains the points over the $F_1$-copies of $\pi(\gamma)$, outside
$(\rfr F_3)^\star$, so only ``$\subseteq$'' holds, which is all that is used.

\section{Theorem 20 from Theorem 15}

\begin{theorem}[Hoehn--Oversteegen, Theorem~15]\label{thm:ho15}
Let $G$ be a connected PL graph, $M\subseteq G\times(0,1)$ a set separating $G\times\set0$ from $G\times\set1$ in $G^\star$, and $U\supseteq M$ open. Then some $S\subseteq U\cap G\times(0,1)$ has a
stairwell structure of odd height.
\end{theorem}

\formal{The statement is the proposition \lean{HPC.HO16.HOThm15}; it is proved as \lean{HPC.HO16.ho16_theorem15} (Chapter~\ref{ch:ladders}).}

Hoehn and Oversteegen state Theorem~15 for any graph. Some nonemptiness is needed: for $G=\emptyset$ and $M=\emptyset$ the separation holds with empty open sets, and no stairwell
lives in $\emptyset$. Connectedness is all that the proof of Theorem~20 uses.

\begin{theorem}\label{thm:lift-stairwell}
Let $G$ be a connected PL graph, $S\subseteq G^\star$ with a stairwell structure of odd height $k$, $X$ a hereditarily indecomposable compact metric space and $f\colon X\to G$ continuous.
For every $\varepsilon>0$ there is a continuous $h\colon X\to S$ with $\nrm{f(x)-\pi h(x)}<\varepsilon$.
\end{theorem}

\begin{proof}
Prove simultaneously the same statement for broken structures of odd height (pit at any level), by strong induction on the measure $(k+1)^2$ for a stairwell of height $k$ and
$(k+1)^2+k+1-i_0$ for a broken structure with a pit at $i_0$, quantified over the normed space $E$ (the induction passes from $E$ to $E\times\R$).
\begin{itemize}
\item Height $1$: $\pi|_S$ is a homeomorphism onto $G$ (Proposition~\ref{prop:stair-remarks}(b)); take $h=(\pi|_S)^{-1}\circ f$.
\item Height $k\ge3$: Proposition~\ref{prop:ho17} gives a broken structure of height $k-2$ with a pit at level~$1$, of smaller measure.
\item Broken, with $\gamma=\emptyset$ or $\mathcal E(P_1)=\emptyset$: Proposition~\ref{prop:rem17} gives a stairwell structure of height $k$ or $1$ on a subset of $S$, of smaller measure.
\item Broken, otherwise: Proposition~\ref{prop:ho19} gives a connected fold $F$ and $S'\subseteq\varphi^{\star-1}(S)$ with a structure of smaller measure in $F^\star$. Theorem~\ref{thm:ho8} gives
$g\colon X\to F$ with $\nrm{f(x)-\varphi(g(x))}<\varepsilon/2$; the induction hypothesis for $F$, $S'$, $g$ and $\varepsilon/2$ gives $h'\colon X\to S'$; put $h=\varphi^\star\circ h'$, which maps into $S$.
Since $\varphi$ is a projection, $\nrm{f(x)-\pi h(x)}\le\nrm{f(x)-\varphi g(x)}+\nrm{g(x)-\pi h'(x)}<\varepsilon$.
\end{itemize}
Oddness is preserved at each step, which excludes height $2$, where Proposition~\ref{prop:ho17} does not apply.
\end{proof}

\formal{\lean{HPC.HO16.ho16_lift_stairwell}, with the induction in \lean{HPC.HO16.ho16_lift_aux}.}

Hoehn and Oversteegen compose the folds explicitly, $h=(\varphi_1)^\star\circ\dots\circ(\varphi_n)^\star\circ\theta\circ g_n$; the recursion does the same without indexing them.

\begin{theorem}\label{thm:20-from-15}
Theorem~\ref{thm:ho15} implies Theorem~\ref{thm:ho20}.
\end{theorem}

\begin{proof}
$G$ has finitely many components $C$, each clopen; $M\cap C^\star$ separates $C\times\set0$ from $C\times\set1$ in $C^\star$ (with the same open sets), and $f^{-1}(C)$ is a clopen, hereditarily
indecomposable compactum. On each piece, Theorem~\ref{thm:ho15} gives $S\subseteq U\cap C\times(0,1)$ with a stairwell structure of odd height and Theorem~\ref{thm:lift-stairwell} gives $h$; glue
the pieces.
\end{proof}

\formal{\lean{HPC.HO16.ho16_theorem20_of_thm15}; with Theorem~15 as the only hypothesis, the classification step is
\lean{HPC.HO16.homeomorph_pseudoArc_of_indec_homogeneous_of_thm15}.}

\section*{Exercises}

\begin{exercise}
Let $G=\I$, $S_1=[0,\frac23]\times\set{\frac14}$, $S_2$ the segment from $(\frac23,\frac14)$ to $(\frac13,\frac12)$ and $S_3$ the segment from $(\frac13,\frac12)$ to $(1,\frac34)$.
Show that $\langle S_1,S_2,S_3\rangle$, with $\alpha_2=\set{(\frac23,\frac14)}$ and $\alpha_3=\set{(\frac13,\frac12)}$, is a stairwell structure of height $3$, and check Claim~14.1.
\end{exercise}

\begin{exercise}
For the stairwell of the previous exercise, carry out Proposition~\ref{prop:ho17} and one application of Proposition~\ref{prop:ho19}. Draw the fold and the new structure.
\end{exercise}

\begin{exercise}
Show that for a stairwell structure of even height on a connected $G$, $\pi(S)$ need not be $G$.
\end{exercise}

\begin{exercise}
Check the five conditions (S1$'$)--(S6$'$) for the configuration of Remark~\ref{rem:p208}.
\end{exercise}

\begin{exercise}
Why must the measure in the proof of Theorem~\ref{thm:lift-stairwell} depend on $i_0$?
\end{exercise}

%======================================================================
\chapter{Theorem 15: ladders}\label{ch:ladders}
%======================================================================

It remains to prove Theorem~\ref{thm:ho15}: inside any neighbourhood $U$ of a separator $M$ of the bottom of $G^\star$ from the top there is a set with a
stairwell structure of odd height. Hoehn and Oversteegen first replace $M$ by a nicer separator $M'$ (their Claim~15.1, ``left to the reader''), make it irreducible,
and then build the stairwell over the arcs of a subdivision, with a zig-zag added wherever two curves of $M'$ are born or die. This chapter replaces Claim~15.1
and the irreducibility by an explicit combinatorial object, a \emph{ladder}, read off from the signs of one function on a grid. The odd counts come from
sign changes, and the parity argument of Hoehn and Oversteegen for (S4) becomes a statement about parities of list lengths.

\section{Ladders}

Fix a decomposition $\mathcal S$ of the PL graph $G$ into oriented edges $s=(s_1,s_2)$, and write $s(\mu)=s_1+\mu(s_2-s_1)$ for $\mu\in\I$.

\begin{definition}[Ladder]\label{def:ladder}
A \emph{ladder} on $G$ consists of a decomposition $\mathcal S$ of $G$, a list $y(v)=(y_1(v),\dots,y_{j(v)}(v))$ of heights in $(0,1)$ of \emph{odd} length $j(v)$ at each vertex $v$, and
for each edge $s$ an \emph{event}, which is either
\begin{itemize}
\item \emph{none}: then $j(s_1)=j(s_2)$; or
\item a \emph{wedge} $(m,\lambda,h)$ at $s_2$: then $j(s_2)=j(s_1)+2$, $1\le m\le j(s_1)+1$, $\lambda,h\in(0,1)$. The wedge occupies places $m$ and $m+1$ at $s_2$, and its tip is
$(s(\lambda),h)$.
\end{itemize}
The \emph{height} of the ladder is the largest $j(v)$.
\end{definition}

\formal{\lean{HPC.HO16.IsLadder}, \lean{HPC.HO16.ladderY}, \lean{HPC.HO16.ladderHeight} (in \lean{HPC/HO16/LadderDefs.lean}).}

\emph{The sheets.} Let $s$ be an edge and $j=j(s_1)$.
\begin{itemize}
\item No event: for $1\le i\le j$, sheet $i$ contains the segment from $(s_1,y_i(s_1))$ to $(s_2,y_i(s_2))$.
\item A wedge $(m,\lambda,h)$: for $i<m$ as before. For $m\le i\le j$, let $L_i$ be the line from $(s_1,y_i(s_1))$ to $(s_2,y_{i+2}(s_2))$, parametrized by $\mu\in\I$, and put
\[
l_i=\lambda\,\frac{2(j-i)+1}{2(j-m)+3},\qquad r_i=\lambda\,\frac{2(j-i)+2}{2(j-m)+3},
\]
so that $0<l_j<r_j<l_{j-1}<\dots<l_m<r_m<\lambda$. The \emph{zig-zag} of curve $i$ consists of $C_i^1=L_i[0,r_i]$ in sheet $i$, $C_i^2=L_i[l_i,r_i]$ in sheet $i+1$, and $C_i^3=L_i[l_i,1]$ in
sheet $i+2$. Finally the two halves of the wedge, the segments from the tip to $(s_2,y_m(s_2))$ and to $(s_2,y_{m+1}(s_2))$, go into sheets $m$ and $m+1$.
\end{itemize}
Sheet $i$ is the union of its pieces over all edges.

\formal{\lean{HPC.HO16.ladderPiece}, \lean{HPC.HO16.ladderSheet}; the zig-zag parameters are \lean{HPC.HO16.zzL}, \lean{HPC.HO16.zzR} and the lines
\lean{HPC.HO16.zzLine}.}

Over a wedge edge, a curve $i\ge m$ entering at rank $i$ leaves at rank $i+2$: it climbs two sheets by running forward in sheet $i$, back in sheet $i+1$ and forward in
sheet $i+2$. The three pieces lie on one line and overlap. Definitions~\ref{def:straight}--\ref{def:stairwell} do not forbid sheets to meet, and the end sets
come out of the construction, so nothing needs to be separated. Compared with Hoehn and Oversteegen's picture, the wedge is put at the end of the edge and the zig-zags
before its tip.

\begin{example}\label{ex:ladder}
$G=[0,2]\subseteq\R$ with vertices $0,1,2$ and lists $(\frac12)$, $(\frac12)$, $(\frac14,\frac12,\frac34)$; no event on $(0,1)$ and a wedge $(1,\frac12,\frac38)$ at $2$ on $(1,2)$. Here $j=1$, $m=1$,
$l_1=\frac16$, $r_1=\frac13$ in the parameter of $[1,2]$. Then $\pi(S_1)=[0,\frac43]\cup[\frac32,2]$, $\pi(S_2)=[\frac76,\frac43]\cup[\frac32,2]$, $\pi(S_3)=[\frac76,2]$, $\pi(\alpha_2)=\set{\frac43,\frac32}$,
$\pi(\alpha_3)=\set{\frac76}$: a stairwell structure of height~$3$.
\end{example}

\formal{\lean{HPC.HO16.ladder_example}, \lean{HPC.HO16.ladder_example_height}, \lean{HPC.HO16.ladder_example_stairwell}.}

\section{A ladder is a stairwell}

\begin{theorem}\label{thm:ladder-stairwell}
Let a ladder on $G$ with at least one edge have height $k$. Then $k$ is odd, and the sheets $S_1,\dots,S_k$ form a stairwell structure of height $k$ for their union.
\end{theorem}

\begin{proof}
$k$ is the length of some list, hence odd.

\emph{The end sets.} On a wedge edge write $L_i(\mu)$ for the point of $L_i$ at parameter $\mu$. Let $\alpha_i$ consist of the points $L_{i-1}(r_{i-1})$ (the common right end of $C^1_{i-1}$ and
$C^2_{i-1}$, if $m\le i-1\le j$), $L_{i-2}(l_{i-2})$ (the common left end of $C^2_{i-2}$ and $C^3_{i-2}$, if $m\le i-2\le j$), and the tip (if $i=m+1$), over all wedge edges. Then
$\alpha_1=\emptyset$, and $\alpha_{k+1}=\emptyset$ because a wedge edge has $j(s_1)\le k-2$.

(S1) Over an edge, the pieces of sheet $i$ lie over the parameter intervals $[0,r_i]$, $[l_{i-1},r_{i-1}]$, $[l_{i-2},1]$ and, for $i\in\set{m,m+1}$, $[\lambda,1]$, which are pairwise disjoint;
over a vertex $v$ every edge at $v$ contributes the same point $(v,y_i(v))$ if $i\le j(v)$, and none otherwise. So $\pi$ is injective on each sheet, and $\pi(S_i)$ is a finite union of
nondegenerate segments, a PL graph (Proposition~\ref{prop:pl-library}(a)). A vertex with $j(v)=k$ gives every sheet a point.

(S2) A vertex $v$ is interior to $\pi(S_i)$ if $i\le j(v)$ (every edge at $v$ has a piece of sheet $i$ starting at $v$) and outside it otherwise. So $\rfr\pi(S_i)$ consists of the parameters
where pieces stop inside edges: $r_i$, $l_{i-1}$, $r_{i-1}$, $l_{i-2}$ and the tip. These are exactly the points of $\alpha_i\cup\alpha_{i+1}$, and they lie in $S_i$.

(S3) At each point of $\beta_i=\alpha_{i+1}$ the two sheets leave in the same direction: $C_i^1$ and $C_i^2$ both end at $r_i$ from the left; $C_{i-1}^2$ and $C_{i-1}^3$ both start at
$l_{i-1}$ to the right; the two wedge halves both start at the tip. Every other piece of the two sheets on that edge is at positive distance.

(S4) Let $C$ be a component of $G\setminus\pi(S_i)$. Inside one edge, $C$ is one of the gaps $(r_i,l_{i-1})$ (both ends in $\pi(\beta_i)$), $(r_{i-1},l_{i-2})$ (both ends in $\pi(\alpha_i)$), or a gap
next to the tip (both ends in $\pi(\beta_m)$, respectively in $\pi(\alpha_{m+1})$). Otherwise $C$ contains vertices, all with $j(v)<i$. Its frontier lies on wedge edges $s$ with $j(s_1)<i\le
j(s_1)+2$, at the first piece of sheet $i$ starting from $s_1$, and that end point is in $\beta_i$ if $i-j(s_1)=1$ and in $\alpha_i$ if $i-j(s_1)=2$. \emph{All $j(v)$ are odd}, so which case occurs
depends only on the parity of $i$: $\rfr C\subseteq\pi(\beta_i)$ for $i$ even and $\rfr C\subseteq\pi(\alpha_i)$ for $i$ odd. As $\pi(\alpha_i)\cap\pi(\beta_i)=\emptyset$, both halves of (S4) follow.

(S5) All end points lie over interior points of edges, which are ordinary, at pairwise different parameters.
\end{proof}

\formal{\lean{HPC.HO16.ho16_stairwell_of_ladder}, with $\alpha$ given by \lean{HPC.HO16.ladderAlphaSet}; the parity argument for (S4) is \lean{HPC.HO16.ladder_S4}.}

This is Hoehn and Oversteegen's parity argument of p.~198 (``we must have that $j_1$ and $j_n$ have the same parity and, hence, $j_1=j_n$''). Connectedness of $G$ is not
needed, the lists need not be sorted, and the heights need not be distinct.

\section{A ladder inside a neighbourhood of a separator}

\begin{lemma}[A sign function]\label{lem:sign-fn}
Let $G$ be a nonempty compact set, $M\subseteq G\times(0,1)$ separate $G\times\set0$ from $G\times\set1$ in $G^\star$, and $U\supseteq M$ open. There are a continuous $u\colon E\times\R\to\R$ and $\delta>0$
with $u\le0$ on $G\times\set0$, $u>0$ on $G\times\set1$, and: if $p\in G^\star$ is within $\delta$ of a point $c_1$ with $u(c_1)\le0$ and of a point $c_2$ with $u(c_2)>0$, then $p\in U\cap G\times(0,1)$.
\end{lemma}

\begin{proof}
Let $O_0\supseteq G\times\set0$ and $O_1\supseteq G\times\set1$ be the open sets of the separation. $K=G^\star\setminus(U\cap E\times(0,1))$ is compact and lies in $G^\star\setminus M\subseteq O_0\cup O_1$. Put
$K_0=K\setminus O_1\supseteq G\times\set0$ and $K_1=K\setminus O_0\supseteq G\times\set1$: disjoint compact sets with union $K$. Let $u=\dist(\cdot,K_0)-\dist(\cdot,K_1)$, a $2$-Lipschitz function, negative on
$K_0$ and positive on $K_1$, so $\omega=\min_K\abs u>0$. Note $G\times\set0\subseteq K_0$: it lies in $K$, in $O_0$, and hence not in $O_1$, as $O_0\cap O_1$ misses $G^\star\setminus M$; likewise
$G\times\set1\subseteq K_1$. Take $\delta=\omega/4$. If $p$ is as in the statement, then $\abs{u(p)-u(c_i)}<2\delta=\omega/2$. If $u(p)\ge0$, then $u(p)<u(c_1)+\omega/2\le\omega/2$; if $u(p)<0$,
then $u(p)>u(c_2)-\omega/2>-\omega/2$. So $\abs{u(p)}<\omega$, hence $p\notin K$, that is, $p\in U\cap G\times(0,1)$.
\end{proof}

\formal{\lean{HPC.HO16.exists_sign_fn}.}

\begin{theorem}[Replacing Claim~15.1]\label{thm:exists-ladder}
Let $G$ be a connected PL graph, $M\subseteq G\times(0,1)$ separate $G\times\set0$ from $G\times\set1$ in $G^\star$, and $U\supseteq M$ open. There is a ladder on $G$ with at least one edge all of
whose sheets lie in $U\cap G\times(0,1)$.
\end{theorem}

\begin{proof}
Take $u$ and $\delta$ from Lemma~\ref{lem:sign-fn}. Choose a decomposition $\mathcal S_0$ of $G$ with edges shorter than $\delta$, and $N$ with $1/N<\delta$; put $t_q=q/N$. At a vertex $v$ let
$\sigma(v,q)=-$ if $u(v,t_q)\le0$ and $+$ otherwise; so $\sigma(v,0)=-$ and $\sigma(v,N)=+$.

\emph{Heights at vertices.} Let $y(v)$ be the list of midpoints $\frac12(t_q+t_{q+1})$ of those $q$ with $\sigma(v,q)\neq\sigma(v,q+1)$, in increasing order. Since the signs go from $-$ to $+$, there is an
odd number of them.

\emph{One flip per edge.} For an edge $(v,w)$ of $\mathcal S_0$ let $q_1<\dots<q_r$ be the indices with $\sigma(v,q)\neq\sigma(w,q)$. Subdivide the edge into $r$ edges (one if $r=0$), and give the $p$-th new
vertex the sign sequence of $v$ with $q_1,\dots,q_p$ flipped (so the last one has that of $w$), and the list of midpoints of its sign changes. Each new edge flips one sign $q$, which changes the
sign changes at $q-1\mid q$ and at $q\mid q+1$:
\begin{itemize}
\item one appears and one disappears: one curve moves from $\frac12(t_{q-1}+t_q)$ to $\frac12(t_q+t_{q+1})$ or back, and no other curve moves; no event;
\item two appear: a wedge at the far end, with $\lambda=\frac12$, $h=t_q$ and $m$ the rank of the new lower curve;
\item two disappear: a wedge at the near end; orient the edge backwards.
\end{itemize}

\emph{The sheets lie in $U$.} At every new vertex the sign at each level $t_q$ is that of $v$ or of $w$. A curve at height $\frac12(t_q+t_{q+1})$ at a vertex means a sign change between $t_q$ and $t_{q+1}$
there. So every piece of a sheet lies in a set $s_0\times[t_q,t_{q+1}]$ ($s_0$ the original edge) whose corners carry both signs, or, near a flip at $q$, in $s_0\times[t_{q-1},t_{q+1}]$, where
$(v,t_q)$ and $(w,t_q)$ carry both signs. Every point of such a set is within $\max(\abs{s_0},1/N)<\delta$ of each of its corners (in the sup-norm on $E\times\R$), so by Lemma~\ref{lem:sign-fn} it lies in
$U\cap G\times(0,1)$.
\end{proof}

\formal{\lean{HPC.HO16.exists_ladder} (in \lean{HPC/HO16/LadderSep.lean}, \lean{HPC/HO16/LadderSubdiv.lean}, \lean{HPC/HO16/LadderSign.lean},
\lean{HPC/HO16/LadderExists.lean}). Connectedness of $G$ is used only for $G\neq\emptyset$.}

\section{The end of the proof}

\begin{proof}[Proof of Theorem~\ref{thm:ho15}]
Theorem~\ref{thm:exists-ladder} gives a ladder whose sheets lie in $U\cap G\times(0,1)$, and Theorem~\ref{thm:ladder-stairwell} makes their union a set with a stairwell structure of odd height.
\end{proof}

\formal{\lean{HPC.HO16.ho16_theorem15}.}

\begin{corollary}\label{cor:final}
Theorem~\ref{thm:ho20} holds; hence Theorem~\ref{thm:ho1} holds without hypothesis, and every nondegenerate homogeneous indecomposable plane continuum is homeomorphic to the
pseudo-arc (Theorem~\ref{thm:indec-pseudoarc}).
\end{corollary}

\begin{proof}
Theorem~\ref{thm:20-from-15} and Theorem~\ref{thm:ho15}.
\end{proof}

\formal{\lean{HPC.HO16.ho16_theorem20}, \lean{HPC.HO16.ho16_indec_homogeneous_pseudoArc}, \lean{HPC.HO16.ho16_indec_homogeneous_iff}.}

\section{On Hoehn and Oversteegen's proof}

For the reader comparing with the paper (pp.~194--198):
\begin{itemize}
\item Claim~15.1, the existence of a separator $M'$ with properties (1), (3), (5), (6), is left to the reader. The passage to an irreducible separator keeps (1) and (3), but this is not argued.
In the definition of a side wedge, ``$C_2$ are arcs which both have $x$ as an endpoint'' should read $(x,y)$. And in the paragraph on property~(4), the reason given for ``$(x,y)$ cannot be an
endpoint of $M'$'' (property~(3)) does not apply at points of $T$; irreducibility gives the step.
\item The odd number of points of $M'$ over each $a\in Z'$ (p.~196) is not argued. Here it comes from sign changes.
\item In the zig-zag construction nothing excludes $b\in J_j$; the arcs $J_i$ must lie in the open arc.
\item The case split in the (S4) argument omits a component not inside one arc whose frontier points lie on one arc (around a cycle of $G$); the same chain argument covers it.
\item That the final $S$ irreducibly separates is claimed but not argued, and is not needed: the end sets come out of the construction.
\end{itemize}

\section*{Exercises}

\begin{exercise}
Verify the end sets in Example~\ref{ex:ladder}, and check (S3) and (S4) for them.
\end{exercise}

\begin{exercise}
Construct a ladder on the circle (a triangle in $\R^2$) of height~$3$ with two wedges. Is the result consistent with Claim~14.1?
\end{exercise}

\begin{exercise}
Run the construction of Theorem~\ref{thm:exists-ladder} for $G=\I$, $M=\I\times\set{\frac12}$ and $u(x,t)=t-\frac12$ with $N=4$. What is the ladder?
\end{exercise}

\begin{exercise}
Why are all list lengths odd in Theorem~\ref{thm:exists-ladder}? What goes wrong in the (S4) argument of Theorem~\ref{thm:ladder-stairwell} if lists of even length were allowed?
\end{exercise}


\appendix
%======================================================================
\chapter{Correspondence with the formal proof}\label{app:formal}
%======================================================================

The table lists the definitions and results of this volume that carry a ``Formal counterpart''
line, together with the corresponding declarations of the Lean development. All declaration
names start with \lean{HPC.}. The formal statements are sometimes slightly more general, or are
split differently into lemmas; the table points to the place where the content of each result is
proved. Rows are in the order of the text; rows labelled ``Text of \dots'' belong to formal counterpart lines in the
running text of the named section.

The main theorems of the volume are, in the Lean development (prefix \lean{HPC.} omitted):
\begin{center}
\begin{tabular}{@{}p{0.42\textwidth}p{0.54\textwidth}@{}}
\toprule
Hagopian 1975, Lemma~1 (Theorem~\ref{thm:hagopian-lemma1}) & \lean{Hagopian.hagopian75_lemma1}\\
No proper subcontinuum separates (Theorem~\ref{thm:nosep}) & \lean{Hagopian.no_proper_subcontinuum_separates}\\
Moore's triod theorem (Theorem~\ref{thm:moore-triod}) & \lean{Hagopian.moore_triod_theorem'}\\
Type $\lambda$ continua exist (Theorem~\ref{thm:exists-lambda}) & \lean{Hagopian.hagopian75_exists_typeA'}\\
Hagopian 1984, Theorem~1 (Theorem~\ref{thm:h84}) & \lean{Hagopian.hagopian84_decomp}\\
No arc (Theorem~\ref{thm:noarc}) & \lean{Hagopian.jones75_noArc}\\
Hagopian 1976, Theorem~1 (Theorem~\ref{thm:HI}) & \lean{Hagopian.hagopian76_theorem1}\\
Cook's theorem (Theorem~\ref{thm:cook}) & \lean{Span.cook67_theorem4}\\
Strip theorem (Theorem~\ref{thm:ot32}) & \lean{Span.ot82_theorem32}\\
Theorem~C (Theorem~\ref{thm:theoremC}) & \lean{Span.theoremC}\\
Hoehn--Oversteegen, Theorem~4 (Theorem~\ref{thm:ho4}) & \lean{HO16.ho16_theorem4}\\
Krasinkiewicz--Minc (Lemma~\ref{lem:km}) & \lean{HO16.km77_lemma2}\\
Hoehn--Oversteegen, Theorem~8 (Theorem~\ref{thm:ho8}) & \lean{HO16.ho16_theorem8}\\
Hoehn--Oversteegen, Proposition~19 (Proposition~\ref{prop:ho19}) & \lean{HO16.ho16_prop19}\\
Theorem~20 from Theorem~15 (Theorem~\ref{thm:20-from-15}) & \lean{HO16.ho16_theorem20_of_thm15}\\
Hoehn--Oversteegen, Theorem~15 (Theorem~\ref{thm:ho15}) & \lean{HO16.ho16_theorem15}\\
Hoehn--Oversteegen, Theorem~20 (Theorem~\ref{thm:ho20}) & \lean{HO16.ho16_theorem20}\\
Hoehn--Oversteegen, Theorem~1 (Theorem~\ref{thm:ho1}) & \lean{HO16.ho16_theorem1_plane} with \lean{HO16.ho16_theorem20}\\
The indecomposable case (Theorem~\ref{thm:indec-pseudoarc}) & \lean{HO16.ho16_indec_homogeneous_pseudoArc}, \lean{HO16.ho16_indec_homogeneous_iff}\\
\bottomrule
\end{tabular}
\end{center}
All of them are proved in the Lean development without \lean{sorry} and without axioms beyond the
standard ones of Lean and Mathlib. The printed forms that are false (Section~\ref{sec:intro-diff}) are
refuted by the declarations whose names end in \lean{_false} or begin with \lean{not_}.

\begin{longtable}{@{}p{0.40\textwidth}p{0.56\textwidth}@{}}
\toprule
In these notes & In the Lean development\\
\midrule
\endhead
\bottomrule
\endfoot
\multicolumn{2}{@{}l}{\emph{Chapter~\ref{ch:homog}}}\\
Text of ``Standing conventions'' & \lean{HPC.Fold.Q78.IsContinuum}, \lean{HPC.Fold.Q78.IsIndecomposable}, \lean{HPC.Jones.IsDecomposableSpace}, \lean{HPC.Plane.IsHereditarilyIndecomposable}, \lean{HPC.Jones.IsHomogeneous}.\\
Text of ``Standing conventions'' & \lean{HPC.Plane.SeparatesPlane}.\\
Theorem~\ref{thm:composants} & \lean{HPC.Hagopian.composant_dense'}, \lean{HPC.Hagopian.union_ne_univ_of_indec}, \lean{HPC.Hagopian.composant_eq_or_disjoint'}, \lean{HPC.Hagopian.interior_eq_empty_of_proper}, \lean{HPC.Hagopian.composant_isMeagre'}, \lean{HPC.Hagopian.uncountable_composants'} (in \lean{HPC/Hagopian/H75Composants.lean}).\\
Lemma~\ref{lem:copies} & \lean{HPC.Hagopian.isContinuum_homeo_image}, \lean{HPC.Hagopian.homeo_image_ne_univ}, \lean{HPC.Hagopian.subset_composant}.\\
Proposition~\ref{prop:effros-nbhd} & \lean{HPC.Hagopian.hagopian75_lemma4}, \lean{HPC.Hagopian.hagopian76_lemma2}, \lean{HPC.Hagopian.hagopian84_lemma1}; all from \lean{HPC.Effros.effros_homogeneous}.\\
Proposition~\ref{prop:moore22} & \lean{HPC.Hagopian.separatesPlane_of_not_isPreconnected_inter}.\\
Lemma~\ref{lem:disjoint-domains} & \lean{HPC.Hagopian.disjoint_components_missing}.\\
Lemma~\ref{lem:sep-composants} & The argument appears twice in the formal development: in \lean{HPC.Hagopian.hereditarilyUnicoherent_of_homogeneous_indecomposable} (\lean{HPC/Hagopian/H75Lemma1.lean}) and in \lean{HPC.Hagopian.no_proper_subcontinuum_separates_aux} (\lean{HPC/Hagopian/H76NoSep.lean}); the countability step is \lean{HPC.Jones.countable_of_pairwiseDisjoint_open}, the Baire step \lean{HPC.Hagopian.exists_notMem_of_isMeagre}.\\
Theorem~\ref{thm:hu} & \lean{HPC.Hagopian.hereditarilyUnicoherent_of_homogeneous_indecomposable}, \lean{HPC.Hagopian.hagopian75_lemma1_hu}.\\
Theorem~\ref{thm:nosep} & \lean{HPC.Hagopian.exists_separating_subcontinuum_of_homeomorph}, \lean{HPC.Hagopian.separating_in_copy}, \lean{HPC.Hagopian.no_proper_subcontinuum_separates}.\\
Proposition~\ref{prop:moore73} & \lean{HPC.Hagopian.comp_subset_comp}, \lean{HPC.Hagopian.countable_of_separating}.\\
Lemma~\ref{lem:interlace} & \lean{HPC.Hagopian.interlace}, with \lean{HPC.Hagopian.mem_comp_of_mem_closure}, \lean{HPC.Hagopian.mem_of_mem_frontier_comp}, \lean{HPC.Hagopian.subset_of_disjoint_frontier}.\\
Text of ``Moore's triod theorem'' & \lean{HPC.Hagopian.exists_loop_param}, \lean{HPC.Hagopian.interlace_param}.\\
Lemma~\ref{lem:interlace-three} & \lean{HPC.Hagopian.interlace_three}.\\
Proposition~\ref{prop:moore79} & \lean{HPC.Hagopian.TriodData}, \lean{HPC.Hagopian.exists_leg}, \lean{HPC.Hagopian.exists_cross_continuum}, \lean{HPC.Hagopian.exists_rat_intervals}, \lean{HPC.Hagopian.countable_of_nonseparating}.\\
Text of ``Moore's triod theorem'' & \lean{HPC.Hagopian.moore_triod_theorem'}, \lean{HPC.Hagopian.mooreTriodStatement}; the statement as a proposition is \lean{HPC.Hagopian.MooreTriodStatement}.\\
Theorem~\ref{thm:hagopian-lemma1} & \lean{HPC.Hagopian.atriodic_of_homogeneous_indecomposable}, \lean{HPC.Hagopian.not_isTriod_univ_of_indecomposable}, \lean{HPC.Hagopian.isTriodSet_image}, \lean{HPC.Hagopian.hagopian75_lemma1}.\\
Corollary~\ref{cor:nested-or-arc} & \lean{HPC.Hagopian.isContinuum_inter_of_HU}; \lean{HPC.Plane.isTriod_of_core}.\\[4pt]
\multicolumn{2}{@{}l}{\emph{Chapter~\ref{ch:typelambda}}}\\
Text of ``Admissible maps and layers'' & \lean{HPC.IsIrreducibleBetween}; for the space itself, \lean{HPC.Hagopian.IrrBetween}.\\
Definition~\ref{def:admissible} & \lean{HPC.Hagopian.IsAdmissibleMap}, \lean{HPC.Hagopian.layer}, \lean{HPC.Hagopian.IsTypeA'}, \lean{HPC.Hagopian.IsMinimalAdmissible}, \lean{HPC.Hagopian.layerSet}.\\
Text of ``Admissible maps and layers'' & \lean{HPC.Hagopian.admissible_quotient_homeomorph}.\\
Lemma~\ref{lem:adm-basic} & \lean{HPC.Hagopian.IsAdmissibleMap.isPreconnected_preimage}, \lean{HPC.Hagopian.isContinuum_img}, \lean{HPC.Hagopian.isContinuum_layer}, \lean{HPC.Hagopian.end_layers}, \lean{HPC.Hagopian.layer_subset_of_meets'}, \lean{HPC.Hagopian.layer_subset_of_meets}.\\
Proposition~\ref{prop:end-sets} & \lean{HPC.Hagopian.dense_middle}, \lean{HPC.Hagopian.eq_of_meets_ends}, \lean{HPC.Hagopian.admissible_end_sets'}, \lean{HPC.Hagopian.hagopian76_end_sets}.\\
Remark~\ref{rem:end-sets-false} & \lean{HPC.Hagopian.not_admissible_end_sets}.\\
Lemma~\ref{lem:split-conn} & \lean{HPC.Hagopian.isPreconnected_of_union_closed}.\\
Proposition~\ref{prop:lemma2-nec} & \lean{HPC.Hagopian.not_isIndecomposable_of_typeA'}.\\
Theorem~\ref{thm:unique-decomp} & \lean{HPC.Hagopian.fibre_subset_fibre}, \lean{HPC.Hagopian.hagopian75_lemma3}, \lean{HPC.Hagopian.fibre_eq_iff}, \lean{HPC.Hagopian.hagopian75_lemma3_unique}.\\
Text of ``Kuratowski's family \texorpdfstring{$\mathcal D$}{D}'' & \lean{HPC.Hagopian.IsDElt}.\\
Lemma~\ref{lem:kur} & \lean{HPC.Hagopian.irr_compl_continuum}, \lean{HPC.Hagopian.irr_no_hole}, \lean{HPC.Hagopian.isDElt_closure_interior}, \lean{HPC.Hagopian.IsDElt.chain}, \lean{HPC.Hagopian.IsDElt.subset_interior}.\\
Proposition~\ref{prop:no-jumps} & \lean{HPC.Hagopian.IsDElt.closure_diff}, \lean{HPC.Hagopian.IsDElt.exists_between}.\\
Theorem~\ref{thm:scale} & \lean{HPC.Hagopian.exists_countable_dense_D}, \lean{HPC.Hagopian.IsDScale}, \lean{HPC.Hagopian.exists_scale}, \lean{HPC.Hagopian.scaleFun}, \lean{HPC.Hagopian.continuous_scaleFun}, \lean{HPC.Hagopian.isPreconnected_scaleFun_fibre}, \lean{HPC.Hagopian.interior_scaleFun_fibre}, \lean{HPC.Hagopian.scaleMap_spec}, \lean{HPC.Hagopian.exists_typeA'_map}.\\
Theorem~\ref{thm:lemma2} & \lean{HPC.Hagopian.isTypeA'_of_decomposable}, \lean{HPC.Hagopian.hagopian75_lemma2}.\\
Remark~\ref{rem:thomas21} & \lean{HPC.Hagopian.not_arc_of_degenerate_layers}, \lean{HPC.Hagopian.arc_of_degenerate_layers'}, \lean{HPC.Hagopian.exists_arc_of_degenerate_layers}.\\
Theorem~\ref{thm:exists-lambda} & \lean{HPC.Hagopian.no_indecomposable_interior_aux}, \lean{HPC.Hagopian.hagopian75_no_indecomposable_interior}, \lean{HPC.Hagopian.hagopian75_exists_typeA'}.\\
Theorem~\ref{thm:layer-invariant} & \lean{HPC.Hagopian.homeo_image_subset_of_meets}, \lean{HPC.Hagopian.homeo_image_layer_subset}, \lean{HPC.Hagopian.hagopian75_layer_invariant_aux}, \lean{HPC.Hagopian.hagopian75_layer_invariant}.\\
Text of ``Interior layers under small homeomorphisms'' & \lean{HPC.Hagopian.layer_homogeneous_aux}, \lean{HPC.Hagopian.hagopian75_layer_homogeneous}.\\
Theorem~\ref{thm:h76-lemma3} & \lean{HPC.Hagopian.hagopian76_lemma3}, proved in \lean{HPC/Hagopian/H76Lemma3.lean}.\\[4pt]
\multicolumn{2}{@{}l}{\emph{Chapter~\ref{ch:decomp}}}\\
Theorem~\ref{thm:h84} & \lean{HPC.Hagopian.Hagopian84Decomp} (the statement), \lean{HPC.Hagopian.hagopian84_decomp} (the proof).\\
Lemma~\ref{lem:rigid} & \lean{HPC.Hagopian.subset_of_meets_avoids}, \lean{HPC.Hagopian.no_arm_at_interior_layer}.\\
Theorem~\ref{thm:h76-lemma5} & \lean{HPC.Hagopian.IsExtensionAway}, \lean{HPC.Hagopian.hagopian76_lemma4_chain}, \lean{HPC.Hagopian.hagopian76_lemma4}, \lean{HPC.Hagopian.hagopian76_lemma5}.\\
Corollary~\ref{cor:restrict} & \lean{HPC.Hagopian.minimalAdmissible_restrict}.\\
Theorem~\ref{thm:h84-lemma6} & \lean{HPC.Hagopian.IsEssential84}, \lean{HPC.Hagopian.hagopian84_lemma6}.\\
Theorem~\ref{thm:layers-agree} & \lean{HPC.Hagopian.layers_agree}, \lean{HPC.Hagopian.hagopian84_lemma7_aux}.\\
Text of ``Rigidity of layers'' & \lean{HPC.Hagopian.hagopian76_lemma6_X}, \lean{HPC.Hagopian.hagopian84_lemma11}.\\
Definition~\ref{def:chain5} & \lean{HPC.Hagopian.IsChain5}, \lean{HPC.Hagopian.midUnion}, \lean{HPC.Hagopian.IsNormal5}, \lean{HPC.Hagopian.IsFree5}, \lean{HPC.Hagopian.RunsStraight}, \lean{HPC.Hagopian.IsRegular5}, \lean{HPC.Hagopian.IsOrderedRefinement5}.\\
Lemma~\ref{lem:fold} & \lean{HPC.Hagopian.exists_fold}.\\
Theorem~\ref{thm:h84-lemma13} & \lean{HPC.Hagopian.WsetOf}, \lean{HPC.Hagopian.chainLinks84}, \lean{HPC.Hagopian.chainLinks84_chain}, \lean{HPC.Hagopian.chainLinks84_runsStraight}, \lean{HPC.Hagopian.chainLinks84_regular}, \lean{HPC.Hagopian.hagopian84_lemma12}, \lean{HPC.Hagopian.hagopian84_lemma13}.\\
Text of ``The layers of a chain'' & \lean{HPC.Hagopian.Omega5}, \lean{HPC.Hagopian.Delta5}.\\
Proposition~\ref{prop:delta} & \lean{HPC.Hagopian.delta5_partition}, \lean{HPC.Hagopian.delta5_dist_pos}, \lean{HPC.Hagopian.hagopian84_lemma14}, \lean{HPC.Hagopian.hagopian84_lemma15}, \lean{HPC.Hagopian.exists_interval_in_delta}.\\
Theorem~\ref{thm:lemma16p} & \lean{HPC.Hagopian.strand_layer_mem_delta}, \lean{HPC.Hagopian.not_strand_subset_extension_layer}, \lean{HPC.Hagopian.strand_layer_of_extension}; Hagopian's Lemma~16 as printed is \lean{HPC.Hagopian.hagopian84_lemma16}.\\
Proposition~\ref{prop:bridges} & \lean{HPC.Hagopian.omega_not_delta_side}, \lean{HPC.Hagopian.exists_bridge84}, \lean{HPC.Hagopian.mem_delta_of_bridge_chain}, \lean{HPC.Hagopian.exists_chainZ84}.\\
Proposition~\ref{prop:finite-cover} & \lean{HPC.Hagopian.exists_finite_cover84}, \lean{HPC.Hagopian.hagopian84_cover}, \lean{HPC.Hagopian.CoverData84}, \lean{HPC.Hagopian.exists_coverData84}.\\
Text of ``The decomposition'' & \lean{HPC.Hagopian.sUnion_cover84_eq_univ}.\\
Theorem~\ref{thm:frakD} & \lean{HPC.Hagopian.omega_layer_of_mem_mid}, \lean{HPC.Hagopian.omega_mem_delta_of_mem_mid}, \lean{HPC.Hagopian.CoverData84.members}, \lean{HPC.Hagopian.CoverData84.members_disjoint}, \lean{HPC.Hagopian.CoverData84.omega_subset_members}, \lean{HPC.Hagopian.CoverData84.invariant}, \lean{HPC.Hagopian.CoverData84.isJonesDecomp}.\\
Text of ``The quotient'' & \lean{HPC.Hagopian.IsJonesDecomp84}.\\
Text of ``The quotient'' & \lean{HPC.Hagopian.exists_comp_small_homeo}, \lean{HPC.Hagopian.hagopian84Decomp_of_package}, \lean{HPC.Hagopian.hagopian84_decomp}.\\[4pt]
\multicolumn{2}{@{}l}{\emph{Chapter~\ref{ch:noarc}}}\\
Definition~\ref{def:jones-setting} & \lean{HPC.Hagopian.JonesSetting}, \lean{HPC.Hagopian.jonesSetting_of_plane}.\\
Proposition~\ref{prop:jones-c} & \lean{HPC.Hagopian.JonesSetting.inter_isPreconnected}, \lean{HPC.Hagopian.JonesSetting.no_three_arms}, \lean{HPC.Hagopian.JonesSetting.union_arc}, \lean{HPC.Hagopian.JonesSetting.nested}.\\
Text of ``Arcs and lines'' & \lean{HPC.Hagopian.JonesSetting.arc_unique}, \lean{HPC.Hagopian.JonesSetting.exists_arc_through}.\\
Text of ``Arcs and lines'' & \lean{HPC.Hagopian.arcComp}, \lean{HPC.Hagopian.IsLine}.\\
Theorem~\ref{thm:lines} & \lean{HPC.Hagopian.JonesSetting.rayArcs_total}, \lean{HPC.Hagopian.JonesSetting.extend}, \lean{HPC.Hagopian.JonesSetting.exists_ray_seq}, \lean{HPC.Hagopian.JonesSetting.exists_line}, \lean{HPC.Hagopian.JonesSetting.line_arc}, \lean{HPC.Hagopian.JonesSetting.arc_in_line}, \lean{HPC.Hagopian.JonesSetting.homeo_on_line}.\\
Lemma~\ref{lem:R0} & \lean{HPC.Hagopian.JonesSetting.no_reversal_arc}.\\
Theorem~\ref{thm:RO} & \lean{HPC.Hagopian.JonesSetting.local_no_reversal}, \lean{HPC.Hagopian.JonesSetting.exists_no_reversal}, \lean{HPC.Hagopian.JonesSetting.orientation}.\\
Text of ``Recurrence'' & \lean{HPC.Hagopian.omegaLim}, \lean{HPC.Hagopian.mem_omegaLim_iff}.\\
Theorem~\ref{thm:recurrent} & \lean{HPC.Hagopian.JonesSetting.omegaLim_saturated}, \lean{HPC.Hagopian.JonesSetting.zero_mem_omegaLim}, \lean{HPC.Hagopian.JonesSetting.line_recurrent}, \lean{HPC.Hagopian.JonesSetting.jones75_d}.\\
Text of ``Passes through a box'' & \lean{HPC.Hagopian.IsCollar}, \lean{HPC.Hagopian.exists_collar}, \lean{HPC.Hagopian.collarCoord}, \lean{HPC.Hagopian.flipC}.\\
Definition~\ref{def:pass} & \lean{HPC.Hagopian.IsPass}, \lean{HPC.Hagopian.PassMon}, \lean{HPC.Hagopian.PassAbove}, \lean{HPC.Hagopian.PassBelow}.\\
Lemma~\ref{lem:passes} & \lean{HPC.Hagopian.IsPass.disjoint}, \lean{HPC.Hagopian.IsPass.finite}, \lean{HPC.Hagopian.exists_first_pass}, \lean{HPC.Hagopian.isPass_base}, \lean{HPC.Hagopian.IsPass.above_or_below}, \lean{HPC.Hagopian.IsPass.mono}.\\
Lemma~\ref{lem:shadow} & \lean{HPC.Hagopian.shadow_pass}.\\
Corollary~\ref{cor:scales} & \lean{HPC.Hagopian.flow_box}, \lean{HPC.Hagopian.FBProp}, \lean{HPC.Hagopian.PassThin}, \lean{HPC.Hagopian.scale_two}, \lean{HPC.Hagopian.exists_late_thin_pass}.\\
Lemma~\ref{lem:region} & \lean{HPC.Hagopian.rect_domain_aux}, \lean{HPC.Hagopian.rect_domain_low}, \lean{HPC.Hagopian.seg_subset_rect_domain}, \lean{HPC.Hagopian.rect_domain_crosscut}, \lean{HPC.Hagopian.thinPass_arc}.\\
Lemma~\ref{lem:R1R2} & \lean{HPC.Hagopian.region_pass}, \lean{HPC.Hagopian.region_cross}.\\
Theorem~\ref{thm:trap} & \lean{HPC.Hagopian.region_lemmaA}, \lean{HPC.Hagopian.two_sides}, \lean{HPC.Hagopian.trap_jordan}, \lean{HPC.Hagopian.trap_pass}; the Jordan curve theorem is \lean{HPC.Jones.jordan_curve_theorem}.\\
Theorem~\ref{thm:noarc} & \lean{HPC.Hagopian.jones75_final_aux}, \lean{HPC.Hagopian.jones75_contra}, \lean{HPC.Hagopian.NoArcIndecHomPlane}, \lean{HPC.Hagopian.jones75_noArc}.\\
Lemma~\ref{lem:monotone-image} & \lean{HPC.Hagopian.indecomposable_of_monotone_image}.\\
Theorem~\ref{thm:HI} & \lean{HPC.Hagopian.exists_plane_quotient_of_monotone}, \lean{HPC.Hagopian.hagopian76_theorem1_of_jones}, \lean{HPC.Hagopian.hagopian76_theorem1_jones}, \lean{HPC.Hagopian.Hagopian76Thm1}, \lean{HPC.Hagopian.hagopian76_theorem1}.\\
Corollary~\ref{cor:bing60} & \lean{HPC.Hagopian.bing60_plane_uncond}, \lean{HPC.Hagopian.hagopian76_theorem2_uncond}.\\[4pt]
\multicolumn{2}{@{}l}{\emph{Chapter~\ref{ch:span}}}\\
Theorem~\ref{thm:theoremC} & \lean{HPC.Span.TheoremC}, \lean{HPC.Span.theoremC}.\\
Definition~\ref{def:span} & \lean{HPC.Span.OffDiag}, \lean{HPC.Span.SurjSpanPos}, \lean{HPC.Span.SurjSemiSpanPos}, \lean{HPC.Span.SpanZero}, \lean{HPC.Span.SemiSpanZero}, \lean{HPC.Span.HOSpanZero}.\\
Proposition~\ref{prop:span-basic} & \lean{HPC.Span.otSemiSpanPos_iff}, \lean{HPC.Span.otSpanPos_iff}, \lean{HPC.Span.SurjSpanPos.image}, \lean{HPC.Span.SurjSemiSpanPos.image}, \lean{HPC.Span.semiSpanZero_iff_hoSpanZero}, \lean{HPC.Span.hoSpanZero_iff_maps}.\\
Theorem~\ref{thm:cook} & \lean{HPC.Span.nested_of_hi}, \lean{HPC.Span.cook67_theorem4}.\\
Lemma~\ref{lem:brouwer} & \lean{HPC.Span.isContinuum_sInter_chain}, \lean{HPC.Span.brouwer_reduction}.\\
Theorem~\ref{thm:davis} & \lean{HPC.Span.exists_surjSpanPos_of_surjSemiSpanPos}, \lean{HPC.Span.semiSpanZero_of_spanZero}.\\
Text of ``Theorem 3.6 from Corollary 3.5'' & \lean{HPC.Span.OTCor35}.\\
Lemma~\ref{lem:dense-proper} & \lean{HPC.Span.union_proper_of_indec}, \lean{HPC.Span.exists_dense_proper_subcontinuum}.\\
Theorem~\ref{thm:ot36} & \lean{HPC.Span.exists_reduced_witness}, \lean{HPC.Span.not_surjSemiSpanPos_univ}.\\
Corollary~\ref{cor:C-from-35} & \lean{HPC.Span.ot82_theorem36_of_cor35}, \lean{HPC.Span.theoremC_of_cor35}.\\
Definition~\ref{def:otthm32} & \lean{HPC.Span.OTThm32}.\\
Lemma~\ref{lem:reembed} & \lean{HPC.Span.exists_homeomorph_unique_leftmost}, \lean{HPC.Span.exists_homeomorph_ball_step}, \lean{HPC.Span.exists_homeomorph_push}, \lean{HPC.Span.isConnected_rightHalf_diff}.\\
Lemma~\ref{lem:crossing} & \lean{HPC.Span.exists_crossing_subcontinuum}, \lean{HPC.Span.exists_uncountable_crossing}.\\
Theorem~\ref{thm:ot35} & \lean{HPC.Span.cor35_core}, \lean{HPC.Span.ot82_cor35_of_thm32}, \lean{HPC.Span.theoremC_of_thm32}.\\[4pt]
\multicolumn{2}{@{}l}{\emph{Chapter~\ref{ch:strips}}}\\
Chapter introduction & \lean{HPC.Plane.parity}, \lean{HPC.Plane.polygonOf}.\\
Lemma~\ref{lem:parity-tools} & \lean{HPC.Span.parity_eq_zero_of_not_mem_convex}, \lean{HPC.Span.parity_eq_of_perturb}, \lean{HPC.Span.parity_flip_vertical_edge}.\\
Lemma~\ref{lem:common-zero} & \lean{HPC.Span.common_zero_rect}, using \lean{HPC.Moore.exists_log_of_convex}.\\
Text of ``The strip lemma'' & \lean{HPC.Span.plPath}, \lean{HPC.Span.plPath_piece}, \lean{HPC.Span.stripPoly}.\\
Lemma~\ref{lem:strip-basic} & \lean{HPC.Span.parity_strip_eventually}.\\
Text of ``The strip lemma'' & \lean{HPC.Span.StripInner}, \lean{HPC.Span.stripSign}.\\
Theorem~\ref{thm:strip-lemma} & \lean{HPC.Span.strip_lemma}.\\
Proposition~\ref{prop:ordered} & \lean{HPC.Span.crossing_traces_separated}.\\
Theorem~\ref{thm:ot32} & \lean{HPC.Span.exists_not_countable_of_cover}, \lean{HPC.Span.exists_between_three}, \lean{HPC.Span.ot82_theorem32}.\\
Text of ``Proof of the strip theorem'' & \lean{HPC.Span.theoremC}.\\[4pt]
\multicolumn{2}{@{}l}{\emph{Chapter~\ref{ch:snake}}}\\
Definition~\ref{def:plgraph} & \lean{HPC.HO16.IsPLGraph}, \lean{HPC.HO16.SeparatesIn}, \lean{HPC.HO16.IsFanAt}.\\
Text of ``Statements'' & \lean{HPC.HO16.IsPLGraph.isCompact}, \lean{HPC.HO16.IsPLGraph.exists_fan}, \lean{HPC.HO16.separatesIn_of_closed}.\\
Theorem~\ref{thm:ho20} & The statement is the proposition \lean{HPC.HO16.HOThm20}; it is proved as \lean{HPC.HO16.ho16_theorem20} (Chapter~\ref{ch:ladders}).\\
Theorem~\ref{thm:ho1} & \lean{HPC.HO16.ho16_theorem1_plane}, which takes Theorem~\ref{thm:ho20} as a hypothesis \lean{h20 : HOThm20}; span zero is \lean{HPC.Span.HOSpanZero} (Definition~\ref{def:span}).\\
Theorem~\ref{thm:indec-pseudoarc} & \lean{HPC.HO16.isPseudoArc_of_indec_homogeneous} and \lean{HPC.HO16.homeomorph_pseudoArc_of_indec_homogeneous} (with \lean{h20} as hypothesis); unconditionally \lean{HPC.HO16.ho16_indec_homogeneous_pseudoArc}; the converse \lean{HPC.HO16.ho16_pseudoArc_indec_homogeneous}; both directions \lean{HPC.HO16.ho16_indec_homogeneous_iff}. Non-vacuity of the hypotheses: \lean{HPC.HO16.exists_indec_homogeneous_plane}.\\
Remark~\ref{rem:ho1-point} & \lean{HPC.HO16.ho16_theorem1_printed_false}, \lean{HPC.HO16.ho16_theorem1prime_printed_false}.\\
Theorem~\ref{thm:ho4} & \lean{HPC.HO16.ho16_theorem4}; Step~1 is \lean{HPC.HO16.thm4_continuum_of_not_separates}, Steps~2--4 \lean{HPC.HO16.thm4_limit_contra}, using \lean{HPC.HO16.subset_of_tendsto_of_eventually} for limits of compacta.\\
Lemma~\ref{lem:hi-interior} & \lean{HPC.HO16.interior_eq_empty_of_hi}.\\
Lemma~\ref{lem:grid} & \lean{HPC.HO16.exists_plGraph_approx}, with \lean{HPC.HO16.radProj}, \lean{HPC.HO16.rectGauge}, \lean{HPC.HO16.radProj_of_onBdry}, \lean{HPC.HO16.radProj_mem}, \lean{HPC.HO16.gridEdge_inter_subset} (in \lean{HPC/HO16/Grid.lean}).\\
Lemma~\ref{lem:path-approx} & \lean{HPC.HO16.exists_path_approx}, with \lean{HPC.HO16.exists_chain_through}.\\
Text of ``The proof of Theorem 1 from Theorem 20'' & \lean{HPC.HO16.ho16_theorem1_plane}; the $\varepsilon$-map criterion is \lean{HPC.Plane.isSnakeLike_univ_of_epsMap}.\\[4pt]
\multicolumn{2}{@{}l}{\emph{Chapter~\ref{ch:folds}}}\\
Lemma~\ref{lem:km} & \lean{HPC.HO16.km77_lemma2}, with \lean{HPC.HO16.isClosed_biUnion_connectedComponentIn} for Step~2.\\
Theorem~\ref{thm:ho8-core} & \lean{HPC.HO16.ho16_theorem8_core}, with \lean{HPC.HO16.exists_center_fun}, \lean{HPC.HO16.continuousOn_center_comp}, \lean{HPC.HO16.continuous_correction} and \lean{HPC.HO16.exists_open_nbhd_side}.\\
Proposition~\ref{prop:pl-library} & \lean{HPC.HO16.isPLGraph_biUnion}, \lean{HPC.HO16.IsPLGraph.union}, \lean{HPC.HO16.IsPLGraph.exists_refinement}, \lean{HPC.HO16.IsPLGraph.closure_diff}, \lean{HPC.HO16.IsPLGraph.relFrontier_finite}, \lean{HPC.HO16.IsPLGraph.connectedComponentIn}, \lean{HPC.HO16.IsPLGraph.finite_components}, \lean{HPC.HO16.IsPLGraph.image_fst}, \lean{HPC.HO16.Subdiv.exists_common_subdivision}; the relative frontier is \lean{HPC.HO16.relFrontier}.\\
Definition~\ref{def:fold} & \lean{HPC.HO16.IsFoldTriple}, \lean{HPC.HO16.IsSimpleFoldOf}.\\
Theorem~\ref{thm:fold-exists} & \lean{HPC.HO16.exists_simpleFold}, with the lifted graphs \lean{HPC.HO16.liftG} and the midpoint subdivision \lean{HPC.HO16.exists_midpoint_subdivision}.\\
Lemma~\ref{lem:ho6} & \lean{HPC.HO16.ho16_lemma6}; Hoehn and Oversteegen leave the proof to the reader.\\
Theorem~\ref{thm:ho8} & \lean{HPC.HO16.ho16_theorem8}, with \lean{HPC.HO16.continuousOn_invFunOn_fst} for $\psi_i$.\\
Proposition~\ref{prop:ho7} & \lean{HPC.HO16.ho16_prop7}, with \lean{HPC.HO16.exists_comp_meets_both}, \lean{HPC.HO16.image_comp_eq}, \lean{HPC.HO16.lift_subset_of_meets_both}.\\[4pt]
\multicolumn{2}{@{}l}{\emph{Chapter~\ref{ch:stairwells}}}\\
Definition~\ref{def:sides} & \lean{HPC.HO16.HasConsCompl}, \lean{HPC.HO16.sideOf}.\\
Text of ``Sides'' & \lean{HPC.HO16.IsPLGraph.finite_components_diff}, \lean{HPC.HO16.IsPLGraph.comp_diff_mem_nhdsWithin}, \lean{HPC.HO16.IsPLGraph.closure_comp_diff}, \lean{HPC.HO16.IsPLGraph.relFrontier_comp_diff}.\\
Proposition~\ref{prop:sides-char} & \lean{HPC.HO16.ho16_obs9_emptyB}, \lean{HPC.HO16.ho16_obs9_nhds}, \lean{HPC.HO16.ho16_obs9_char}, \lean{HPC.HO16.ho16_obs9_unique}; the common core is \lean{HPC.HO16.IsPLGraph.sideOf_char_aux}.\\
Proposition~\ref{prop:ho10} & \lean{HPC.HO16.ho16_prop10}.\\
Proposition~\ref{prop:ho11} & \lean{HPC.HO16.ho16_prop11}, \lean{HPC.HO16.ho16_prop11_frontier}, \lean{HPC.HO16.ho16_prop11_remark}.\\
Definition~\ref{def:straight} & \lean{HPC.HO16.IsOrdinaryPt}, \lean{HPC.HO16.IsGeneric}, \lean{HPC.HO16.IsStraight}, \lean{HPC.HO16.endSet}.\\
Text of ``Straight sets and stairwell structures'' & \lean{HPC.HO16.isOrdinaryPt_iff_order}, \lean{HPC.HO16.IsPLGraph.finite_not_ordinary}, \lean{HPC.HO16.ho16_obs12}.\\
Definition~\ref{def:stairwell} & \lean{HPC.HO16.IsStairwellStr} (with the sets $\alpha_i$ as one sequence, so that $\beta_i=\alpha_{i+1}$ holds by construction) and its existential form \lean{HPC.HO16.HasStairwell}.\\
Proposition~\ref{prop:stair-remarks} & \lean{HPC.HO16.ho16_obs13_comp}, \lean{HPC.HO16.ho16_height_one}, \lean{HPC.HO16.ho16_claim14_1}, \lean{HPC.HO16.ho16_odd_image}, \lean{HPC.HO16.ho16_p193_remark_false}.\\
Definition~\ref{def:broken} & \lean{HPC.HO16.IsBrokenStairwellStr}; the restriction to a component is \lean{HPC.HO16.ho16_obs16_comp}.\\
Proposition~\ref{prop:ho17} & \lean{HPC.HO16.ho16_prop17}, \lean{HPC.HO16.ho16_prop17_exists}. The hypothesis $k\ge3$ is not printed; for $k\le2$ the conclusion has height $\le0$.\\
Proposition~\ref{prop:rem17} & \lean{HPC.HO16.ho16_rem17_gamma}, \lean{HPC.HO16.ho16_rem17_P1}.\\
Lemma~\ref{lem:ho18} & \lean{HPC.HO16.ho16_lemma18}, with the preimage \lean{HPC.HO16.foldPre} and $\varphi^\star$ as \lean{HPC.HO16.phiStar}; PL-graph facts for preimages: \lean{HPC.HO16.IsPLGraph.inter_preimage_fst}, \lean{HPC.HO16.isPLGraph_union_finite}.\\
Proposition~\ref{prop:ho19} & \lean{HPC.HO16.ho16_prop19}, whose statement is the proposition \lean{HPC.HO16.HOProp19}; for the full (possibly disconnected) fold, \lean{HPC.HO16.ho16_prop19_full}.\\
Text of ``Proposition 19: unwinding one turn'' & The checks are spread over \lean{HPC/HO16/Prop19A.lean} to \lean{HPC/HO16/Prop19D.lean}; the fold is \lean{HPC.HO16.ho16_prop19_fold}; the sheets are \lean{HPC.HO16.p19T}, with a companion definition for the end sets in \lean{HPC/HO16/P19Defs.lean}; ordinary points of folds: \lean{HPC.HO16.IsSimpleFoldOf.isOrdinaryPt_iff}.\\
Remark~\ref{rem:p208} & \lean{HPC.HO16.ho16_p208_false}.\\
Theorem~\ref{thm:ho15} & The statement is the proposition \lean{HPC.HO16.HOThm15}; it is proved as \lean{HPC.HO16.ho16_theorem15} (Chapter~\ref{ch:ladders}).\\
Theorem~\ref{thm:lift-stairwell} & \lean{HPC.HO16.ho16_lift_stairwell}, with the induction in \lean{HPC.HO16.ho16_lift_aux}.\\
Theorem~\ref{thm:20-from-15} & \lean{HPC.HO16.ho16_theorem20_of_thm15}; with Theorem~15 as the only hypothesis, the classification step is \lean{HPC.HO16.homeomorph_pseudoArc_of_indec_homogeneous_of_thm15}.\\[4pt]
\multicolumn{2}{@{}l}{\emph{Chapter~\ref{ch:ladders}}}\\
Definition~\ref{def:ladder} & \lean{HPC.HO16.IsLadder}, \lean{HPC.HO16.ladderY}, \lean{HPC.HO16.ladderHeight} (in \lean{HPC/HO16/LadderDefs.lean}).\\
Text of ``Ladders'' & \lean{HPC.HO16.ladderPiece}, \lean{HPC.HO16.ladderSheet}; the zig-zag parameters are \lean{HPC.HO16.zzL}, \lean{HPC.HO16.zzR} and the lines \lean{HPC.HO16.zzLine}.\\
Example~\ref{ex:ladder} & \lean{HPC.HO16.ladder_example}, \lean{HPC.HO16.ladder_example_height}, \lean{HPC.HO16.ladder_example_stairwell}.\\
Theorem~\ref{thm:ladder-stairwell} & \lean{HPC.HO16.ho16_stairwell_of_ladder}, with $\alpha$ given by \lean{HPC.HO16.ladderAlphaSet}; the parity argument for (S4) is \lean{HPC.HO16.ladder_S4}.\\
Lemma~\ref{lem:sign-fn} & \lean{HPC.HO16.exists_sign_fn}.\\
Theorem~\ref{thm:exists-ladder} & \lean{HPC.HO16.exists_ladder} (in \lean{HPC/HO16/LadderSep.lean}, \lean{HPC/HO16/LadderSubdiv.lean}, \lean{HPC/HO16/LadderSign.lean}, \lean{HPC/HO16/LadderExists.lean}). Connectedness of $G$ is used only for $G\neq\emptyset$.\\
Text of ``The end of the proof'' & \lean{HPC.HO16.ho16_theorem15}.\\
Corollary~\ref{cor:final} & \lean{HPC.HO16.ho16_theorem20}, \lean{HPC.HO16.ho16_indec_homogeneous_pseudoArc}, \lean{HPC.HO16.ho16_indec_homogeneous_iff}.\\[4pt]
\end{longtable}


\end{document}
